\documentclass[11pt]{book}

\usepackage{fontspec}
\setmainfont{TeX Gyre Pagella}

\usepackage{amsmath}
\usepackage{amssymb}
\usepackage{amsthm}
\usepackage{mathtools}
\usepackage{mathrsfs}
\usepackage{bm}

\usepackage{microtype}

\usepackage{geometry}
\geometry{margin=1.1in}

\usepackage{xcolor}

\usepackage{graphicx}
\usepackage{tikz}
\usetikzlibrary{
arrows.meta,
calc,
positioning,
matrix,
fit,
decorations.pathmorphing
}

\usepackage{booktabs}
\usepackage{longtable}
\usepackage{array}
\usepackage{multirow}

\usepackage{enumitem}

\usepackage{hyperref}
\hypersetup{
colorlinks=true,
linkcolor=blue!60!black,
citecolor=blue!60!black,
urlcolor=blue!60!black
}

\usepackage[nameinlink]{cleveref}

\numberwithin{equation}{chapter}

\theoremstyle{definition}
\newtheorem{definition}{Definition}[chapter]
\newtheorem{example}[definition]{Example}
\newtheorem{remark}[definition]{Remark}

\theoremstyle{plain}
\newtheorem{theorem}[definition]{Theorem}
\newtheorem{proposition}[definition]{Proposition}
\newtheorem{lemma}[definition]{Lemma}
\newtheorem{corollary}[definition]{Corollary}

\DeclareMathOperator{\Adm}{Adm}
\DeclareMathOperator{\Reach}{Reach}
\DeclareMathOperator{\id}{id}
\DeclareMathOperator{\Vol}{Vol}
\DeclareMathOperator{\Tr}{Tr}
\DeclareMathOperator{\diag}{diag}

\newcommand{\Sspace}{\mathcal{S}}
\newcommand{\Tspace}{\mathcal{T}}
\newcommand{\Mspace}{\mathcal{M}}
\newcommand{\Xspace}{\mathcal{X}}
\newcommand{\Pspace}{\mathcal{P}}

\newcommand{\R}{\mathbb{R}}
\newcommand{\N}{\mathbb{N}}

\title{\Huge Admissible Continuation\\[1em]
\Large A Mathematical Theory of Persistent Systems}

\author{Flyxion}

\date{\today}

\begin{document}

\frontmatter

\maketitle

\begin{abstract}

Persistent systems exhibit an apparently paradoxical property. They undergo continual change while nevertheless remaining identifiable through time. Organisms replace their constituent matter, scientific theories are repeatedly revised, software systems evolve through successive deployments, institutions adapt to new environments, and intelligent agents continually update their internal representations. Although the mechanisms governing these systems differ dramatically, they share a common structural problem: they must execute state transitions without destroying the invariant conditions that make further transitions possible.

This monograph develops a mathematical theory of persistent systems founded upon the notion of admissible continuation. Rather than treating states or actions as primitive objects, the theory adopts transitions as its fundamental unit of analysis and characterizes persistence through a distinguished relation on the transition space. A transition is said to be admissible precisely when it preserves an explicitly defined family of structural invariants governing the continued existence of the system. Under this formulation, repair, refusal, intervention, rollback, verification, and maintenance emerge as natural operators acting upon admissible transition spaces rather than as independent engineering mechanisms.

The mathematical framework developed throughout this volume proceeds from the concepts of distinction, information, boundary preservation, and continuation. Admissibility is formulated as a relational property of transitions relative to invariant families, leading to an operator calculus describing proposal, evaluation, commitment, repair, rollback, maintenance, and state-preserving refusal. General aggregation operators are introduced for combining multiple invariant valuations, allowing both strict bottleneck and probabilistic formulations of admissibility. Geometric interpretations are subsequently developed in terms of admissibility manifolds, boundary metrics, repair flows, hysteresis, and trajectory deformation. These constructions naturally extend from isolated transitions to histories, recursive processes, and indefinitely continuing systems.

The resulting theory is intentionally domain-independent. Physical conservation laws, biological homeostasis, logical consistency, computational correctness, institutional stability, scientific methodology, and artificial intelligence are treated as distinct realizations of a common mathematical structure in which persistence depends upon the preservation of invariant relations under successive transformations. Modern agentic software architectures appear not as the primary subject of the theory but as contemporary examples of a broader class of admissibility-preserving systems whose emergence reflects the increasing necessity of separating unconstrained proposal generation from warranted state commitment.

The central claim of this work is that persistence cannot be understood solely in terms of reachable states, optimal actions, or generated outputs. It is fundamentally a property of admissible trajectories through structured state spaces. Intelligence, adaptation, and long-term viability consequently emerge not from the unrestricted capacity to generate possibilities but from the disciplined preservation of those invariant conditions that permit meaningful continuation across time.

\end{abstract}

\newpage
\tableofcontents

\mainmatter

\chapter*{Preface}
\addcontentsline{toc}{chapter}{Preface}

Every enduring scientific theory begins by identifying an object that previous theories treated only implicitly. Classical mechanics elevated trajectories above isolated observations. Thermodynamics shifted attention from individual particles to macroscopic state variables. Information theory demonstrated that communication could be studied independently of the semantic content being transmitted. Category theory replaced the isolated mathematical object with the structure-preserving morphism. In each case, a change in primitive concepts produced a corresponding expansion in explanatory power.

This monograph is motivated by the observation that many seemingly unrelated systems confront an identical structural problem. Biological organisms maintain their identity despite continuous molecular turnover. Scientific disciplines revise their theories while preserving standards of evidence. Distributed computer systems execute millions of concurrent state mutations without violating consistency constraints. Constitutional governments modify laws without abandoning the institutional structures upon which those laws depend. Artificial agents repeatedly generate candidate actions while attempting to avoid irreversible corruption of memory, execution state, or external environments. Although these domains employ distinct vocabularies, mathematical techniques, and explanatory traditions, they all depend upon a common requirement: they must continue changing without destroying the conditions that make continued change possible.

The central objective of this work is to construct a unified mathematical framework for studying such systems. Rather than beginning with optimization, utility, equilibrium, control, or computation, we begin with the more primitive notion of persistence. The fundamental question is not how rapidly a system can move through its state space, nor how efficiently it can optimize an objective function, but under what conditions it remains recognizably the same system while undergoing continual transformation.

This shift in perspective requires a corresponding shift in mathematical primitives. Traditional approaches frequently regard states as the primary objects of analysis and actions as operations that transform one state into another. Here, neither states nor actions occupy the foundational level. The primitive object of the theory is the transition itself. A transition is not merely an ordered pair of states but a structured relation whose validity depends upon the preservation of a family of invariants. Persistence is therefore not characterized by remaining stationary, nor by maximizing change, but by selecting those transitions that preserve the conditions required for future transitions.

The term \emph{admissible continuation} is introduced to describe this phenomenon. A continuation is admissible whenever the transition from one configuration to another preserves the invariant structures upon which the system's future existence depends. These invariant structures need not be identical across domains. In physics they may arise from conservation laws, symmetry principles, or dynamical constraints. In biology they appear as homeostatic envelopes and viability conditions. In formal systems they emerge as consistency, soundness, and type preservation. In computation they appear as correctness conditions, memory safety, and execution semantics. In institutions they arise from constitutions, procedural rules, and normative constraints. Despite their diversity, each serves the same abstract role: separating persistent trajectories from those that culminate in structural dissolution.

The development of the theory proceeds in several stages. The opening chapters investigate the origin of distinctions and the emergence of invariant families from the requirement that a system remain distinguishable from its environment. Subsequent chapters introduce admissible transition relations, an operator calculus governing proposal, evaluation, commitment, repair, maintenance, rollback, and refusal, and a generalized aggregation framework for combining multiple invariant valuations. The discussion then turns toward geometric representations of admissibility, including boundary structures, repair trajectories, hysteresis, and admissible path spaces. Later chapters examine recursive continuation, self-modifying systems, institutional persistence, biological adaptation, scientific methodology, and computational architectures as successive realizations of the same mathematical framework.

Although contemporary developments in artificial intelligence provide many illustrative examples throughout the text, the theory presented here is not intended as a work on machine learning or alignment alone. The recent emergence of agentic runtimes, verification layers, repair loops, transactional memory systems, and execution guardrails simply illustrates that practical engineering is increasingly rediscovering principles that appear to hold across a much broader class of persistent systems. The mathematical constructions developed in this monograph are therefore intended to stand independently of any particular technological implementation.

The exposition assumes familiarity with undergraduate mathematics, including linear algebra, real analysis, topology, probability, and abstract algebra, while drawing selectively upon differential geometry, dynamical systems, information theory, category theory, control theory, and theoretical computer science where appropriate. Wherever possible, new constructions are introduced from first principles and developed through formal definitions, propositions, lemmas, and theorems. Applications are presented only after the underlying mathematical structure has been established.

Several themes recur throughout the text. Distinction precedes information. Information precedes invariants. Invariants determine admissibility. Admissibility governs continuation. Continuation, rather than static equilibrium, becomes the primary mathematical expression of persistence. This hierarchy is not merely organizational; it reflects the logical dependence of each concept upon those that precede it.

The theory remains intentionally ambitious. It proposes that persistence itself may be treated as a mathematical object worthy of independent investigation, and that many phenomena traditionally studied within separate disciplines can be understood as different manifestations of a common theory of admissible continuation. Whether this proposal ultimately proves successful depends not upon philosophical appeal but upon its capacity to unify existing results, generate new mathematical questions, and illuminate persistent structures across disciplines that have historically developed in isolation.

The chapters that follow begin, therefore, not with intelligent agents, optimization algorithms, or computational architectures, but with the more elementary question from which every subsequent construction emerges.

\begin{center}
\emph{Why do some systems survive repeated change while others disappear?}
\end{center}

\chapter{The Persistence Problem}

\section{Persistence as a Mathematical Problem}

Every mathematical theory begins by specifying the class of objects it seeks to describe. Euclidean geometry studies spatial figures together with their invariant metric relations. Group theory studies algebraic structures whose operations satisfy prescribed axioms. Dynamical systems theory studies trajectories generated by evolution laws acting upon state spaces. Information theory studies the transmission and transformation of distinguishable configurations through noisy channels. Although these disciplines differ substantially in subject matter, they share a common methodological principle: the identification of an appropriate primitive object from which all subsequent constructions are derived.

The present work begins from a question whose simplicity conceals its remarkable generality.

How do systems remain themselves while continually changing?

The question appears in nearly every scientific discipline, although it is rarely posed in identical language. A biological organism replaces the overwhelming majority of its constituent molecules during its lifetime while nevertheless remaining recognizably the same organism. A language changes continuously across centuries without ceasing to be identified as the same language by historians and linguists. A software system undergoes thousands of revisions while preserving compatibility with previously constructed components. Scientific theories absorb new empirical observations, revise mathematical models, and discard obsolete hypotheses while maintaining an identifiable methodological continuity. Political constitutions are amended, interpreted, and reformed across generations without necessarily ceasing to define the same political order.

At first glance these examples appear to belong to unrelated disciplines whose mathematical foundations possess little common structure. Biology invokes metabolism, genetics, and natural selection. Computer science employs operational semantics, algorithms, and data structures. Physics appeals to conservation laws and differential equations. Jurisprudence relies upon procedural rules, constitutional interpretation, and institutional practice. Yet beneath these domain-specific mechanisms lies a common structural phenomenon. Each system repeatedly undergoes transformations while simultaneously preserving those properties necessary for further transformation.

This observation suggests that persistence itself may constitute a mathematical object deserving independent investigation.

Historically, persistence has often appeared indirectly through concepts such as stability, equilibrium, robustness, resilience, homeostasis, fault tolerance, or conservation. Although each captures an important aspect of long-term behavior, none provides a sufficiently general framework for describing systems whose defining characteristic is continual transformation rather than static maintenance. Equilibrium emphasizes immobility or balanced dynamics. Robustness measures resistance to perturbation. Stability concerns asymptotic behavior under local disturbances. Homeostasis applies primarily to biological regulation. None of these concepts directly addresses the more general question of which transformations preserve the possibility of future transformations.

The distinction is subtle but fundamental.

A system may remain perfectly stable while possessing no capacity for adaptation. Conversely, a highly adaptive system may rapidly destroy its own structural identity if unrestricted change is permitted. Persistence therefore occupies a conceptual position between rigidity and unconstrained variability. A completely rigid system cannot respond to environmental change and eventually fails through inflexibility. A completely unconstrained system dissolves because every perturbation is equally permissible. Persistent systems occupy an intermediate regime in which change is permitted only insofar as it preserves the conditions necessary for subsequent change.

This observation motivates the central thesis of the present work.

Persistence is fundamentally a property of constrained continuation.

The theory developed throughout this monograph therefore seeks to identify the mathematical conditions under which sequences of transformations preserve the structures that define the continuing existence of a system.

\section{The Inadequacy of State-Centered Descriptions}

Most mathematical treatments of dynamical systems begin by specifying a state space
\[
\mathcal{S},
\]
whose elements represent complete descriptions of the system at individual moments in time. Evolution is then represented by a transition function, differential equation, stochastic kernel, or control policy mapping one state into another.

This formulation has proven enormously successful across physics, engineering, economics, and computer science. Nevertheless, it embeds a philosophical assumption that will be questioned throughout this monograph.

The assumption is that the state constitutes the primary mathematical object.

Under this viewpoint, transitions appear merely as secondary constructions relating already existing states. Questions concerning safety, optimization, stability, and control are therefore typically formulated as properties of subsets of the state space. One seeks safe regions, equilibrium regions, attractor basins, feasible regions, reachable sets, or optimal configurations.

While these constructions remain valuable, they obscure an important observation.

No isolated state possesses intrinsic mathematical significance independent of the transitions through which it is reached and the transitions that remain possible thereafter.

To illustrate this point, consider a computational process whose memory contains no data. Such a configuration may represent successful initialization, catastrophic memory corruption, intentional garbage collection, or complete hardware failure. The state itself is insufficient to distinguish these possibilities. Its interpretation depends entirely upon the trajectory through which it was produced and the admissible continuations available from it.

A similar phenomenon appears throughout the sciences.

A particular chemical concentration may represent healthy metabolism in one organism and terminal pathology in another. An identical logical proposition may constitute a theorem in one formal system and a contradiction in another. An identical institutional decision may preserve constitutional order under one procedural history while violating it under another.

The mathematical meaning of a configuration therefore depends not only upon its internal structure but upon its relationship to preceding and succeeding configurations.

This dependence suggests that persistence cannot be understood solely through properties of individual states.

Instead, persistence arises through relationships among successive configurations.

The primitive mathematical object of the present theory will therefore not be the state itself, but the transition relating one state to another.

The consequences of this seemingly modest change in perspective will occupy the remainder of this volume.

\section{The Continuation Principle}

Suppose that a system occupies some configuration
\[
s_0 \in \mathcal{S}.
\]

Traditional formulations ask which states can be reached from
\[
s_0.
\]

This question concerns reachability.

The present work asks a different question.

Among all reachable configurations, which permit continued existence of the system as the same identifiable entity?

This distinction immediately separates two notions that are frequently conflated.

The first concerns possibility.

The second concerns persistence.

A transition may be physically possible while destroying every structural condition required for future operation. Likewise, a transition may maximize an immediate objective function while irreversibly eliminating future alternatives. Neither possibility nor optimization alone guarantees persistence.

Instead, persistence requires the existence of an indefinitely extensible sequence of structurally acceptable transformations.

This observation motivates what will later become one of the fundamental principles of the theory.

A system persists throughout an interval precisely when there exists a continuation whose successive transformations preserve the structures upon which further continuation depends.

The remainder of this monograph develops the mathematical language necessary to make this statement precise.

Rather than beginning with optimization objectives, reward functions, utility measures, or action spaces, we begin with the more elementary concepts of distinction, invariance, and continuation.

Only after these foundations have been established shall we introduce the formal notion of admissibility and construct the operator calculus governing persistent systems.

The progression is deliberate.

Distinctions give rise to information.

Information gives rise to invariant structures.

Invariant structures determine which continuations remain viable.

Viable continuations generate persistence.

Persistence, rather than static equilibrium, becomes the fundamental mathematical object of investigation.

\section{Distinction and the Emergence of Systems}

The preceding discussion deliberately avoided specifying what constitutes a system. This omission was not accidental. Before a theory can investigate the persistence of systems, it must first explain what distinguishes a system from its surroundings. Without such a distinction, persistence becomes meaningless, for there exists no entity whose continued existence may be discussed.

The notion of distinction therefore precedes every subsequent construction in this monograph.

One cannot preserve an object that has never been distinguished.

One cannot describe a trajectory before identifying the entity that traverses it.

One cannot formulate admissibility before determining what is to be preserved.

Historically, distinctions have appeared in many different mathematical forms. In topology they arise as subsets endowed with boundaries. In measure theory they emerge as measurable regions separated by sigma-algebras. In graph theory they appear as induced subgraphs separated by cut sets. In theoretical computer science they arise as namespaces, memory regions, and execution contexts. In information theory they appear through distinguishable messages, symbols, and codewords. Although these constructions differ technically, each performs the same conceptual operation: the separation of one collection of configurations from another.

The present work adopts this operation as its logical starting point.

We therefore refrain, at least initially, from assuming the existence of states, trajectories, or dynamics. Instead, we begin with a universe of possible configurations and ask how an identifiable system can emerge within it.

Let
\[
\Omega
\]
denote an ambient configuration space whose elements represent all configurations relevant to the discussion. No assumptions are yet imposed concerning topology, differentiability, algebraic structure, or metric properties. The space may be finite or infinite, discrete or continuous, deterministic or stochastic. At this stage it serves only as the background within which distinctions may be drawn.

The emergence of a system requires the identification of some subset whose behavior is to be studied independently of its surroundings. This act is more fundamental than assigning coordinates or specifying dynamics. It establishes the object of investigation itself.

The distinction need not correspond to a physical object. It may represent an organism, a software process, a mathematical proof, a legal institution, a scientific theory, a distributed protocol, or an entire civilization. The theory intentionally remains agnostic regarding ontology. It concerns only the structural consequences of distinguishing one region of configuration space from another.

The consequences of this apparently simple operation are profound.

Once a distinction has been established, every subsequent transformation acquires two possible interpretations. It may preserve the distinction, allowing the identified system to remain recognizable across time, or it may destroy the distinction by erasing the boundary separating the system from its environment.

Persistence therefore becomes meaningful only after distinction has been introduced.

The theory consequently proceeds in the opposite direction from many traditional approaches. Rather than assuming systems and asking how they evolve, we first investigate how systems become mathematically distinguishable. Evolution is then understood as the preservation or destruction of those distinctions through successive transformations.

\section{The Primitive Distinction}

We now introduce the first formal definition.

\begin{definition}[Primitive Distinction]
Let
\[
\Omega
\]
be a non-empty configuration space.

A \emph{distinction} is an ordered pair
\[
D=(S,S^{c}),
\]
where
\[
S\subseteq\Omega,
\]
\[
S^{c}=\Omega\setminus S,
\]
and both
\[
S\neq\varnothing,
\qquad
S^{c}\neq\varnothing.
\]

The subset
\[
S
\]
will be called the \emph{system region}, while
\[
S^{c}
\]
will be called the \emph{environment}.
\end{definition}

This definition is intentionally elementary.

Nothing has yet been said regarding dynamics, geometry, probability, or information. The definition merely formalizes the existence of an identifiable entity whose behavior may subsequently be investigated.

The importance of the distinction lies not in the subsets themselves but in the existence of the separating boundary that they induce.

Whenever a distinction exists, there is necessarily some criterion by which configurations lying within the system region are recognized as belonging to the system rather than the environment. The precise mathematical realization of this criterion depends upon the surrounding structure. In topology it becomes a boundary. In graph theory it becomes a cut. In information theory it becomes distinguishability. In category theory it may appear through objects and morphisms. In computation it becomes namespace isolation or memory protection.

The abstract theory developed here deliberately postpones these specializations.

Instead, we retain only the universal observation that every identifiable system depends upon the preservation of some distinction separating internal configurations from external ones.

This observation leads naturally to the next question.

What must remain true for the distinction itself to survive repeated transformation?

\section{Persistence Requires Boundary Preservation}

Suppose that a distinguished system undergoes a transformation.

If the transformation preserves the separation between the system and its environment, then subsequent observers continue to recognize the transformed configuration as belonging to the same evolving entity.

If, however, the transformation destroys every meaningful distinction between the system and its surroundings, then the system has not merely changed.

It has ceased to exist as a distinguishable object.

This statement should not be interpreted philosophically but mathematically.

Whenever the characteristic properties defining the distinguished subset become indistinguishable from those of the surrounding space, the distinction introduced in the previous section collapses. The mathematical object originally selected for investigation can no longer be uniquely identified.

Persistence therefore cannot simply mean remaining close to an initial configuration.

Nor can it mean maximizing similarity according to some arbitrary metric.

Instead, persistence concerns the preservation of those structural properties that continue to separate the evolving system from its environment.

The distinction itself becomes the first invariant.

Everything else follows from it.

Information requires distinguishable alternatives.

Memory requires distinguishable historical states.

Observation requires distinguishable outcomes.

Control requires distinguishable futures.

Proof requires distinguishable propositions.

Computation requires distinguishable machine configurations.

Without distinction, every subsequent mathematical construction degenerates into undifferentiated noise.

This observation motivates the first informal principle of the theory.

\begin{quote}
A system persists only insofar as the distinctions defining that system remain recoverable throughout its evolution.
\end{quote}

The remainder of this chapter is devoted to transforming this qualitative principle into precise mathematical language.

To accomplish this, we must next investigate how distinctions generate information and why the preservation of distinguishability naturally gives rise to invariant structures.

\section{Distinction Generates Information}

The preceding sections established that every mathematical investigation of persistence presupposes the existence of a distinction. Before one may ask whether a system continues through time, one must first specify what is to be distinguished from its surroundings. The primitive distinction introduced previously therefore occupies a position logically prior to dynamics, optimization, computation, and even information itself.

The next objective is to demonstrate that information is not an independent primitive. Rather, it emerges as a quantitative expression of distinguishability.

This observation has appeared in many forms throughout the history of mathematics and science. Shannon demonstrated that information measures the reduction of uncertainty associated with distinguishable alternatives. Statistical mechanics relates entropy to the cardinality of distinguishable microscopic configurations compatible with a macroscopic state. Algorithmic information theory studies the shortest effective description capable of distinguishing one object from all others. Model theory distinguishes structures through satisfaction of logical formulae. Across these diverse formulations, the common mathematical ingredient is not probability, compression, or computation in isolation. It is the existence of alternatives that can be distinguished.

Without distinction, there exists no uncertainty.

Without uncertainty, there exists no information.

Suppose that the ambient configuration space consists of a single element,
\[
\Omega=\{\omega\}.
\]

No observation performed upon such a universe can produce information, because every possible observation necessarily yields the same result. Every measurement, regardless of precision or methodology, returns the identical configuration. There exists nothing to discriminate, classify, or encode.

Now suppose instead that
\[
|\Omega|>1.
\]

The introduction of multiple distinguishable configurations immediately creates the possibility of observation. A measurement now partitions the configuration space into alternatives. The outcome of the observation identifies which element, or more generally which subset, corresponds to reality. Information therefore appears not because probabilities have been assigned, but because distinctions exist.

This logical relationship is more fundamental than any particular quantitative definition.

Probability distributions measure uncertainty over distinctions.

Entropy measures the expected uncertainty remaining before distinctions are resolved.

Mutual information measures the preservation of distinctions between correlated systems.

Compression exploits regularities among distinctions.

Error-correcting codes preserve distinctions against noise.

The common mathematical substrate beneath each construction is distinguishability.

Consequently, the preservation of information should not be understood as an isolated engineering objective. It is a direct consequence of preserving the distinctions that define the system itself.

This viewpoint reverses the historical development of several disciplines.

Information theory traditionally begins by assuming distinguishable symbols and proceeds to study efficient communication among them. The present theory instead asks why distinguishable symbols exist in the first place. Their existence derives from the more primitive act of partitioning an ambient configuration space into identifiable regions.

Information is therefore secondary.

Distinction is primary.

This hierarchy is reflected throughout the remainder of the monograph.

Whenever information is lost, some distinction has become less recoverable.

Whenever entropy increases, distinguishability has decreased relative to the observational framework under consideration.

Whenever repair succeeds, distinguishability has been restored.

Whenever persistence fails, the defining distinctions of the system have become irrecoverable.

Thus the mathematics of persistence naturally inherits the mathematics of information through the preservation of distinction.

\section{Recoverability and the Meaning of Persistence}

The preceding discussion suggests that persistence cannot simply mean the continued existence of physical matter, computational memory, or symbolic representations. Such quantities may remain present while nevertheless losing the structural organization that originally distinguished the system from its surroundings.

A computer file whose bits have been randomly permuted still occupies storage media, yet no longer preserves the information originally encoded within it.

A biological organism whose molecular constituents remain present after death nevertheless ceases to function as the organism previously identified.

A legal institution may retain its buildings, personnel, and archives while losing the constitutional relationships that originally defined its authority.

These examples illustrate a common phenomenon.

Persistence concerns not the survival of material substrate but the recoverability of organized distinction.

The notion of recoverability deserves careful emphasis.

A distinction need not remain perfectly unchanged in order to persist. Biological cells continually replace proteins. Software systems modify source code. Languages alter grammar and vocabulary. Scientific theories undergo repeated revision. In each case the observable configuration changes substantially while remaining sufficiently organized that earlier and later states may still be recognized as belonging to the same evolving entity.

Persistence therefore admits gradual transformation.

It excludes only those transformations that render the defining distinctions irrecoverable.

The concept of recoverability consequently occupies an intermediate position between exact identity and unrestricted change.

Exact identity requires no alteration whatsoever.

Complete dissolution destroys every recognizable relationship between past and future configurations.

Persistent systems inhabit the broad mathematical region between these extremes.

This observation suggests that persistence should not be represented by equality.

Rather, persistence concerns the existence of sufficiently structured transformations allowing one configuration to be related to another without destroying the distinctions upon which identification depends.

The mathematical object of primary interest is therefore not the state itself, nor the equality relation between states, but the family of transformations under which distinguishability remains recoverable.

This realization marks the transition from static mathematics to relational mathematics.

The subsequent chapters will formalize this observation by introducing admissible transition relations, invariant families, and continuation operators. Before doing so, however, one further conceptual step remains necessary.

We must distinguish between changes that merely alter observable properties of a system and changes that destroy the structural conditions required for the system to remain identifiable.

The mathematical language for expressing this distinction will be provided by the concept of invariants.

\section{Toward Invariant Structure}

Every persistent system exhibits two complementary aspects.

On the one hand, it changes continuously. Components are replaced, parameters drift, memories accumulate, environments fluctuate, and trajectories evolve through time. Without such change, the system would be incapable of adaptation and would eventually fail under environmental variation.

On the other hand, some structural properties remain sufficiently stable that successive configurations continue to be recognized as manifestations of the same underlying entity. These preserved properties need not be immutable, nor must they remain numerically constant. They need only survive in a manner that maintains the recoverability of the defining distinctions introduced earlier in this chapter.

The existence of simultaneously changing and preserved structure is the defining characteristic of persistence.

This observation motivates the introduction of invariants.

Informally, an invariant is any property whose preservation is necessary for the continued recoverability of the distinctions defining a system. Different domains instantiate different invariant families. Physical systems preserve conservation laws and admissible dynamical trajectories. Biological systems preserve metabolic viability and homeostatic regulation. Formal systems preserve logical consistency. Computational systems preserve memory safety and semantic correctness. Institutions preserve procedural legitimacy. Scientific communities preserve evidential standards.

Although these examples differ substantially in implementation, they perform the same abstract mathematical role.

They constrain the permissible transformations of the system.

In the chapters that follow, invariants will become the central organizing principle of the theory. Rather than beginning with objectives to optimize or actions to perform, we shall begin with the invariant structures that must survive any admissible continuation. Transitions will then be evaluated not according to their immediate utility alone, but according to their capacity to preserve the invariant conditions required for continued evolution.

The progression developed throughout this opening chapter may therefore be summarized as a sequence of logical dependence.

Distinction gives rise to information.

Information makes recoverability meaningful.

Recoverability requires the preservation of invariant structure.

Invariant structure determines which transformations may legitimately occur.

Those transformations, and only those transformations, constitute the admissible continuations from which persistent histories are assembled.

The next chapter develops this idea formally by investigating the mathematical origin of invariants themselves. Rather than treating invariants as externally imposed constraints or domain-specific axioms, we shall derive them systematically from the requirement that distinctions remain recoverable across successive transformations.

\chapter{The Origin of Invariants}

\section{The Derivation Problem}

The preceding chapter established a hierarchy of logical dependence beginning with distinction and culminating in persistence. We argued that persistence cannot be understood as the preservation of static configurations, but rather as the preservation of those structural conditions that permit distinguishable systems to continue through successive transformations. Information emerged from distinction, recoverability emerged from information, and the notion of invariant structure appeared as the mathematical bridge connecting recoverability to persistence.

At this point the theory encounters its first genuinely foundational question.

What is an invariant?

More importantly, where do invariants come from?

These questions are considerably more subtle than they initially appear. Across mathematics, physics, computer science, and engineering, invariant quantities are introduced almost reflexively. One assumes conservation of energy, preservation of type correctness, maintenance of referential integrity, constitutional constraints, memory safety, or consistency conditions, and then proceeds to study the consequences of those assumptions. The invariants themselves are seldom derived. They are accepted as primitive features of the domain under consideration.

For many specialized theories this approach is entirely appropriate. Classical mechanics does not attempt to derive conservation of momentum from more elementary mathematical principles. Programming language semantics does not ordinarily derive type safety from epistemic distinguishability. Constitutional law does not prove the necessity of due process from first principles of information theory. Each discipline inherits its own collection of foundational constraints and develops increasingly sophisticated consequences from them.

The present work adopts a different objective.

If persistence is to be treated as a genuinely universal mathematical phenomenon rather than as a family of unrelated engineering heuristics, then invariant structure cannot simply be assumed independently within every discipline. Instead, there must exist a more primitive principle from which the necessity of invariants itself emerges.

The purpose of this chapter is to identify that principle.

The argument developed throughout the remainder of this work is that invariants do not originate as arbitrary restrictions imposed upon otherwise unrestricted systems. Nor are they merely convenient conventions adopted for computational efficiency or human convenience. Rather, they arise inevitably from the requirement that distinguishable systems remain distinguishable throughout their evolution.

The distinction is essential.

Constraint-based theories often begin by specifying a set of allowable transitions and then asking how systems behave within those constraints. Such an approach treats constraints as external objects. They appear as rules imposed from outside the mathematical system itself.

The present theory proceeds in the opposite direction.

We begin with the existence of distinguishable entities and ask what structural conditions must necessarily hold if those entities are to remain identifiable through time.

The resulting invariants are therefore not external limitations.

They are internal consequences of persistence.

This reversal transforms the conceptual role of invariants.

Instead of restricting behavior that would otherwise be permissible, they define the mathematical conditions under which the notion of a continuing system remains meaningful.

Without invariants there is no persistence.

Without persistence there is no identifiable system.

Without identifiable systems there exists no meaningful object whose evolution may be studied.

Thus the derivation of invariants is not an auxiliary problem.

It is the logical foundation upon which the remainder of the theory depends.

\section{Against Arbitrary Constraint Systems}

Before constructing a positive theory of invariant emergence, it is instructive to examine the limitations of arbitrary constraint systems.

Suppose that one begins with an arbitrary collection of functions

\[
\Phi_1,\Phi_2,\ldots,\Phi_n,
\]

defined over a state space or transition space.

Nothing prevents us from declaring these functions to represent admissibility conditions.

Likewise, nothing prevents us from constructing evaluation operators, optimization procedures, or verification algorithms whose sole purpose is to preserve them.

Such a formalism may be internally consistent.

It may even prove useful for particular engineering applications.

Nevertheless, it possesses an important philosophical weakness.

The origin of the constraints remains unexplained.

If the functions are introduced solely by stipulation, then any alternative collection of functions enjoys equal mathematical legitimacy. One investigator may define admissibility through conservation of computational resources. Another may define it through aesthetic preferences. A third may define it through arbitrary symbolic patterns having no relationship whatsoever to persistence.

Without an independent derivation, the mathematical formalism possesses no criterion by which one collection of constraints can be distinguished from another.

The resulting theory therefore explains only how to preserve constraints once they have been specified.

It does not explain why those constraints deserve preservation.

This limitation is not merely philosophical.

It has practical consequences.

Modern engineering systems frequently accumulate constraints over time in an ad hoc fashion. Safety rules, software policies, verification conditions, organizational procedures, and regulatory requirements are continually appended as new problems arise. Although each additional constraint may solve an immediate difficulty, the growing collection often lacks any unifying mathematical principle.

Consequently, systems become increasingly difficult to reason about.

Interactions among independently introduced constraints generate unexpected behavior.

Repair procedures become increasingly complex.

Optimization becomes progressively more expensive.

Eventually the system evolves into a collection of exceptions rather than an expression of coherent mathematical structure.

A theory of persistence should avoid this difficulty.

Rather than accumulating arbitrary constraints, it should identify the small collection of primitive principles from which necessary constraints naturally emerge.

The search for invariant origins is therefore simultaneously a search for mathematical economy.

If successful, the resulting theory replaces an indefinitely expanding catalogue of domain-specific restrictions with a comparatively small family of structurally necessary invariants.

Such economy has historically characterized the most successful mathematical theories.

Maxwell unified electricity and magnetism.

Shannon unified communication and coding.

Category theory unified diverse branches of abstract algebra.

The aspiration of the present work is considerably more modest but structurally similar.

Instead of unifying physical forces or algebraic constructions, it seeks to unify the mathematical conditions under which distinguishable systems remain capable of continued existence.

\section{Persistence Before Mechanics}

One consequence of this perspective deserves particular emphasis.

Traditional mathematical theories frequently begin by specifying the mechanics governing a system.

One defines equations of motion, transition probabilities, production rules, algorithms, or control policies.

Only afterward does one ask whether the resulting dynamics preserve particular quantities.

The present theory reverses this order.

Persistence precedes mechanics.

This statement should not be misunderstood.

It does not imply that persistence causes physical dynamics or replaces ordinary scientific explanation.

Rather, it reflects the logical ordering of concepts required for a general mathematical theory.

Before specifying how a system evolves, one must first specify what it means for the evolving object to remain the same identifiable system.

Only after this criterion has been established can one meaningfully investigate equations governing its evolution.

Consequently, the derivation of invariant structure is intentionally independent of any particular physical law, computational architecture, or biological mechanism.

Whether the underlying dynamics arise from differential equations, stochastic processes, symbolic rewriting systems, distributed consensus protocols, or evolutionary selection is irrelevant at this stage.

The question remains identical in every case.

What properties must survive successive transformations if the system is to remain distinguishable from its surroundings?

Everything else follows from the answer to this question.

The remainder of this chapter develops that answer by returning once again to the primitive distinction introduced previously.

The goal is to show that invariants emerge naturally as mathematical expressions of boundary preservation rather than as externally imposed restrictions.

In doing so, we shall obtain a domain-independent framework capable of accommodating physical conservation laws, informational integrity, biological homeostasis, logical consistency, computational correctness, and institutional continuity as different realizations of the same underlying principle.

\section{Boundary Preservation as the Source of Invariants}

The conclusion of the previous chapter was that persistence cannot be defined in terms of static identity. Systems survive not because they remain unchanged, but because successive transformations preserve the distinctions that permit the system to remain identifiable throughout its evolution. The present chapter has further argued that the invariant structures governing these transformations should not be regarded as arbitrary restrictions imposed upon an otherwise unconstrained dynamics. Instead, they must arise as necessary consequences of the requirement that distinctions remain recoverable through time.

We are therefore led to a more primitive question.

What does it mean for a distinction itself to persist?

To answer this question we return to the primitive distinction introduced previously. Let

\[
D=(S,S^{c}),
\]

where

\[
S\subseteq\Omega,
\qquad
S^{c}=\Omega\setminus S.
\]

The distinction itself is not the subset \(S\) alone. Nor is it merely the complement \(S^{c}\). Rather, it is the relational separation between them. A system exists only insofar as configurations belonging to \(S\) remain mathematically distinguishable from those belonging to \(S^{c}\).

Suppose now that the system evolves through some transformation

\[
\tau:S\longrightarrow S'.
\]

Several possibilities immediately arise.

The transformed region \(S'\) may remain entirely distinguishable from its surrounding environment. In this case the original distinction survives, although the internal organization of the system may have changed substantially.

Alternatively, the transformation may deform the system until portions of its boundary become ambiguous. Some configurations may no longer admit a unique classification as belonging either to the system or to the environment. The distinction weakens, although it has not yet disappeared entirely.

Finally, the transformation may destroy the distinction altogether. The transformed region becomes indistinguishable from the ambient configuration space, and the mathematical object originally identified as the system can no longer be recovered.

These three possibilities correspond respectively to persistence, degradation, and dissolution.

The crucial observation is that none of them depends primarily upon the internal coordinates of the system.

Instead, each depends upon the behavior of the boundary separating the system from its surroundings.

The preservation of this boundary therefore becomes more fundamental than the preservation of any particular state variable.

This represents a significant conceptual shift.

Traditional conservation laws preserve quantities.

The present framework preserves distinctions.

Quantities become important only insofar as they participate in maintaining those distinctions.

To illustrate the point, consider a living cell.

No individual molecule remains permanently inside the organism. Water molecules, proteins, ions, and metabolic intermediates are continually exchanged with the surrounding environment. If persistence required preservation of constituent particles, every organism would cease to exist within moments.

Instead, biological persistence depends upon preservation of organizational boundaries.

The membrane remains distinguishable.

The regulatory architecture remains distinguishable.

The metabolic network remains distinguishable.

Material changes continuously.

The distinction survives.

An analogous phenomenon appears in computation.

The individual voltages stored within memory cells fluctuate continually. Processor registers are repeatedly overwritten. Temporary variables are created and destroyed millions of times each second. Nevertheless, the executing process persists because the logical boundaries separating executable state, protected memory, and external processes remain intact.

Again, persistence is achieved through preservation of distinction rather than preservation of individual components.

Scientific theories exhibit the same structural behavior.

Individual hypotheses are continually revised. Mathematical techniques evolve. Experimental methods improve. Entire explanatory frameworks may undergo radical reconstruction. Yet the scientific enterprise itself persists because the distinction between empirical warrant and unsupported speculation continues to organize the process of theory revision.

Across these examples the common mathematical structure is unmistakable.

Persistence is governed not by immutable substance but by the preservation of organizational boundaries.

The notion of boundary employed here should be interpreted broadly.

In topological spaces it corresponds to the ordinary topological boundary.

In graph-theoretic systems it appears as separating cut sets.

In information theory it becomes distinguishability between message ensembles.

In computation it manifests as namespace isolation, memory protection, or interface separation.

In institutions it corresponds to procedural constraints separating legitimate from illegitimate transitions.

Although these realizations differ technically, they perform the identical mathematical function.

They preserve distinguishability.

Consequently, the first universal invariant is not conservation of energy, momentum, probability, information, or computational correctness.

It is the preservation of distinction itself.

Every subsequent invariant may be understood as a refinement of this more primitive requirement.

Physical conservation laws preserve distinguishable physical trajectories.

Logical consistency preserves distinguishable propositions.

Biological regulation preserves distinguishable living organization.

Computational correctness preserves distinguishable execution histories.

Institutional continuity preserves distinguishable legal authority.

Each domain introduces its own specialized mathematical machinery, yet every one of these structures ultimately serves the same abstract purpose: preventing the collapse of the distinction upon which the existence of the system depends.

This observation suggests that invariant structure is neither accidental nor conventional.

It is inevitable.

Whenever a distinguishable system persists through successive transformations, there necessarily exists some family of properties whose preservation prevents the collapse of the defining distinction.

Conversely, whenever every such property is destroyed, the distinction itself becomes unrecoverable, and the system dissolves into its surrounding environment.

We therefore arrive at the central thesis of this chapter.

Invariant structures are not external constraints imposed upon persistent systems.

They are the mathematical expression of boundary preservation.

The remainder of this chapter is devoted to formalizing this claim. Beginning with the preservation of distinction, we shall derive progressively richer families of invariants governing information, dynamics, computation, biology, logic, and institutions. These domain-specific invariant families will ultimately be unified under a common mathematical framework in which admissibility is defined not by arbitrary policy but by the preservation of the boundaries that make persistence itself possible.

\section{The Axiom of Distinguishability Preservation}

The preceding discussion has remained intentionally informal. Although the language of boundaries, recoverability, and persistence provides substantial intuition, a mathematical theory cannot ultimately rest upon intuition alone. The transition from conceptual exposition to formal development therefore begins here.

The objective of this section is not to introduce the full machinery of admissibility. Rather, it is to identify the weakest mathematical principle from which every subsequent construction may be derived.

The principle is remarkably simple.

A persistent transformation must preserve the distinguishability of the system it transforms.

This statement should be interpreted carefully. It does not require that every observable feature remain unchanged. Persistent systems continually modify internal structure, replace components, accumulate history, and interact with their environments. Indeed, complete absence of change would preclude adaptation and render persistence impossible under varying external conditions.

Instead, the principle requires that the transformation preserve sufficient structure that the evolving system remains mathematically identifiable as the continuation of the system that preceded it.

The emphasis therefore lies upon preservation of distinguishability rather than preservation of identity.

To formalize this idea, let

\[
D=(S,S^{c})
\]

be a primitive distinction within an ambient configuration space

\[
\Omega.
\]

Suppose that

\[
\tau:\Omega\rightarrow\Omega
\]

represents an admissible candidate transformation.

The image

\[
\tau(S)
\]

need not coincide with

\[
S.
\]

Indeed, requiring equality would reduce persistence to stasis.

Instead, the transformed region need only remain distinguishable from

\[
\tau(S^{c}),
\]

the transformed environment.

Persistence therefore concerns the preservation of separation rather than geometric coincidence.

This observation motivates the first axiom.

\begin{definition}[Distinguishability Preservation]
Let
\[
D=(S,S^{c})
\]
be a distinction on
\[
\Omega.
\]

A transformation
\[
\tau:\Omega\rightarrow\Omega
\]
is said to preserve the distinction \(D\) if there exists a criterion by which every element of
\[
\tau(S)
\]
remains distinguishable from every element of
\[
\tau(S^{c}).
\]

When this condition holds, the distinction is said to survive the transformation.
\end{definition}

Several remarks are immediately in order.

First, the definition deliberately avoids specifying the nature of the distinguishing criterion. In a metric space it may arise from positive distance. In topology it may arise from neighborhood separation. In graph theory it may arise from edge cuts. In logic it may arise from satisfaction classes. In computation it may arise from namespace isolation or execution semantics.

The abstract theory requires only that such a criterion exist.

Second, distinguishability need not be absolute.

Many systems exhibit partial degradation while remaining recoverable. Biological organisms accumulate mutations without immediate loss of identity. Software systems accumulate technical debt while remaining executable. Institutions undergo procedural irregularities without immediate constitutional collapse.

The definition therefore establishes only the binary notion of survival.

Later chapters will introduce quantitative measures describing the degree to which distinctions are preserved, weakened, or repaired.

Finally, the definition is independent of time.

The transformation may represent an infinitesimal evolution generated by a differential equation, a discrete algorithmic update, a stochastic transition, a legal amendment, or an evolutionary generation.

Persistence concerns the structural properties of the transformation itself rather than the temporal mechanism producing it.

The definition therefore applies equally to continuous and discrete systems.

Having established the fundamental preservation principle, we may now investigate its immediate mathematical consequences.

\begin{proposition}[Persistence Requires Distinguishability]
Suppose a system evolves through a finite sequence of transformations

\[
\tau_1,\tau_2,\ldots,\tau_n.
\]

If any transformation fails to preserve the defining distinction of the system, then every subsequent configuration fails to represent a continuation of the original distinguished system.
\end{proposition}

\begin{proof}

Let

\[
D_0=(S_0,S_0^{c})
\]

denote the initial distinction.

Assume that for some index

\[
k,
\]

the transformation

\[
\tau_k
\]

fails to preserve distinguishability.

Then, by definition,

\[
\tau_k(S_{k-1})
\]

cannot be uniquely distinguished from

\[
\tau_k(S_{k-1}^{c}).
\]

Consequently, the image

\[
S_k
\]

no longer possesses a well-defined boundary separating the system from its environment.

Every subsequent transformation

\[
\tau_{k+1},\ldots,\tau_n
\]

acts upon an object that is no longer mathematically identifiable as the original distinguished system.

These later transformations may produce additional structure, but they cannot restore continuity with the original system merely through continued evolution. Once distinguishability has been completely lost, subsequent evolution lacks a uniquely defined object whose persistence may be asserted.

Therefore, preservation of distinction is a necessary condition for persistence across every finite evolution sequence.

\end{proof}

The proposition establishes an important asymmetry.

The preservation of distinction guarantees neither optimality nor stability.

A system may remain perfectly distinguishable while becoming increasingly inefficient, fragile, or maladapted.

Conversely, destruction of distinction immediately terminates persistence regardless of any improvements achieved according to secondary objectives.

Persistence therefore occupies a more primitive logical position than optimization.

Optimization presupposes the continued existence of an identifiable system capable of realizing improved states.

Persistence merely requires that such a system continue to exist at all.

This asymmetry explains why optimization problems frequently become meaningless after catastrophic structural failure. A destroyed computation cannot maximize utility. A dissolved institution cannot enact policy. An extinct biological lineage cannot adapt further.

Optimization belongs to surviving systems.

Persistence determines whether there remains any system to optimize.

The remainder of this chapter develops this observation into a hierarchy of invariant families. We shall show that every invariant introduced later in this monograph may be interpreted as preserving distinguishability with respect to some particular observational framework. Conservation laws preserve distinguishable physical histories. Logical consistency preserves distinguishable formal derivations. Computational correctness preserves distinguishable execution trajectories. Biological regulation preserves distinguishable living organization.

Thus, although these invariants appear disparate when studied within their individual disciplines, they become unified when viewed through the more primitive mathematics of distinction preservation.

\section{Observation and Relative Distinguishability}

The preceding section introduced the preservation of distinguishability as the first structural requirement for persistence. The discussion, however, deliberately left one important question unresolved. Distinguishable to whom, or with respect to what?

This question is unavoidable. Distinguishability is never an intrinsic property of an isolated configuration considered in complete abstraction. Rather, it arises through some process capable of extracting information from that configuration. Two physical systems may be distinguishable under microscopic observation while appearing identical at macroscopic scales. Two computational states may differ internally while producing identical outputs. Two mathematical objects may admit distinct constructions yet satisfy precisely the same observable properties within a given logical language.

The existence of such examples demonstrates that distinguishability is fundamentally relational.

The present theory therefore introduces observation as an abstract mathematical operation. The purpose of this construction is not to model human perception, experimental apparatus, or computational sensors in particular. Instead, it provides a unified framework within which every process of discrimination may be represented.

\begin{definition}[Observation Map]

Let
\[
\Omega
\]
be a configuration space.

An \emph{observation map} is a function

\[
O:\Omega\rightarrow\mathcal{Y},
\]

where
\[
\mathcal{Y}
\]
is an arbitrary observation space.

The value
\[
O(x)
\]
is called the \emph{observation} of the configuration
\[
x\in\Omega.
\]

\end{definition}

The definition intentionally imposes almost no structure upon the codomain. The observation space may consist of numerical measurements, symbolic descriptions, probability distributions, logical propositions, feature vectors, graph invariants, equivalence classes, or any other mathematical object capable of representing observable information.

The flexibility of this definition reflects the domain-independent ambition of the theory. A thermometer, a camera, a theorem prover, a compiler, a scientific instrument, a biological receptor, and a software verification routine all become special cases of observation maps.

The differences lie entirely in the choice of codomain and the mathematical properties of the mapping itself.

Once an observation map has been specified, distinguishability acquires a precise meaning.

\begin{definition}[Relative Distinguishability]

Let

\[
O:\Omega\rightarrow\mathcal{Y}
\]

be an observation map.

Two configurations

\[
x,y\in\Omega
\]

are said to be \emph{distinguishable with respect to}
\[
O
\]

whenever

\[
O(x)\neq O(y).
\]

Otherwise they are said to be observationally indistinguishable.

\end{definition}

Several important consequences follow immediately.

First, distinguishability is observer-relative.

This statement should not be interpreted psychologically. The "observer" is simply the mathematical map \(O\). Different maps may preserve different structural features of the same configuration space.

For example, consider two digital images differing only in one least significant bit. A cryptographic hash distinguishes them immediately, whereas the human visual system typically does not. Conversely, two photographs captured under different lighting conditions may produce substantially different pixel arrays while remaining immediately recognizable as depicting the same scene.

Neither observation is incorrect.

They merely arise from different observation maps.

The same phenomenon appears throughout mathematics.

Two graphs may be non-isomorphic while sharing identical degree sequences.

Two manifolds may possess identical homology groups while differing in smooth structure.

Two computer programs may have entirely different implementations while exhibiting identical observable behavior.

Observation therefore induces an equivalence relation upon the configuration space.

\begin{definition}[Observational Equivalence]

Given an observation map

\[
O:\Omega\rightarrow\mathcal{Y},
\]

define

\[
x\sim_O y
\]

whenever

\[
O(x)=O(y).
\]

The equivalence classes of

\[
\sim_O
\]

are called the \emph{observational classes} induced by
\[
O.
\]

\end{definition}

The configuration space is thereby partitioned into collections of states that cannot be distinguished by the chosen observation.

This partition is not an accidental byproduct of observation.

It is the mathematical object from which information itself will emerge.

Indeed, every measurement may now be understood as determining the observational class containing the true configuration. Information is therefore not attached to individual states but to the progressive refinement of these observational partitions.

This viewpoint unifies several apparently unrelated mathematical constructions.

In Shannon's theory, information reduces uncertainty among distinguishable messages.

In statistical mechanics, macrostates partition microscopic configurations.

In model theory, logical theories partition structures into elementary equivalence classes.

In programming language semantics, observational equivalence identifies programs that cannot be distinguished by any permitted execution context.

Each discipline studies different mathematical objects.

Each nevertheless constructs partitions induced by observation.

The present theory recognizes this common structure as fundamental.

The significance of these partitions extends beyond information theory.

Suppose that a transformation

\[
\tau:\Omega\rightarrow\Omega
\]

acts upon the configuration space.

If the transformation maps distinct observational classes into a single class, then information has been irreversibly lost relative to the observation map.

Conversely, if the transformation preserves the partition structure, then distinguishability survives the evolution.

This observation provides the first rigorous interpretation of distinction preservation.

The invariant is no longer merely a subset of the state space.

Rather, it is the preservation of the observational partition induced by the relevant family of observation maps.

The conceptual shift is substantial.

Persistence does not require that configurations remain unchanged.

Nor does it require that every numerical quantity remain constant.

Instead, persistence requires that those observational distinctions necessary for identifying the evolving system remain recoverable throughout its history.

The theory therefore acquires an explicitly relational character.

Configurations possess meaning only through the distinctions preserved by observation.

Information becomes the refinement of observational partitions.

Invariant structures become those properties that prevent the collapse of these partitions under admissible evolution.

Admissibility itself will ultimately be defined by the preservation of an appropriate family of observational structures rather than by arbitrary collections of externally imposed constraints.

The next section strengthens this perspective by introducing families of observation maps and demonstrating how invariant structures arise naturally as quantities preserved across every member of such a family.

\section{Families of Observation and the Emergence of Invariants}

The introduction of observation maps provides a precise mathematical meaning for distinguishability. A single observation map partitions the configuration space into equivalence classes whose members cannot be distinguished by the corresponding observational process. Persistence may therefore be understood as the preservation of those partitions whose continued existence is necessary for identifying the evolving system.

A single observation, however, is rarely sufficient to characterize a persistent object.

Consider a biological organism. Genetic measurements, metabolic activity, anatomical structure, behavioral responses, and ecological interactions each reveal different aspects of the same organism. None alone completely characterizes its persistence. A genetic sequence may remain unchanged while metabolism ceases. Metabolic activity may continue while ecological relationships collapse. Anatomical structure may remain intact after death despite the disappearance of every dynamic process associated with life.

The persistence of the organism therefore cannot be reduced to the preservation of a single observational partition.

A similar phenomenon appears in computation.

A compiler distinguishes programs according to syntactic structure. An operating system distinguishes executing processes according to memory isolation. A user distinguishes software through externally observable behavior. A formal verification system distinguishes implementations according to semantic correctness. Each observation captures a different structural aspect of the computational system.

Again, no individual observation is sufficient.

Persistence depends upon the simultaneous preservation of multiple distinguishable structures.

This suggests replacing the single observation map introduced previously with a collection of mutually complementary observations.

\begin{definition}[Observation Family]

Let

\[
\Omega
\]

be a configuration space.

An \emph{observation family} is a collection

\[
\mathcal{O}
=
\{O_i\}_{i\in I},
\]

where every

\[
O_i:\Omega\rightarrow\mathcal{Y}_i
\]

is an observation map.

The index set

\[
I
\]

may be finite or infinite.

\end{definition}

The observation family represents the total collection of perspectives through which persistence is to be evaluated.

Different choices of

\[
\mathcal{O}
\]

produce different mathematical theories.

A physical theory selects observations corresponding to measurable physical quantities.

A logical theory selects observations preserving formal consequence.

A computational theory selects observations describing executable semantics.

A biological theory selects observations associated with viability and regulation.

An institutional theory selects observations preserving constitutional legitimacy.

The general framework deliberately refrains from privileging any particular family.

Instead, the mathematics concerns structural properties common to every observation family.

The first such property is agreement.

Suppose that two configurations

\[
x,y\in\Omega
\]

produce identical observations for every member of the family.

Then

\[
O_i(x)=O_i(y)
\qquad
\text{for every}
\quad
i\in I.
\]

From the standpoint of the observation family, the two configurations become indistinguishable.

This motivates the following definition.

\begin{definition}[Family Equivalence]

Let

\[
\mathcal{O}
=
\{O_i\}_{i\in I}
\]

be an observation family.

Two configurations satisfy

\[
x\sim_{\mathcal O}y
\]

whenever

\[
O_i(x)=O_i(y)
\]

for every

\[
i\in I.
\]

\end{definition}

Unlike observational equivalence relative to a single map, family equivalence reflects simultaneous agreement across multiple observational perspectives.

As the observation family expands, the equivalence classes generally become smaller.

Each additional observation refines the partition of the configuration space.

Consequently, the observation family determines the resolution at which persistence is studied.

This refinement process admits an intuitive interpretation.

Every new observational capability eliminates certain ambiguities while preserving others.

Scientific instruments increase observational resolution.

Mathematical axioms distinguish additional structures.

Software verification reveals implementation defects invisible to ordinary execution.

Evolution itself may be viewed as progressively refining observational distinctions through environmental selection.

In each case, increasing observational richness corresponds mathematically to refining the induced partition of the configuration space.

Persistence therefore becomes increasingly demanding.

A transformation preserving coarse observational distinctions may nevertheless destroy distinctions visible under more refined observation.

The theory must therefore identify transformations preserving every observational partition judged essential.

This requirement naturally gives rise to invariant structure.

\begin{definition}[Invariant Relative to an Observation Family]

Let

\[
\mathcal{O}
=
\{O_i\}_{i\in I}
\]

be an observation family.

A property

\[
\Phi
\]

is called an \emph{invariant relative to}
\[
\mathcal O
\]

if every admissible transformation preserves the observational distinctions represented by

\[
\mathcal O
\]

through preservation of

\[
\Phi.
\]

\end{definition}

The wording of this definition deserves careful attention.

The invariant is not defined independently of observation.

Nor is observation determined by the invariant.

Instead, the two concepts are mutually constrained.

Observation determines which distinctions matter.

Invariant structure determines how those distinctions survive transformation.

The relationship resembles that between geometry and symmetry.

A geometry specifies the structures under consideration.

Symmetries preserve those structures.

Here, the observation family specifies distinguishable organization.

Invariant structure preserves that organization.

This perspective resolves an ambiguity that frequently appears across scientific disciplines.

Conservation laws are often presented as immutable numerical equalities.

Within the present framework, numerical conservation constitutes only one possible realization of invariant structure.

An invariant may equally consist of preservation of logical validity, maintenance of executable semantics, retention of biological viability, or continuity of institutional authority.

Their mathematical role is identical.

Each prevents observational distinctions from collapsing under admissible evolution.

This observation finally explains why invariant families appear so naturally across otherwise unrelated disciplines.

Their common origin does not lie in shared physical mechanisms.

It lies in the universal requirement that persistent systems remain distinguishable throughout continued transformation.

The mathematics therefore shifts from preserving quantities to preserving recoverable organization.

Numerical conservation becomes one important special case.

The general theory concerns preservation of observational structure itself.

Having established the relationship between observation families and invariant families, we are now prepared to introduce one of the central constructions of the entire monograph.

Given an observation family

\[
\mathcal O,
\]

how should multiple invariant requirements be combined into a single mathematical criterion governing admissible continuation?

The answer to this question leads naturally to the aggregation operators developed in the following section and ultimately to the formal definition of admissibility itself.

\section{Aggregation of Invariant Families}

The preceding sections established that persistence is evaluated relative to an observation family rather than a single observation map. Every observation induces a partition of the configuration space, and every invariant preserves some aspect of the distinguishable structure represented by that partition. The natural consequence is that a persistent system is rarely governed by a single invariant. Instead, persistence emerges from the simultaneous preservation of many structurally distinct requirements.

This observation raises an immediate mathematical problem.

Suppose that a system possesses several invariant families,

\[
\Phi_1,\Phi_2,\ldots,\Phi_n.
\]

Each invariant describes a different aspect of persistence. One may quantify structural integrity, another informational recoverability, another logical consistency, another energetic feasibility, and another computational correctness. Although these invariants may arise from different observational perspectives, admissible evolution requires that they be considered together.

How should such heterogeneous information be combined?

At first glance the answer appears obvious. One simply declares that every invariant must be preserved independently. While logically correct, this formulation proves mathematically inconvenient. It obscures interactions among invariants, complicates optimization, and provides no language for describing gradual degradation, repair, or trade-offs between competing structural requirements.

A more useful approach is to introduce a general aggregation operator.

The purpose of this operator is not to replace individual invariants. Rather, it provides a mathematical mechanism through which multiple invariant families may be evaluated collectively while preserving their individual identities.

We therefore introduce the following definition.

\begin{definition}[Invariant Aggregation Operator]

Let

\[
\Phi_1,\Phi_2,\ldots,\Phi_n
\]

be invariant functionals defined on a common transition space.

An \emph{aggregation operator} is a mapping

\[
F:
\mathbb{R}^{\,n}
\longrightarrow
\mathbb{R}
\]

which assigns to every collection of invariant evaluations

\[
(\Phi_1,\ldots,\Phi_n)
\]

a single aggregate value

\[
F(\Phi_1,\ldots,\Phi_n).
\]

\end{definition}

The definition intentionally avoids prescribing a particular functional form.

This generality is essential.

Different persistent systems require different aggregation mechanisms.

Some systems fail catastrophically whenever any single invariant fails. Others admit graceful degradation. Still others permit temporary violations provided suitable repair remains possible.

No universal aggregation rule can capture every such phenomenon.

Instead, the aggregation operator itself becomes part of the mathematical description of the persistence problem.

Several familiar constructions emerge immediately as special cases.

If persistence requires simultaneous satisfaction of every invariant without compensation, one may define

\[
F(\Phi_1,\ldots,\Phi_n)
=
\min_i \Phi_i.
\]

The aggregate is therefore controlled entirely by the weakest preserved invariant.

Such an operator naturally describes safety-critical systems in which failure of any essential component destroys persistence.

Alternatively, one may consider multiplicative aggregation,

\[
F(\Phi_1,\ldots,\Phi_n)
=
\prod_{i=1}^{n}\Phi_i,
\]

provided the invariant evaluations are suitably normalized.

This construction penalizes deterioration across multiple dimensions simultaneously and has the desirable property that complete failure of any indispensable invariant collapses the aggregate to zero.

Weighted additive constructions,

\[
F(\Phi_1,\ldots,\Phi_n)
=
\sum_{i=1}^{n}
w_i\Phi_i,
\qquad
w_i\ge0,
\]

emphasize relative importance among invariant families while permitting compensatory behavior between partially degraded components.

More sophisticated operators may involve nonlinear optimization, lattice operations, probabilistic combination rules, or category-theoretic limits.

The abstract theory remains compatible with each of these possibilities.

The essential observation is not the algebraic form of \(F\).

It is the existence of a coherent mathematical procedure by which multiple invariant requirements may be evaluated jointly.

This distinction is more important than it may first appear.

Historically, many engineering systems have accumulated independent safety rules, verification procedures, performance objectives, and regulatory constraints without any overarching mathematical principle describing their interaction. As the number of independent requirements increases, reasoning about the resulting system becomes progressively more difficult.

The aggregation operator provides precisely such a principle.

Rather than treating invariant families as unrelated collections of constraints, it interprets them as components of a single persistence structure.

The mathematical consequences become particularly significant when systems evolve.

Suppose that a transformation improves one invariant while degrading another.

Without aggregation there exists no principled method for determining whether the overall transition should be regarded as preserving persistence.

Different observers may emphasize different invariant families, producing incompatible evaluations.

The aggregation operator resolves this ambiguity by specifying how the observation family itself induces a unified persistence criterion.

It therefore serves as the bridge connecting observation to admissibility.

Several structural properties of aggregation operators prove especially useful throughout the remainder of this monograph.

\begin{definition}[Monotone Aggregation]

An aggregation operator

\[
F
\]

is said to be \emph{monotone} whenever

\[
\Phi_i\le\Psi_i
\quad
\text{for every}
\quad
i
\]

implies

\[
F(\Phi_1,\ldots,\Phi_n)
\le
F(\Psi_1,\ldots,\Psi_n).
\]

\end{definition}

Monotonicity formalizes an intuitively desirable property.

Improving every invariant should never decrease the aggregate evaluation.

Although obvious in many applications, this property excludes pathological aggregation schemes in which strengthening one component paradoxically reduces overall persistence.

A second useful property concerns continuity.

Persistent systems rarely evolve through discontinuous jumps in admissibility under infinitesimal perturbations.

Small changes in invariant evaluations should ordinarily produce correspondingly small changes in aggregate persistence.

Accordingly, continuity assumptions imposed upon \(F\) permit later development of geometric repair trajectories, stability analysis, and variational formulations.

Additional algebraic assumptions—such as convexity, associativity, or idempotence—may be imposed whenever dictated by particular applications, but none is required for the abstract development presented here.

The aggregation operator therefore occupies a deliberately intermediate position within the hierarchy of the theory.

It is more primitive than admissibility because admissibility cannot be evaluated until multiple invariant families have been combined.

It is less primitive than observation because invariant families themselves arise from observational distinctions.

The resulting logical progression may now be stated clearly.

Observation determines distinguishability.

Distinguishability determines invariant families.

Invariant families are combined through aggregation.

Aggregation produces a unified persistence functional.

Admissibility will be defined relative to that functional.

The following section completes this construction by introducing admissible transitions themselves. There, for the first time, persistence will be expressed as a mathematically well-defined property of transformations rather than merely as an intuitive principle governing the evolution of distinguishable systems.

\begin{remark}[Comparison with the Yarncrawler Endofunctor]

Readers familiar with the Yarncrawler framework may recognize a close structural analogy between the aggregation operator introduced above and the action of an endofunctor on a category of evolving structures.

The analogy should not be interpreted as an identification.

The aggregation operator
\[
F
\]
acts upon evaluations of invariant families. It combines multiple preservation criteria into a single persistence functional without itself specifying how systems evolve. An endofunctor, by contrast, transforms objects and morphisms while preserving the categorical structure within which those objects reside.

Nevertheless, both constructions express the same underlying mathematical philosophy: evolution should preserve admissible structure rather than merely generate new configurations.

Within the Yarncrawler setting, one may regard an endofunctor
\[
\mathcal{T}:\mathcal{C}\rightarrow\mathcal{C}
\]
as describing an admissible refinement of the category of persistent systems. Objects are transformed into more informative or more developed objects, while morphisms are transformed into corresponding continuation morphisms. Composition and identities are preserved,

\[
\mathcal{T}(g\circ f)
=
\mathcal{T}(g)\circ\mathcal{T}(f),
\qquad
\mathcal{T}(\mathrm{id}_X)
=
\mathrm{id}_{\mathcal{T}(X)},
\]

ensuring that the compositional structure of admissible continuations survives refinement.

The aggregation operator introduced in the present chapter may therefore be interpreted as assigning a scalar evaluation to the persistence properties that such an endofunctor preserves. In this interpretation, the endofunctor specifies \emph{how} admissible structures evolve, whereas the aggregation functional evaluates \emph{whether} the resulting evolution continues to satisfy the invariant family under consideration.

This distinction is important. The endofunctor is structural, while the aggregation operator is evaluative. One describes the transformation of persistent organization; the other measures the preservation of that organization.

From this perspective, the two constructions are complementary rather than competing. The categorical formulation emphasizes compositionality and recursive closure, while the aggregation framework developed in this monograph emphasizes admissibility, repair, and quantitative persistence. Together they suggest a broader picture in which persistent systems evolve through structure-preserving endofunctorial transformations whose admissibility is determined by invariant aggregation.

This relationship will be revisited in Part~III, where recursive continuation is formulated categorically and admissible repair operators are shown to define natural transformations between persistence-preserving endofunctors.
\end{remark}

\section{Admissible Transitions}

The preceding sections have established a sequence of progressively richer mathematical constructions. We began with primitive distinctions, introduced observation maps to formalize distinguishability, generalized these maps into observation families, derived invariant structures as the preservation of observational distinctions, and finally introduced aggregation operators capable of combining multiple invariant families into a single persistence functional.

The purpose of these constructions has been preparatory.

We now possess sufficient machinery to define the central concept of this monograph.

The object of primary interest is not a state.

Nor is it an action.

Neither is it an optimization objective.

Rather, it is the transition relating one distinguishable configuration to another.

The remainder of this work will therefore study systems through the admissibility of their transitions rather than through properties of isolated configurations.

This shift is more than terminological. It reflects a fundamental change in mathematical perspective.

A state possesses no intrinsic temporal direction. It is simply a point in a configuration space. Questions concerning persistence cannot be answered by examining individual points in isolation, because persistence concerns relationships extending across successive configurations.

Likewise, an action possesses no intrinsic guarantee that its consequences preserve the structures necessary for continued evolution. Two actions that appear identical when considered locally may differ dramatically in their long-term effects upon distinguishability, recoverability, and invariant preservation.

Transitions occupy the correct level of abstraction because they simultaneously describe change and preserve the relationship between successive configurations.

We therefore formalize the notion of admissibility.

\begin{definition}[Admissible Transition]

Let

\[
\Omega
\]

be a configuration space,

\[
\mathcal{O}
=
\{O_i\}_{i\in I}
\]

an observation family,

\[
\Phi
=
(\Phi_1,\ldots,\Phi_n)
\]

an associated invariant family, and

\[
F
\]

an aggregation operator.

A transition

\[
\tau:x\longrightarrow y
\]

is called \emph{admissible} whenever the aggregated invariant evaluation satisfies

\[
F\!\left(
\Phi_1(\tau),
\ldots,
\Phi_n(\tau)
\right)
\geq
\theta,
\]

where

\[
\theta
\]

is the admissibility threshold determined by the persistence problem under consideration.

The collection of all admissible transitions will be denoted

\[
\mathcal{A}.
\]

\end{definition}

Several features of this definition deserve careful discussion.

First, admissibility is fundamentally relational.

Neither the initial configuration

\[
x
\]

nor the terminal configuration

\[
y
\]

is classified independently.

Only the transition connecting them receives an admissibility value.

This immediately distinguishes the present framework from approaches that classify individual states as safe, unsafe, valid, or invalid.

Such classifications frequently obscure the fact that an apparently safe configuration may only have been reached through catastrophic loss of recoverability, while an apparently dangerous configuration may lie along an entirely admissible repair trajectory.

Second, admissibility depends upon the observation family.

Changing the observation family changes the invariant family.

Changing the invariant family changes the aggregation functional.

Consequently, admissibility is never an intrinsic property of the underlying configuration space alone.

It is always defined relative to the structures that are intended to persist.

This relativity should not be regarded as a weakness.

On the contrary, it explains why identical physical transformations may be judged differently by different scientific disciplines.

A compiler evaluates semantic correctness.

A physicist evaluates conservation laws.

A biologist evaluates viability.

A constitutional court evaluates procedural legitimacy.

Each employs a different observation family and therefore derives a different admissibility relation.

The underlying mathematical framework remains unchanged.

Third, the admissibility threshold

\[
\theta
\]

need not represent a universal constant.

Different systems tolerate different degrees of degradation.

Highly redundant biological organisms often admit substantial local damage while remaining persistent.

By contrast, cryptographic verification algorithms may fail completely after a single incorrect bit.

The threshold therefore belongs to the mathematical specification of the persistence problem rather than to the abstract theory itself.

The admissible transition relation

\[
\mathcal{A}
\]

induces a directed graph upon the configuration space.

Vertices correspond to configurations.

Directed edges correspond to admissible transitions.

Persistence then becomes a question concerning paths within this graph.

Configurations themselves possess no privileged status.

The mathematical structure resides in the network of admissible continuations connecting them.

This viewpoint immediately suggests an important reinterpretation of evolution.

Traditional dynamical systems begin with an evolution law and investigate the trajectories it generates.

The present theory instead begins with the admissibility relation.

Dynamics become one possible mechanism for selecting admissible transitions, but they no longer define persistence by themselves.

Persistence is determined by the existence of admissible paths.

The particular dynamics merely determine which of those paths are realized.

This distinction will become increasingly important in later chapters dealing with repair, recursion, institutional evolution, and adaptive computation.

In each of these domains, multiple admissible continuations typically exist.

The central mathematical problem is not merely predicting which continuation occurs, but characterizing the structure of the admissible continuation space itself.

The remainder of this chapter investigates the properties of the admissible transition graph. In particular, we shall demonstrate that persistence may be characterized as the existence of arbitrarily extendable admissible paths, thereby replacing static notions of equilibrium with a fundamentally relational theory of continued existence.

\section{The Extension Principle}

The introduction of admissible transitions permits a fundamental reformulation of persistence. Up to this point, persistence has been described informally as the capacity of a distinguishable system to continue evolving without destroying the invariant structures upon which its identity depends. The admissibility relation now allows this intuition to be expressed in entirely graph-theoretic terms.

Rather than asking whether a particular configuration persists, we ask whether it admits continued extension through admissible transitions.

This change in viewpoint is subtle but decisive.

Persistence is no longer identified with occupying a privileged region of configuration space. Nor is it characterized by proximity to equilibrium, optimization of an objective function, or convergence toward a fixed point. Instead, persistence becomes a property of the local and global geometry induced by the admissibility relation.

To make this precise, recall that the admissible transition relation

\[
\mathcal{A}\subseteq\Omega\times\Omega
\]

defines a directed graph whose vertices are configurations and whose directed edges are admissible transitions.

Within this graph, every evolution of the system corresponds to a directed path.

The mathematical question therefore becomes one of extendability.

Can every admissible history be continued?

This motivates the following definition.

\begin{definition}[Admissible Path]

An \emph{admissible path} is a finite or infinite sequence

\[
(x_0,x_1,x_2,\ldots)
\]

such that

\[
(x_k,x_{k+1})\in\mathcal A
\]

for every index \(k\) for which the successor exists.

When the sequence terminates after \(m\) transitions, it is called a finite admissible path of length \(m\).

\end{definition}

Notice that no assumptions are made regarding determinism.

A given configuration may possess no admissible successors, exactly one admissible successor, or infinitely many admissible successors.

The admissibility graph therefore accommodates deterministic dynamics, nondeterministic computation, stochastic evolution, branching developmental processes, and adaptive control within a single mathematical framework.

The distinction between finite and infinite admissible paths is central.

A finite path ending at a configuration having no admissible successor represents the exhaustion of persistence. Such a terminal configuration need not be physically impossible, logically inconsistent, or computationally undefined. It merely possesses no continuation satisfying the invariant structures prescribed by the observation family.

Conversely, an infinite admissible path represents indefinite continuation. The system may change arbitrarily throughout its evolution, provided every transition remains admissible.

Persistence therefore concerns extendability rather than immobility.

The following definition formalizes this observation.

\begin{definition}[Extendable Configuration]

A configuration

\[
x\in\Omega
\]

is called \emph{extendable} if every finite admissible path terminating at \(x\) admits at least one admissible extension.

Equivalently, whenever

\[
(x_0,\ldots,x_{m-1},x)
\]

is an admissible path, there exists some configuration

\[
y\in\Omega
\]

such that

\[
(x,y)\in\mathcal A.
\]

\end{definition}

The definition deliberately quantifies over histories rather than isolated vertices.

Persistence is not merely the existence of one outgoing edge from one configuration viewed in isolation. Instead, it requires that every admissible realization terminating at that configuration remains capable of further admissible evolution.

This historical formulation reflects one of the guiding principles of the present theory.

The admissibility of future evolution depends not only upon the present configuration but upon the recoverability of the structures accumulated throughout its history.

In systems possessing path dependence, hysteresis, memory, or irreversible transitions, identical configurations reached through different admissible histories may admit different future continuations. Later chapters will generalize the formalism to accommodate such history-dependent admissibility relations explicitly. For the present development, however, the simplified state-based relation suffices to establish the central theorem.

\begin{theorem}[Extension Principle]

Suppose that every configuration encountered along an admissible path is extendable.

Then the path admits an arbitrarily long admissible continuation.

If, in addition, every extendable configuration possesses at least one extendable admissible successor, then every finite admissible path extends to an infinite admissible path.

\end{theorem}

\begin{proof}

Let

\[
(x_0,\ldots,x_m)
\]

be an arbitrary finite admissible path.

Since

\[
x_m
\]

is extendable, there exists

\[
x_{m+1}
\]

such that

\[
(x_m,x_{m+1})\in\mathcal A.
\]

Hence

\[
(x_0,\ldots,x_m,x_{m+1})
\]

is itself an admissible path.

If

\[
x_{m+1}
\]

is likewise extendable, the argument may be repeated.

Inductively, one constructs admissible paths of every finite length,

\[
(x_0,\ldots,x_{m+r})
\]

for arbitrary positive integers \(r\).

This proves the existence of arbitrarily long admissible continuations.

Now assume additionally that every extendable configuration admits at least one extendable admissible successor.

Beginning with the original path, choose such a successor at each stage.

The resulting inductive construction never terminates because the inductive hypothesis is preserved after every extension.

Consequently, the finite admissible path generates an infinite sequence

\[
(x_0,x_1,x_2,\ldots)
\]

whose successive transitions all belong to

\[
\mathcal A.
\]

Thus every finite admissible history extends to an infinite admissible continuation.

\end{proof}

The theorem is intentionally elementary.

Its importance lies not in technical sophistication but in conceptual orientation.

Persistence has now been characterized without reference to equilibrium, optimization, conserved quantities, or fixed points.

The only essential ingredient is the continual existence of admissible extension.

This observation represents one of the principal departures of the present framework from many classical theories of dynamics.

Traditional stability theory asks whether trajectories remain near distinguished configurations.

The present theory asks whether trajectories remain capable of further admissible evolution.

These questions coincide in certain special cases, particularly when equilibria themselves possess arbitrarily extendable neighborhoods. In general, however, they differ substantially.

A perfectly stable dead system is not persistent.

An evolving adaptive system may never approach equilibrium and yet remain indefinitely persistent.

The mathematical distinction between stability and continuation will become increasingly significant in later chapters. Repair operators, recursive modification, evolutionary adaptation, institutional development, and scientific progress all derive their persistent character not from remaining unchanged, but from preserving the possibility of further admissible change.

Continuation therefore replaces equilibrium as the primary geometric object of the theory.

\section{The Geometry of Continuation}

The Extension Principle establishes persistence as a property of extendability rather than equilibrium. This result invites a corresponding geometric interpretation. Instead of viewing the configuration space as a collection of isolated states through which trajectories happen to pass, we may regard it as a landscape whose geometry is determined by admissible continuation.

This geometry differs fundamentally from the geometry induced by a metric.

A metric measures proximity.

The geometry of continuation measures possibility.

Two configurations may lie arbitrarily close in Euclidean distance while possessing entirely different admissible futures. Conversely, configurations separated by large metric distances may nevertheless belong to the same persistent continuation if every intermediate transition preserves the relevant invariant family.

The geometry developed in this chapter therefore concerns reachability under admissibility rather than distance alone.

To formalize this perspective, let

\[
\mathcal{A}\subseteq\Omega\times\Omega
\]

denote the admissible transition relation introduced previously.

Associated with every configuration

\[
x\in\Omega
\]

is the collection of configurations that may be reached through admissible evolution.

\begin{definition}[Admissible Reachability]

The \emph{admissible reachability set} of a configuration

\[
x\in\Omega
\]

is defined by

\[
\mathcal{R}(x)
=
\left\{
y\in\Omega
\;\middle|\;
\text{there exists an admissible path from }x\text{ to }y
\right\}.
\]

\end{definition}

Unlike ordinary reachability, this construction depends explicitly upon the invariant family through the admissibility relation.

Changing the observation family changes the invariant family.

Changing the invariant family changes the admissible transition relation.

Consequently, the reachability geometry itself changes.

Persistence is therefore relative to the structures one seeks to preserve.

This dependence is neither arbitrary nor undesirable.

In classical mechanics, accessibility depends upon conservation laws.

In computation, accessibility depends upon semantic correctness.

In biology, accessibility depends upon physiological viability.

In law, accessibility depends upon procedural legitimacy.

The same mathematical construction accommodates each of these domains by modifying only the underlying invariant family.

The reachability sets therefore constitute one of the first genuinely universal objects of the theory.

Their boundaries deserve particular attention.

Suppose that

\[
y\in\mathcal{R}(x).
\]

If every sufficiently small admissible perturbation of

\[
y
\]

also belongs to

\[
\mathcal{R}(x),
\]

then the system occupies an interior region of the continuation geometry.

Minor disturbances do not threaten persistence.

On the other hand, configurations lying near the edge of the reachability set behave differently.

Small perturbations may eliminate every admissible continuation.

Such configurations inhabit what may be called the continuation boundary.

Although no topological assumptions have yet been imposed, the intuition is already clear.

Persistent systems spend much of their existence moving through interior regions where numerous admissible continuations remain available.

Repair becomes necessary only as trajectories approach regions in which admissible continuation begins to disappear.

This observation provides an important reinterpretation of robustness.

Traditionally, robustness is viewed as resistance to perturbation.

Within the present framework, robustness becomes the possession of a sufficiently large neighborhood of admissible continuations.

The distinction is subtle.

Resistance emphasizes what the system prevents.

Continuation emphasizes what the system preserves.

The former is defensive.

The latter is constructive.

This constructive interpretation naturally accommodates adaptive systems whose persistence depends not upon resisting every disturbance but upon continually discovering new admissible trajectories.

Evolutionary biology provides an instructive example.

Species rarely persist by remaining genetically unchanged.

Instead, they persist because viable evolutionary paths continue to exist within changing environments.

Likewise, scientific disciplines survive not because individual theories remain fixed, but because admissible revisions preserve standards of evidence while extending explanatory power.

The same principle governs software maintenance.

Programs evolve through successive modifications.

Persistence depends not upon preserving every line of source code, but upon maintaining a sufficiently rich collection of admissible future revisions.

The common mathematical feature across these examples is the preservation of reachable continuation rather than preservation of static structure.

This observation motivates the introduction of one further concept.

\begin{definition}[Continuation Capacity]

The \emph{continuation capacity} of a configuration

\[
x\in\Omega
\]

is the structural richness of its admissible reachability set

\[
\mathcal{R}(x).
\]

Any quantitative functional assigning greater values to configurations possessing more diverse admissible futures may be regarded as a measure of continuation capacity.

\end{definition}

The present definition intentionally leaves the notion of "structural richness" abstract.

Different applications require different quantitative realizations.

One may measure the cardinality of reachable configurations in finite systems.

One may employ measure-theoretic volume in continuous spaces.

Graph-theoretic branching factors naturally describe discrete computational systems.

Information-theoretic entropy may quantify the diversity of admissible futures.

Topological invariants may describe the connectivity of continuation regions.

The abstract theory remains compatible with each of these realizations.

The important point is not the numerical value assigned by any particular capacity functional.

Rather, it is the recognition that persistence admits degrees.

Some configurations possess only a handful of admissible continuations.

Others possess extraordinarily rich continuation geometries supporting extensive adaptation, repair, exploration, and recursive self-modification.

Continuation capacity therefore provides a mathematical language for describing resilience without reducing persistence to binary survival.

Indeed, many phenomena traditionally treated separately—robustness, adaptability, evolvability, flexibility, optionality, and resilience—may ultimately be interpreted as different manifestations of continuation capacity under different observation families.

This unification will become increasingly important in subsequent chapters.

Repair operators will be shown to increase continuation capacity by restoring admissible transitions previously lost through degradation.

Recursive continuation will enlarge continuation capacity through self-modification that preserves invariant structure.

Institutional and scientific development will likewise be interpreted as the long-term expansion of admissible continuation rather than the accumulation of static knowledge.

The geometry introduced here therefore serves as the bridge between the local concept of admissible transition and the global theory of persistent evolution developed throughout the remainder of this monograph.

\section{Repair as the Restoration of Continuation}

The geometric interpretation developed in the previous section suggests that persistence should not be viewed as a binary property possessed or absent in isolation. Every configuration occupies a position within a landscape of admissible continuations. Some regions admit a rich collection of future trajectories, while others lie precariously close to boundaries beyond which no admissible continuation exists. Between these extremes lies the phenomenon that motivates much of the present theory: repair.

Repair is among the most ubiquitous processes observed in persistent systems. Biological organisms repair damaged tissue. Software systems repair corrupted data structures. Communication protocols repair transmission errors. Scientific communities repair erroneous theories through experiment and revision. Legal institutions repair procedural defects through appeals and constitutional review. Even physical systems exhibit repair-like behavior through relaxation toward dynamically stable configurations after perturbation.

Despite its prevalence, repair is frequently treated as a domain-specific mechanism. Biology studies wound healing, computer science studies fault tolerance, engineering studies maintenance, and jurisprudence studies judicial review. Each discipline develops highly specialized techniques adapted to its own subject matter.

The present theory seeks something more general.

Rather than viewing repair as an application-specific procedure, we regard it as a structural operation acting upon the geometry of admissible continuation.

The essential purpose of repair is not merely to restore a previous state.

Indeed, many successful repairs never reconstruct the exact configuration that existed before the disturbance. A healed bone contains remodeled tissue rather than the original microscopic structure. A corrected software system incorporates new source code rather than recovering the original defective implementation. A revised scientific theory typically differs substantially from its predecessor while preserving the observations that motivated its development.

These examples illustrate an important distinction.

Repair restores the possibility of continued admissible evolution.

It need not restore the past.

This observation motivates the following abstract definition.

\begin{definition}[Repair Operator]

Let

\[
\mathcal{A}
\]

denote the admissible transition relation on the configuration space

\[
\Omega.
\]

A \emph{repair operator} is a mapping

\[
R:\Omega\longrightarrow\Omega
\]

such that whenever

\[
x
\]

possesses reduced continuation capacity relative to the invariant family under consideration, the repaired configuration

\[
R(x)
\]

admits an admissible reachability structure satisfying

\[
\mathcal{R}(R(x))
\supseteq
\mathcal{R}(x),
\]

where the inclusion is interpreted relative to the admissibility relation induced by the chosen observation family.

\end{definition}

Several remarks clarify the intent of this definition.

First, the inclusion need not be strict.

Some repairs merely preserve existing admissible futures after local degradation has been arrested.

Others restore admissible continuations that had previously disappeared.

Still others generate entirely new regions of admissible evolution unavailable before the disturbance.

The definition encompasses all three possibilities.

Second, repair is defined relative to continuation rather than state.

Traditional error-correction procedures frequently specify a target configuration that the repaired system must reproduce exactly. Such formulations are appropriate whenever exact reconstruction is both possible and desirable. Persistent systems, however, often possess many admissible futures compatible with the preservation of their defining invariant structures.

Consequently, repair should not be understood as inversion.

In general,

\[
R
\]

is neither the inverse of the perturbation nor a projection onto an ideal configuration.

Its mathematical purpose is to restore admissible continuation.

This distinction becomes especially significant in adaptive systems.

Suppose that a perturbation destroys one admissible trajectory while simultaneously revealing another previously unavailable route through the continuation geometry.

A repair operator that insists upon exact reconstruction may unnecessarily eliminate these newly available possibilities.

By contrast, a repair operator guided by continuation preserves the enlarged admissible structure whenever doing so remains compatible with the invariant family.

Repair therefore need not oppose adaptation.

Under appropriate circumstances it may actively facilitate it.

The geometry introduced in the previous section makes this relationship transparent.

Perturbations alter the continuation landscape.

Repair reshapes that landscape so that sufficiently rich admissible futures remain available.

The resulting evolution may differ substantially from the original trajectory while nevertheless preserving persistence.

This interpretation also clarifies the distinction between maintenance and repair.

Maintenance operates primarily within the interior of regions possessing high continuation capacity.

Its purpose is preventive.

By preserving invariant structure before substantial degradation occurs, maintenance keeps the evolving system away from the boundaries at which admissible continuation becomes scarce.

Repair, by contrast, operates after degradation has already reduced continuation capacity.

Its purpose is restorative.

It seeks to enlarge the admissible continuation structure that has been diminished through perturbation, error, damage, or environmental change.

Although maintenance and repair differ operationally, they share the same mathematical objective.

Both preserve the possibility of future admissible evolution.

Their difference lies only in the position occupied within the continuation geometry at which intervention occurs.

The repair operator introduced here remains intentionally abstract.

Subsequent chapters will specialize this construction to a variety of settings. In discrete computational systems, repair will become graph transformations restoring semantic correctness. In continuous systems, repair will appear as trajectories ascending continuation gradients. Within recursive continuation, repair operators will themselves become objects subject to admissible modification, giving rise to hierarchies of self-repair and meta-repair.

For the present development, however, one conclusion already emerges clearly.

Repair is not fundamentally the correction of past error.

It is the restoration of future possibility.

The mathematics of persistence therefore shifts attention away from recovering what has been lost and toward preserving what may yet occur. This orientation will remain central throughout the remainder of the theory, culminating in the recursive continuation framework developed in the later parts of this monograph.

\section{Quantitative Invariant Functionals}

The preceding sections established that invariant structure arises from the preservation of observational distinctions rather than from arbitrary collections of externally imposed constraints. This qualitative perspective explains the origin of invariants but does not yet provide a mathematical language for comparing different admissible transitions. In practical applications one rarely encounters systems in which every invariant is either perfectly preserved or completely destroyed. Biological systems accumulate damage gradually, computational systems exhibit partial degradation, scientific theories explain some observations more successfully than others, and engineered systems often continue operating despite localized failures.

The mathematics of persistence must therefore accommodate degrees of preservation.

The objective of the present section is to replace the purely qualitative notion of an invariant by a quantitative functional measuring the extent to which a transition preserves the observational structures introduced earlier in this chapter.

Let

\[
\mathcal{O}
=
\{O_i\}_{i\in I}
\]

be an observation family on a configuration space

\[
\Omega.
\]

Each observation map induces an observational partition of the configuration space. An admissible transition should preserve these partitions as completely as possible. The simplest mathematical representation of this requirement is a real-valued functional assigned to every transition.

\begin{definition}[Invariant Functional]

An \emph{invariant functional} is a mapping

\[
\Phi:\mathcal{A}\rightarrow[0,\infty),
\]

where

\[
\mathcal A
\]

denotes the space of candidate transitions.

The value

\[
\Phi(\tau)
\]

measures the degree to which the transition

\[
\tau
\]

preserves the observational distinctions associated with a specified observation family.

\end{definition}

The codomain has been chosen to be non-negative for convenience rather than necessity. Different applications may normalize invariant functionals differently. Some may assign the value one to perfect preservation and zero to complete failure. Others may employ unbounded quantities representing repair cost, informational divergence, energetic expenditure, or logical inconsistency.

The abstract theory depends only upon the ordering induced by the functional.

Higher values correspond to greater preservation of the invariant under consideration.

This formulation immediately separates two concepts that are often conflated.

An invariant is not itself a numerical quantity.

Rather, the invariant is the structural property whose preservation is measured by the functional.

The functional provides the quantitative representation of that preservation.

The distinction is analogous to the relationship between distance and geometry. Geometry exists independently of any particular coordinate representation, while a metric assigns numerical values describing aspects of that geometry. Likewise, persistence exists independently of any particular invariant functional, while the functional quantifies one aspect of invariant preservation.

The use of functionals rather than Boolean predicates provides several mathematical advantages.

First, it induces an order on admissible transitions.

Given two transitions

\[
\tau_1,\tau_2,
\]

one may compare

\[
\Phi(\tau_1)
\quad\text{and}\quad
\Phi(\tau_2),
\]

thereby expressing that one transition preserves the invariant more effectively than another.

Second, quantitative functionals permit limits.

Suppose that

\[
(\tau_n)
\]

is a sequence of admissible transitions converging, in an appropriate sense, toward a limiting transition

\[
\tau.
\]

If

\[
\Phi
\]

possesses suitable continuity or lower semicontinuity properties, then the preservation of invariant structure behaves predictably under limiting processes. Such assumptions will become essential in the variational theory developed in later chapters.

Third, quantitative functionals permit optimization.

Many engineering systems already attempt to maximize safety, efficiency, robustness, or correctness. Within the present framework, these objectives become special cases of maximizing appropriate invariant functionals while remaining above the admissibility threshold established earlier.

Optimization therefore appears as a secondary construction.

Persistence remains primary.

The existence of quantitative invariant functionals also permits the introduction of defect measures.

Rather than describing only how much structure survives a transition, one may equally measure how much distinguishability has been lost. Throughout the remainder of this monograph we shall frequently pass between preservation functionals and their corresponding defect functionals. Although mathematically equivalent under appropriate normalization, the two viewpoints often illuminate different aspects of persistent evolution.

This duality is particularly important for repair theory.

Repair operators act naturally upon defect functionals by reducing accumulated degradation. Conversely, admissibility is most naturally expressed in terms of preservation functionals whose values remain above prescribed thresholds. The ability to move freely between these complementary descriptions will simplify many later constructions.

The introduction of invariant functionals therefore marks an important transition in the development of the theory. Distinction has now given rise to observation. Observation has generated invariant structure. Invariant structure has acquired quantitative representation. The remaining task is to understand how defects accumulate when invariants fail to be perfectly preserved.

This leads naturally to the concept of reconstruction defect, which provides the first quantitative measure of structural degradation and forms the mathematical foundation of the repair calculus developed in subsequent chapters.

\section{Reconstruction Defect}

The introduction of invariant functionals permits the preservation of observational structure to be measured quantitatively. A complementary perspective, however, often proves more natural. Rather than measuring how much structure survives a transition, one may instead measure the extent to which distinguishability has been degraded.

This complementary quantity will be called the reconstruction defect.

The terminology has been chosen deliberately. The word "reconstruction" emphasizes that persistence ultimately concerns recoverability rather than static preservation. A system need not remain unchanged throughout its evolution. Instead, it must remain sufficiently organized that its defining observational distinctions can, in principle, be reconstructed from its subsequent configurations. Whenever this becomes increasingly difficult, the reconstruction defect increases.

The reconstruction defect therefore measures not the amount of change undergone by a system, but the loss of recoverable structure induced by that change.

This distinction is essential.

Consider a software project that undergoes a substantial architectural refactoring while preserving its externally observable behavior, formal correctness, and maintainability. Such a transition may involve enormous structural change, yet the reconstruction defect remains small because the defining distinctions of the software continue to be recoverable.

Conversely, a single erroneous modification may silently corrupt a critical data structure. The observable state changes very little, yet the reconstruction defect may become large because future computations can no longer reliably recover the intended semantics.

The magnitude of change is therefore not an appropriate measure of persistence.

Recoverability is.

We formalize this observation as follows.

\begin{definition}[Reconstruction Defect]

Let

\[
\Phi
\]

be an invariant functional normalized so that larger values correspond to greater preservation.

A corresponding \emph{reconstruction defect functional} is a mapping

\[
\Delta_R:\mathcal A\rightarrow[0,\infty)
\]

whose value increases as recoverability decreases.

When

\[
\Phi
\]

is normalized to lie within

\[
[0,1],
\]

one convenient choice is

\[
\Delta_R(\tau)
=
1-\Phi(\tau).
\]

More generally,

\[
\Delta_R
\]

may be any monotone decreasing transformation of

\[
\Phi.
\]

\end{definition}

The general formulation is preferable because different applications naturally employ different normalizations. In some settings defect corresponds to probability of reconstruction failure. In others it represents energetic repair cost, information divergence, logical inconsistency, or accumulated semantic ambiguity. The abstract theory depends only upon the ordering induced by the functional.

The reconstruction defect should not be interpreted as a measure of error alone.

A system may possess nonzero defect while remaining entirely admissible.

Indeed, every persistent physical system experiences continual perturbation. Biological organisms accumulate molecular damage. Software accumulates technical debt. Scientific theories acquire unresolved anomalies. Institutions develop procedural imperfections.

Persistence does not require the absence of defect.

It requires that defect remain sufficiently bounded for repair to remain possible.

This observation immediately distinguishes the present framework from idealized theories in which every deviation from perfect structure constitutes failure.

Persistent systems operate within a regime of controlled imperfection.

The mathematics of persistence therefore concerns the management of defect rather than its complete elimination.

The reconstruction defect possesses several structural properties that follow directly from its interpretation.

\begin{proposition}[Minimal Defect]

Suppose

\[
\Delta_R
\]

is normalized by

\[
\Delta_R=1-\Phi,
\]

where

\[
0\le\Phi\le1.
\]

Then

\[
0\le\Delta_R\le1.
\]

Moreover,

\[
\Delta_R=0
\]

if and only if the transition preserves the corresponding invariant perfectly.

\end{proposition}

\begin{proof}

Since

\[
0\le\Phi\le1,
\]

subtracting from unity yields

\[
0
\le
1-\Phi
\le
1.
\]

Therefore

\[
0
\le
\Delta_R
\le
1.
\]

Furthermore,

\[
\Delta_R=0
\]

holds precisely when

\[
\Phi=1,
\]

which is the normalization chosen for perfect preservation.

\end{proof}

Although elementary, this proposition illustrates an important methodological principle that will recur throughout the remainder of the monograph.

Every preservation functional admits an equivalent defect representation.

The choice between the two depends upon the mathematical question under investigation.

Optimization problems often maximize preservation.

Repair problems minimize defect.

Both describe the same underlying phenomenon.

The reconstruction defect also interacts naturally with aggregation.

Suppose that

\[
\Phi_1,\ldots,\Phi_n
\]

represent multiple invariant functionals.

Each induces a corresponding defect functional

\[
\Delta_{R,1},
\ldots,
\Delta_{R,n}.
\]

Rather than aggregating preservation directly, one may instead aggregate defects,

\[
\Delta
=
G(
\Delta_{R,1},
\ldots,
\Delta_{R,n}
),
\]

for an appropriate aggregation operator

\[
G.
\]

Different choices of

\[
G
\]

lead to different notions of overall degradation.

Maximum operators describe systems dominated by their worst failure mode.

Additive operators describe cumulative degradation.

Quadratic operators penalize concentrated defects more strongly than distributed ones.

Information-theoretic operators measure statistical divergence between preserved and degraded observational structures.

The abstract framework remains compatible with each of these choices.

Only monotonicity with respect to defect is generally required.

The introduction of reconstruction defect also clarifies the mathematical role of repair.

A repair operator need not maximize every preservation functional simultaneously.

Instead, its primary objective is to reduce accumulated reconstruction defect while maintaining admissibility.

This formulation proves especially useful in recursive systems, where perfect reconstruction may be impossible or even undesirable. Self-modifying systems often improve future continuation by replacing obsolete structures with new ones rather than restoring previous configurations exactly.

Repair therefore minimizes irrecoverable loss rather than historical deviation.

This perspective aligns naturally with the philosophy developed throughout the preceding chapters.

Persistent systems do not preserve history.

They preserve the capacity to generate coherent future histories.

The reconstruction defect quantifies precisely the extent to which that capacity has been diminished.

The next step in the development of the theory is to examine how reconstruction defects accumulate along extended admissible paths. Individual transitions rarely determine persistence in isolation. Instead, long-term evolution depends upon the cumulative interaction of many small defects distributed across an entire continuation history. Understanding this accumulation process will provide the mathematical foundation for global repair strategies and the variational principles developed in the later parts of this monograph.

\chapter{Primitive Persistence Operators}

\section{The Operator-Theoretic Perspective}

The preceding chapters established three structural facts.

First, observation is necessarily non-invertible: distinct field trajectories may induce identical observable configurations under the projection operator
\[
\mathcal{O}_{\mathrm{obs}} : \mathcal{F}\rightarrow\mathcal{S}.
\]

Second, admissibility defines the subset
\[
\operatorname{Adm}_{\Phi}\subseteq\mathcal S
\]
of configurations compatible with the invariant basis
\[
\mathbf{\Phi}_{\mathrm{base}}.
\]

Third, the Continuation Principle guarantees that admissible trajectories extend locally---and, under the hypotheses of Chapter~3, globally---without requiring reconstruction of their lost field-space provenance.

The remaining question is constructive rather than existential.

Rather than asking whether admissible continuations exist, we ask how they are generated, evaluated, repaired, and committed.

This motivates the introduction of a persistence operator calculus.

The philosophy throughout this chapter is intentionally minimal.

Complex persistence mechanisms are not treated as primitive objects.

Instead, every admissible computation is assembled from a small collection of elementary persistence operators whose algebraic properties determine the behaviour of the complete system.

The resulting operator algebra plays a role analogous to elementary arithmetic operations in numerical analysis or generators of a group in abstract algebra.

\section{Persistence Operators}

Let
\[
\mathcal S
\]
denote the configuration manifold introduced in Chapter~3.

Throughout this chapter all operators are understood to be partial maps

\[
T:
\operatorname{Dom}(T)
\rightharpoonup
\mathcal S,
\]

whose domains are subsets of
\[
\mathcal S.
\]

Only operators preserving admissibility are regarded as persistence operators.

\begin{definition}[Persistence Operator]
A persistence operator is a partial transformation

\[
T:
\operatorname{Adm}_{\Phi}
\rightharpoonup
\operatorname{Adm}_{\Phi}
\]

whose image lies entirely within the admissible manifold.

The collection of all persistence operators is denoted

\[
\mathfrak T_{\operatorname{Adm}}.
\]
\end{definition}

Persistence operators therefore act exclusively within the admissible geometry established previously.

Observation itself is not a persistence operator, since it acts between different spaces,

\[
\mathcal O_{\mathrm{obs}}
:
\mathcal F
\rightarrow
\mathcal S,
\]

and serves instead as the interface between field dynamics and observable configuration dynamics.

\section{Primitive Persistence Alphabet}

The persistence algebra is generated by four primitive operators.

\begin{definition}[Primitive Alphabet]

The primitive persistence alphabet is

\[
\mathcal A_{\mathrm{prim}}
=
\{
\mathcal P,
\mathcal Q,
\mathcal C,
\operatorname{id}
\},
\]

together with the evaluation functional

\[
\mathcal E
:
\mathcal S
\rightarrow
\mathbb R_{\ge0}.
\]

The primitives have the following interpretations.

\begin{enumerate}

\item
Proposal

\[
\mathcal P
:
\mathcal S
\rightharpoonup
\mathcal S,
\]

producing candidate successor configurations.

\item
Correction

\[
\mathcal Q
:
\mathcal S
\rightharpoonup
\mathcal S,
\]

performing admissibility-preserving repair.

\item
Commitment

\[
\mathcal C
:
\mathcal S
\rightharpoonup
\mathcal S,
\]

which finalizes an admissible configuration.

\item
Identity

\[
\operatorname{id}_{\mathcal S},
\]

leaving configurations unchanged.

\item
Evaluation

\[
\mathcal E
:
\mathcal S
\rightarrow
\mathbb R_{\ge0},
\]

assigning the local reconstruction defect

\[
\delta_R(s).
\]

\end{enumerate}

\end{definition}

Unlike the remaining primitives,
\[
\mathcal E
\]
is a functional rather than an operator.
It measures admissibility but does not transform configurations.

\section{Admissibility Preservation}

Not every partial transformation belongs to
\[
\mathfrak T_{\operatorname{Adm}}.
\]

Persistence operators are distinguished by preservation of admissibility.

\begin{definition}[Admissibility Preservation]

A persistence operator

\[
T
\]

is admissibility preserving whenever

\[
T
\left(
\operatorname{Adm}_{\Phi}
\right)
\subseteq
\operatorname{Adm}_{\Phi}.
\]

\end{definition}

This property is the defining structural constraint of the persistence calculus.

Unlike arbitrary state transformations, persistence operators never intentionally leave the admissible manifold.

Violation of admissibility is handled by explicit repair operators rather than by enlarging the admissible space itself.

\begin{proposition}[Closure]

The collection

\[
\mathfrak T_{\operatorname{Adm}}
\]

is closed under composition.

\end{proposition}

\begin{proof}

Let

\[
T_1,
T_2
\in
\mathfrak T_{\operatorname{Adm}}.
\]

Then

\[
T_1
\left(
\operatorname{Adm}_{\Phi}
\right)
\subseteq
\operatorname{Adm}_{\Phi},
\]

and likewise

\[
T_2
\left(
\operatorname{Adm}_{\Phi}
\right)
\subseteq
\operatorname{Adm}_{\Phi}.
\]

Consequently,

\[
(T_2\circ T_1)
\left(
\operatorname{Adm}_{\Phi}
\right)
=
T_2
\left(
T_1
(
\operatorname{Adm}_{\Phi}
)
\right)
\subseteq
\operatorname{Adm}_{\Phi},
\]

establishing closure.

\end{proof}

\section{The Monoid of Persistence Operators}

The preceding definitions endow the collection of admissibility-preserving
transformations with a natural algebraic structure.

\begin{theorem}[Persistence Monoid]

The triple

\[
\left(
\mathfrak T_{\operatorname{Adm}},
\circ,
\operatorname{id}_{\mathcal S}
\right)
\]

forms a monoid.

\end{theorem}

\begin{proof}

Associativity follows immediately from associativity of function
composition.

The identity operator

\[
\operatorname{id}_{\mathcal S}
\]

belongs to
\(
\mathfrak T_{\operatorname{Adm}}
\)
since

\[
\operatorname{id}_{\mathcal S}
(
\operatorname{Adm}_{\Phi}
)
=
\operatorname{Adm}_{\Phi}.
\]

Closure was established in Proposition~4.1.

Therefore

\[
\left(
\mathfrak T_{\operatorname{Adm}},
\circ,
\operatorname{id}_{\mathcal S}
\right)
\]

is a monoid.

\end{proof}

\begin{remark}

The persistence monoid is generally noncommutative.

The order in which persistence operators are applied is therefore an
essential component of the computation.

In particular,

\[
T_1\circ T_2
\neq
T_2\circ T_1
\]

for generic persistence operators.

This noncommutativity reflects the directional character of continuation:
evaluation before repair is not equivalent to repair before evaluation,
nor is commitment interchangeable with proposal.

\end{remark}

\begin{proposition}[Noncommutativity]

There exist persistence operators

\[
T_1,T_2
\in
\mathfrak T_{\operatorname{Adm}}
\]

such that

\[
T_1\circ T_2
\neq
T_2\circ T_1.
\]

\end{proposition}

\begin{proof}

Let

\[
T_1=\mathcal P,
\qquad
T_2=\mathcal C.
\]

Proposal generates a candidate continuation,

whereas commitment finalizes an already admissible configuration.

Applying

\[
\mathcal C\circ\mathcal P
\]

commits the proposed continuation.

Conversely,

\[
\mathcal P\circ\mathcal C
\]

creates a new proposal after commitment,
which is generally distinct.

Therefore

\[
\mathcal C\circ\mathcal P
\neq
\mathcal P\circ\mathcal C.
\]

\end{proof}

\section{The Persistence Pipeline}

Although the persistence monoid contains many admissible compositions,
one composition occupies a distinguished position.

It represents the canonical progression by which admissible continuations
are generated, validated, repaired when necessary, and incorporated into
the evolving system.

\begin{definition}[Persistence Pipeline]

The persistence pipeline is the composite morphism

\[
\Pi
=
\mathcal C
\circ
\mathcal Q
\circ
\mathcal P,
\]

together with the auxiliary evaluation functional

\[
\mathcal E
:
\mathcal S
\rightarrow
\mathbb R_{\ge0},
\]

which determines the branch selected by

\[
\mathcal Q.
\]

\end{definition}

The evaluation functional is external to the compositional chain.

It does not transform configurations but determines the admissibility
branch followed by the correction operator.

Consequently,

the persistence pipeline is not merely a composition of operators,

but a controlled computation driven by defect evaluation.

\section{Soundness of the Persistence Pipeline}

The first fundamental property of the persistence pipeline is that it
never leaves the admissible manifold.

\begin{theorem}[Pipeline Soundness]

Suppose

\[
s
\in
\operatorname{Adm}_{\Phi}.
\]

Then

\[
\Pi(s)
\in
\operatorname{Adm}_{\Phi}.
\]

\end{theorem}

\begin{proof}

Proposal produces a candidate continuation

\[
s'
=
\mathcal P(s).
\]

The evaluation functional computes

\[
d
=
\delta_R(s').
\]

If

\[
d=0,
\]

the correction operator acts as the identity.

Otherwise,

\[
\mathcal Q
\]

constructs an admissible repair.

By definition,

\[
\mathcal Q
(
s'
)
\in
\operatorname{Adm}_{\Phi}.
\]

Finally,

commitment preserves admissibility,

yielding

\[
\Pi(s)
\in
\operatorname{Adm}_{\Phi}.
\]

\end{proof}

\section{Completeness of the Primitive Alphabet}

The primitive persistence operators are not merely sufficient for
constructing admissible computations.

They are universal.

\begin{theorem}[Pipeline Completeness]

Every admissibility-preserving persistence operator

\[
T
\in
\mathfrak T_{\operatorname{Adm}}
\]

admits a representation through the primitive persistence alphabet.

\end{theorem}

\begin{proof}

Let

\[
T
:
\operatorname{Adm}_{\Phi}
\rightarrow
\operatorname{Adm}_{\Phi}
\]

be arbitrary.

Any such transformation necessarily consists of four conceptual stages.

First,

a candidate continuation is generated.

Second,

the candidate is evaluated against the invariant basis through the defect
functional.

Third,

if the candidate violates admissibility,

a corrective transformation is applied.

Finally,

the resulting admissible configuration becomes the current state.

These four stages correspond precisely to

\[
\mathcal P,
\quad
\mathcal E,
\quad
\mathcal Q,
\quad
\mathcal C,
\]

respectively.

Hence every admissible persistence transformation factors through the
primitive alphabet.

\end{proof}

\section{The Persistence Rewrite System}

The Pipeline Completeness Theorem establishes that every admissible
persistence transformation admits a factorization through the primitive
alphabet.

Such factorizations are generally not unique.

Indeed, the same persistence computation may be represented by many
different compositions of primitive operators, differing only by
redundant evaluations, repeated commitments, or unnecessary intermediate
repairs.

To remove this syntactic ambiguity we introduce a rewrite system on the
free monoid generated by the primitive alphabet.

The resulting normal forms provide canonical representatives of
equivalence classes of persistence computations.

\begin{definition}[Persistence Rewrite System]

Let

\[
\mathcal A_{\mathrm{prim}}^*
\]

denote the free monoid generated by the primitive persistence alphabet.

A persistence rewrite system is a collection of reduction rules

\[
W
\longrightarrow
W'
\]

acting on finite words

\[
W\in\mathcal A_{\mathrm{prim}}^*
\]

such that both words represent the same persistence transformation.

\end{definition}

\subsection{Elementary Rewrite Rules}

The persistence calculus is generated by the following elementary
reductions.

\begin{enumerate}

\item
Repeated commitment is redundant:

\[
\mathcal C
\circ
\mathcal C
\;\longrightarrow\;
\mathcal C.
\]

\item
Repeated observation is redundant:

\[
\mathcal O_{\mathrm{obs}}
\circ
\mathcal O_{\mathrm{obs}}
\;\longrightarrow\;
\mathcal O_{\mathrm{obs}}.
\]

\item
Consecutive proposals collapse into a single effective proposal:

\[
\mathcal P_2
\circ
\mathcal P_1
\;\longrightarrow\;
\mathcal P_{\mathrm{eff}}.
\]

\item
Repeated evaluation produces no additional information:

\[
\mathcal E
\circ
\mathcal E
\;\longrightarrow\;
\mathcal E.
\]

\item
Evaluation following commitment is redundant:

\[
\mathcal E
\circ
\mathcal C
\;\longrightarrow\;
\mathcal C.
\]

\end{enumerate}

These reductions express algebraic identities rather than computational
heuristics.

Each preserves the induced persistence morphism.

\begin{definition}[Reduction Length]

For every word

\[
W
\in
\mathcal A_{\mathrm{prim}}^*,
\]

define its reduction length

\[
\ell(W)
\]

to be the number of primitive operators appearing in the expression.

Likewise define the redundancy count

\[
r(W)
\]

to be the total number of reducible adjacent operator pairs.

The complexity measure

\[
\mu(W)
=
\ell(W)+r(W)
\]

is called the rewrite complexity.

\end{definition}

\begin{lemma}[Strong Normalization]

Every reduction sequence generated by the persistence rewrite system
terminates after finitely many steps.

\end{lemma}

\begin{proof}

Each elementary rewrite strictly decreases either the reduction length

\[
\ell(W)
\]

or the redundancy count

\[
r(W).
\]

Consequently,

\[
\mu(W)
=
\ell(W)+r(W)
\]

strictly decreases under every reduction.

Since

\[
\mu(W)\in\mathbb N,
\]

there can exist no infinite descending chain.

Hence every reduction sequence terminates.

\end{proof}

\begin{lemma}[Local Confluence]

Whenever

\[
W
\rightarrow
W_1,
\qquad
W
\rightarrow
W_2,
\]

there exists a word

\[
W_3
\]

such that

\[
W_1
\rightarrow^*
W_3,
\qquad
W_2
\rightarrow^*
W_3.
\]

\end{lemma}

\begin{proof}

The only possible ambiguities arise from overlapping elementary rewrite
rules.

These overlap classes are finite.

Each critical pair admits a common reduction after finitely many
applications of the rewrite rules.

The complete verification is given in Appendix~F.

Therefore the rewrite system is locally confluent.

\end{proof}

\begin{theorem}[Canonical Normal Form]

Every persistence computation

\[
W
\in
\mathcal A_{\mathrm{prim}}^*
\]

reduces to a unique canonical normal form.

\end{theorem}

\begin{proof}

By Lemma~4.X, the persistence rewrite system is terminating.

By Lemma~4.Y, it is locally confluent.

Newman's Lemma therefore implies global confluence.

Hence every word possesses a unique irreducible representative.

This representative is called the canonical persistence normal form.

\end{proof}

\begin{corollary}[Algorithmic Canonicalization]

There exists a deterministic normalization algorithm

\[
\operatorname{NF}
:
\mathcal A_{\mathrm{prim}}^*
\rightarrow
\mathcal A_{\mathrm{prim}}^*
\]

such that

\[
\operatorname{NF}(W)
=
\operatorname{NF}(V)
\]

if and only if

\[
W
\]

and

\[
V
\]

represent the same persistence transformation.

\end{corollary}

\section{Semantic Equivalence of Persistence Computations}

The preceding section established that every persistence computation admits
a unique syntactic normal form.

We now show that this normal form is not merely a convenient symbolic
representation, but determines the semantic action of the computation on
the admissible state space.

Consequently, two persistence computations are regarded as equivalent
precisely when they induce the same admissible transformation.

\begin{definition}[Semantic Equivalence]

Let

\[
W_1,W_2
\in
\mathcal A_{\mathrm{prim}}^*
\]

be persistence words.

We write

\[
W_1
\equiv
W_2
\]

whenever

\[
W_1(s)
=
W_2(s)
\]

for every

\[
s
\in
\operatorname{Adm}_{\Phi}
\]

for which both computations are defined.

The relation

\[
\equiv
\]

is called semantic equivalence.

\end{definition}

\begin{proposition}

Semantic equivalence is an equivalence relation.

\end{proposition}

\begin{proof}

Reflexivity follows immediately since every computation agrees with
itself.

Symmetry follows because equality of functions is symmetric.

Transitivity follows from transitivity of equality.

Therefore

\[
\equiv
\]

is an equivalence relation.

\end{proof}

\begin{theorem}[Correctness of Normalization]

For every persistence computation

\[
W
\in
\mathcal A_{\mathrm{prim}}^*,
\]

the canonical normal form satisfies

\[
W
\equiv
\operatorname{NF}(W).
\]

\end{theorem}

\begin{proof}

Each elementary rewrite preserves the induced persistence
transformation.

Since normalization is obtained by a finite sequence of elementary
rewrites,

semantic equivalence is preserved throughout the reduction.

Therefore

\[
W
\equiv
\operatorname{NF}(W).
\]

\end{proof}

\begin{corollary}

Two persistence computations are semantically equivalent if and only if
they possess the same canonical normal form.

\end{corollary}

\begin{proof}

The forward implication follows from uniqueness of normal forms.

The reverse implication follows immediately from the correctness of
normalization.

\end{proof}

\section{Universality of the Primitive Alphabet}

The primitive persistence alphabet serves as a generating system for the
entire persistence calculus.

The following theorem formalizes this universality.

\begin{theorem}[Universal Generation]

Every admissibility-preserving persistence morphism

\[
T
\in
\mathfrak T_{\operatorname{Adm}}
\]

is generated by the primitive persistence alphabet.

Equivalently,

the evaluation map

\[
\pi
:
\mathcal A_{\mathrm{prim}}^*
\rightarrow
\mathfrak T_{\operatorname{Adm}}
\]

defined by operator composition is surjective.

\end{theorem}

\begin{proof}

Pipeline Completeness establishes that every admissibility-preserving
morphism factors through the primitive operators.

Consequently every

\[
T
\]

lies in the image of

\[
\pi.
\]

Hence

\[
\pi
\]

is surjective.

\end{proof}

\begin{theorem}[Minimality of the Primitive Alphabet]

The primitive persistence alphabet is a minimal generating set for

\[
\mathfrak T_{\operatorname{Adm}}.
\]

\end{theorem}

\begin{proof}

Let

\[
X
\in
\mathcal A_{\mathrm{prim}}.
\]

Removing

\[
X
\]

produces the generating set

\[
\mathcal A_{\mathrm{prim}}
\setminus
\{X\}.
\]

For each primitive operator we exhibit a persistence morphism that
cannot be generated by the reduced alphabet.

\begin{enumerate}

\item
Without the proposal operator, no nontrivial continuation can be
generated.

\item
Without the evaluation functional, admissibility cannot be measured,
making correction undefined.

\item
Without the correction operator, inadmissible trajectories cannot be
returned to the admissible manifold.

\item
Without commitment, admissible continuations cannot become persistent
states.

\end{enumerate}

Therefore no proper subset of the primitive alphabet generates the full
persistence monoid.

\end{proof}

\begin{theorem}[Representation Theorem for Persistence Calculus]

The persistence monoid

\[
\mathfrak T_{\operatorname{Adm}}
\]

is isomorphic to the quotient monoid

\[
\mathcal A_{\mathrm{prim}}^*/
{\equiv},
\]

where

\[
\equiv
\]

denotes semantic equivalence.

Consequently, every admissibility-preserving persistence computation is
represented uniquely by a canonical normal form.

\end{theorem}

\chapter{Repair Geometry}

\section{The Geometry of Repair}

The operator calculus developed in the previous chapter describes how
admissible continuations are generated and manipulated.

The present chapter studies a complementary question.

Rather than asking how persistence operators act, we ask what geometric
structure they induce on the admissible manifold.

The central observation is that repair is not an arbitrary correction.

Instead, every repair operator attempts to move an inadmissible state
toward the admissible manifold while minimizing reconstruction defect.

Consequently, repair naturally acquires a metric and variational
interpretation.

Throughout this chapter we assume that

\[
(\mathcal S,d)
\]

is a metric space together with an admissibility manifold

\[
\operatorname{Adm}_{\Phi}
\subseteq
\mathcal S.
\]

\section{Repair Metric}

The reconstruction defect introduced in Chapter~2 naturally induces a
distance measuring the cost of transforming one configuration into
another through admissible continuation.

\begin{definition}[Repair Metric]

A repair metric is a function

\[
d_R:
\mathcal S\times\mathcal S
\rightarrow
\mathbb R_{\ge0}
\]

satisfying

\begin{enumerate}

\item
Non-negativity,

\[
d_R(x,y)\ge0;
\]

\item
Identity,

\[
d_R(x,y)=0
\iff
x=y;
\]

\item
Symmetry,

\[
d_R(x,y)=d_R(y,x);
\]

\item
Triangle inequality,

\[
d_R(x,z)
\le
d_R(x,y)+d_R(y,z).
\]

\end{enumerate}

\end{definition}

Unlike ordinary Euclidean distance,

the repair metric measures the minimal admissible effort required to
transform one configuration into another.

\begin{definition}[Repair Length]

Let

\[
\sigma:
[0,T]
\rightarrow
\mathcal S
\]

be a piecewise continuously differentiable trajectory.

Its repair length is

\[
L_R(\sigma)
=
\int_0^T
\delta_R(\sigma(t))
\,
\|\dot\sigma(t)\|
\,dt.
\]

\end{definition}

The quantity

\[
L_R(\sigma)
\]

weights geometric length by instantaneous reconstruction defect.

Trajectories remaining near the admissible manifold therefore incur
smaller repair cost than trajectories wandering through highly defective
regions.

\begin{definition}[Repair Distance]

For configurations

\[
x,y
\in
\mathcal S,
\]

define

\[
d_R(x,y)
=
\inf_{\sigma}
L_R(\sigma),
\]

where the infimum is taken over all admissible piecewise smooth
trajectories joining

\[
x
\]

to

\[
y.
\]

\end{definition}

\begin{theorem}[Existence of Shortest Repair Paths]

Assume

\begin{enumerate}

\item
\[
(\mathcal S,d)
\]
is complete,

\item
the admissible manifold is closed,

\item
the reconstruction defect

\[
\delta_R
\]

is lower semicontinuous,

\item
the repair length functional is coercive.

\end{enumerate}

Then every pair of configurations joined by an admissible trajectory
admits a shortest repair path.

\end{theorem}

\begin{proof}

The proof follows the direct method of the calculus of variations.

Choose a minimizing sequence

\[
\{\sigma_n\}
\]

for

\[
L_R.
\]

Coercivity bounds the sequence.

Compactness yields a convergent subsequence.

Lower semicontinuity implies

\[
L_R(\sigma)
\le
\liminf
L_R(\sigma_n),
\]

so the limit realizes the infimum.

\end{proof}

\section{Repair Neighborhoods}

Repair is fundamentally local.

Accordingly we introduce neighborhoods determined by repair distance
rather than ambient geometry.

\begin{definition}[Repair Ball]

For

\[
x
\in
\mathcal S
\]

and

\[
r>0,
\]

define the repair ball

\[
B_R(x,r)
=
\{
y\in\mathcal S:
d_R(x,y)<r
\}.
\]

\end{definition}

Repair balls generally differ from metric balls induced by the ambient
geometry.

Regions separated by high reconstruction defect may be nearby in
Euclidean distance while remaining far apart in repair distance.

\begin{proposition}

Repair balls form a basis for the repair topology.

\end{proposition}

\begin{proof}

Every point belongs to some repair ball.

Furthermore,

if

\[
z
\in
B_R(x,r)
\cap
B_R(y,s),
\]

then there exists

\[
\varepsilon>0
\]

such that

\[
B_R(z,\varepsilon)
\subseteq
B_R(x,r)
\cap
B_R(y,s).
\]

Hence the repair balls satisfy the basis axioms.

\end{proof}

\section{Repair Geodesics}

Shortest repair paths occupy the same role in repair geometry that
geodesics occupy in Riemannian geometry.

\begin{definition}[Repair Geodesic]

A repair geodesic is a trajectory

\[
\gamma
:
[0,1]
\rightarrow
\mathcal S
\]

such that every sufficiently small subarc minimizes repair length.

\end{definition}

Unlike ordinary geodesics,

repair geodesics minimize reconstruction effort rather than geometric
distance.

They therefore follow valleys of low defect instead of straight lines.

%%%%%%%%%%%%%%%%%%%%%%%%%%%%%%%%%%%%%%%%%%%%%%%%%%%%%%%%%%%%%%
\part{Constructive Persistence}

\chapter{Repair Geometry}

\epigraph{%
Persistence is not merely the absence of failure.
It is the geometry by which failure is corrected.
}{}

\section{From Algebra to Geometry}

The preceding chapters established the analytical and algebraic
foundations of persistence.

Chapter~1 introduced the observation map and demonstrated that
projection necessarily destroys the recoverability of field-space
provenance.

Chapter~2 identified the invariant structures whose preservation
defines admissibility.

Chapter~3 showed that persistence belongs to admissible trajectories
rather than isolated configurations.

Finally, Chapter~4 constructed an operator calculus in which every
admissible computation factors through a canonical persistence
pipeline.

The present chapter changes perspective.

Rather than studying operators themselves, we study the geometric
structure induced by their action.

The central idea is that repair is not an arbitrary transformation.

Instead, every repair operator attempts to minimize reconstruction
defect while remaining inside the admissible manifold.

Repair therefore possesses an intrinsic geometry.

The purpose of this chapter is to develop that geometry.

\section{The Principle of Minimal Repair}

Every persistence operator produces a trajectory through the
configuration manifold.

Among all admissible continuations, certain trajectories require
less reconstruction effort than others.

These trajectories represent the natural evolution of the system.

This observation motivates the following principle.

\begin{axiom}[Minimal Repair Principle]

Whenever multiple admissible repairs exist, persistence selects a
trajectory minimizing accumulated reconstruction defect.

\[
\sigma^*
=
\operatorname*{arg\,min}_{\sigma}
\mathcal D[\sigma].
\]

\end{axiom}

The Minimal Repair Principle is not an optimization heuristic.

Rather, it provides the geometric interpretation of persistence.

Admissible systems evolve along trajectories of minimal repair cost.

\section{Repair Energy}

The reconstruction functional introduced previously may now be
interpreted geometrically.

\begin{definition}[Repair Energy]

Let

\[
\sigma:[0,T]\rightarrow\mathcal S
\]

be an admissible trajectory.

Its repair energy is

\[
E_R(\sigma)
=
\int_0^T
\delta_R(\sigma(t))
\,dt.
\]

\end{definition}

Unlike the repair length,

which weights geometric motion,

repair energy measures the total accumulated reconstruction effort
experienced during continuation.

Low-energy trajectories therefore correspond to naturally stable
persistent evolution.

\begin{proposition}

Every repair geodesic minimizes repair energy locally.

\end{proposition}

\begin{proof}

Suppose

\[
\gamma
\]

is a repair geodesic.

By definition,

every sufficiently short subarc minimizes repair length.

Since repair energy differs from repair length only by
reparameterization of admissible motion,

any local perturbation increases accumulated defect.

Therefore

\[
\gamma
\]

is locally energy minimizing.

\end{proof}

\section{Repair Gradient Flow}

Repair operators admit a continuous approximation.

Rather than applying discrete corrections,

one may evolve continuously in the direction of decreasing defect.

\begin{definition}[Repair Gradient Flow]

The repair gradient flow is defined by

\[
\frac{ds}{dt}
=
-
\nabla
\delta_R(s).
\]

Whenever

\[
\delta_R
\]

is continuously differentiable,

the resulting trajectory decreases reconstruction defect
monotonically.

\end{definition}

\begin{theorem}[Monotonicity of Repair]

Along every repair gradient flow,

\[
\delta_R(s(t))
\]

is non-increasing.

\end{theorem}

\begin{proof}

Differentiating,

\[
\frac{d}{dt}
\delta_R(s(t))
=
\nabla\delta_R(s)
\cdot
\frac{ds}{dt}.
\]

Substituting the repair flow yields

\[
\frac{d}{dt}
\delta_R(s(t))
=
-
\|\nabla\delta_R(s)\|^2
\le
0.
\]

Hence reconstruction defect decreases monotonically.

\end{proof}

\section{Repair Basins}

The repair gradient induces a decomposition of the admissible
manifold.

\begin{definition}[Repair Basin]

The repair basin associated with an admissible equilibrium

\[
s^*
\]

is

\[
\mathcal B(s^*)
=
\{
x\in\mathcal S
:
\lim_{t\rightarrow\infty}
\Phi_t(x)
=
s^*
\},
\]

where

\[
\Phi_t
\]

denotes the repair flow.

\end{definition}

Repair basins partition the configuration manifold into regions
sharing the same asymptotic repaired configuration.

They provide the geometric analogue of attraction basins in
dynamical systems while remaining defined through admissibility
rather than stability alone.

\section{Toward Recursive Continuation}

Repair geometry completes the transition from the analytical
foundations established in Part~I to the recursive systems
developed in the following chapter.

Observation determines what survives projection.

Invariant theory determines what must be preserved.

Continuation determines which trajectories remain admissible.

Operator calculus determines how admissible continuations are
constructed.

Repair geometry determines why certain continuations are naturally
preferred.

The remaining question is no longer geometric but recursive.

Can a persistence system repeatedly repair itself while preserving
admissibility?

Chapter~6 develops the mathematical theory of recursive
continuation and answers this question through fixed-point theory,
recursive repair operators, and self-maintaining persistence
architectures.

\section{Convexity of Repair Landscapes}

The geometry induced by reconstruction defect admits additional
structure whenever the defect functional satisfies suitable convexity
properties.

Convexity guarantees that repair trajectories are well behaved,
eliminates spurious local minima, and ensures that admissible repairs
vary continuously under perturbation.

\begin{definition}[Repair Convexity]

A reconstruction defect functional

\[
\delta_R :
\mathcal S
\rightarrow
\mathbb R_{\ge0}
\]

is said to be repair-convex if, for every admissible geodesic

\[
\gamma:[0,1]\rightarrow\mathcal S,
\]

the inequality

\[
\delta_R(\gamma(t))
\le
(1-t)\delta_R(\gamma(0))
+
t\delta_R(\gamma(1))
\]

holds for every

\[
0\le t\le1.
\]

\end{definition}

Repair convexity implies that interpolation between admissible
configurations never introduces defect exceeding the linear
interpolation of endpoint defects.

Consequently, the defect landscape possesses no interior barriers
along admissible geodesics.

\begin{theorem}[Uniqueness of Local Repair]

Suppose the reconstruction defect functional is strictly repair-convex.

Then every sufficiently small repair neighborhood possesses a unique
minimal repair trajectory.

\end{theorem}

\begin{proof}

Assume two distinct minimizing repair trajectories connect the same
pair of nearby admissible configurations.

Strict repair convexity implies that every convex interpolation
between the two trajectories possesses strictly smaller accumulated
defect.

This contradicts the assumed minimality.

Hence the minimizing repair trajectory is unique.

\end{proof}

\section{Repair Projection Operators}

Repair may alternatively be viewed as a projection onto the admissible
manifold.

Rather than modifying trajectories arbitrarily, the repair operator
returns the nearest admissible continuation according to the repair
metric.

\begin{definition}[Repair Projection]

The repair projection is the mapping

\[
\Pi_R:
\mathcal S
\rightharpoonup
\operatorname{Adm}_{\Phi}
\]

defined by

\[
\Pi_R(s)
=
\operatorname*{arg\,min}_{x\in\operatorname{Adm}_{\Phi}}
d_R(s,x),
\]

whenever the minimizer exists.

\end{definition}

Unlike ordinary orthogonal projection in Euclidean geometry, the repair
projection depends upon the geometry induced by reconstruction defect
and therefore reflects admissibility rather than ambient distance.

\begin{proposition}

If the admissible manifold is closed and the repair metric is complete,
then every sufficiently small neighborhood admits a unique repair
projection.

\end{proposition}

\begin{proof}

Completeness guarantees existence of minimizing sequences.

Closedness of

\[
\operatorname{Adm}_{\Phi}
\]

ensures that every convergent minimizing sequence converges to an
admissible point.

Local repair convexity establishes uniqueness.

Therefore the repair projection exists and is unique.

\end{proof}

\section{Stability Under Perturbation}

An admissible persistence system should vary continuously under small
perturbations.

Repair stability formalizes this intuitive requirement.

\begin{definition}[Repair Stability]

A repair operator

\[
\mathcal R
\]

is stable if, for every

\[
\varepsilon>0,
\]

there exists

\[
\delta>0
\]

such that

\[
d_R(x,y)<\delta
\]

implies

\[
d_R(\mathcal R(x),\mathcal R(y))
<
\varepsilon.
\]

\end{definition}

\begin{theorem}[Continuity of Repair]

Suppose

\[
\delta_R
\]

is continuous and locally Lipschitz.

Then the repair projection

\[
\Pi_R
\]

is continuous on every neighborhood in which the repair minimizer is
unique.

\end{theorem}

\begin{proof}

Let

\[
x_n\rightarrow x.
\]

By uniqueness of the repair projection,

every convergent subsequence of

\[
\Pi_R(x_n)
\]

must converge to

\[
\Pi_R(x).
\]

Since every subsequence possesses the same limit,

the entire sequence converges.

Therefore

\[
\Pi_R
\]

is continuous.

\end{proof}

\section{Composition of Repairs}

Repair operators naturally compose.

The resulting compositions remain admissible under mild analytical
assumptions.

\begin{proposition}

Let

\[
\mathcal R_1,
\mathcal R_2
:
\mathcal S
\rightharpoonup
\operatorname{Adm}_{\Phi}
\]

be repair operators.

If both preserve admissibility, then

\[
\mathcal R_2
\circ
\mathcal R_1
\]

is again a repair operator.

\end{proposition}

\begin{proof}

Since

\[
\mathcal R_1
(\mathcal S)
\subseteq
\operatorname{Adm}_{\Phi},
\]

and

\[
\mathcal R_2
(
\operatorname{Adm}_{\Phi}
)
\subseteq
\operatorname{Adm}_{\Phi},
\]

composition yields

\[
(\mathcal R_2
\circ
\mathcal R_1)
(\mathcal S)
\subseteq
\operatorname{Adm}_{\Phi},
\]

establishing admissibility preservation.

\end{proof}

\section{The Repair Semigroup}

Repeated repair naturally generates a discrete dynamical system.

\begin{definition}[Repair Semigroup]

Let

\[
\mathcal R
\]

be a repair operator.

The associated repair semigroup is

\[
\{
\mathcal R^n
:
n\in\mathbb N_0
\},
\]

where

\[
\mathcal R^0
=
\operatorname{id}_{\mathcal S},
\qquad
\mathcal R^{n+1}
=
\mathcal R
\circ
\mathcal R^n.
\]

\end{definition}

The repair semigroup describes repeated correction of an evolving
configuration and forms the bridge between static repair geometry and
the recursive continuation theory developed in the following chapter.

\begin{theorem}[Closure of the Repair Semigroup]

If

\[
\mathcal R
\]

preserves admissibility, then

\[
\mathcal R^n
\]

preserves admissibility for every

\[
n\ge0.
\]

\end{theorem}

\begin{proof}

The proof proceeds by induction.

The identity operator preserves admissibility.

Assume

\[
\mathcal R^n
\]

preserves admissibility.

Then

\[
\mathcal R^{n+1}
=
\mathcal R
\circ
\mathcal R^n
\]

is the composition of two admissibility-preserving operators.

By Proposition~5.X, composition preserves admissibility.

Therefore every iterate

\[
\mathcal R^n
\]

belongs to the repair semigroup.

\end{proof}

\section{Summary}

Repair geometry completes the geometric foundation of persistence.

Observation projects high-dimensional fields onto reduced
configurations.

Invariant theory determines admissibility.

Continuation establishes the existence of persistent trajectories.

Operator calculus constructs admissible transformations.

Repair geometry equips those transformations with an intrinsic metric,
variational, and topological structure.

The next chapter extends these ideas from single repairs to recursive
self-maintaining persistence systems, where repair operators become
iterated processes and admissibility evolves through recursive
continuation.

\chapter{Recursive Continuation}

\epigraph{%
Persistence is not the existence of a repair.
Persistence is the indefinite continuation of repair.
}{}

\section{From Repair to Recursion}

The preceding chapter established the geometry of repair.

A repair operator transforms inadmissible configurations into
admissible configurations while minimizing reconstruction defect.

However, persistence rarely consists of a single repair event.

Living systems, computational systems, physical systems, and social
institutions continually experience perturbations.

Persistence therefore requires repeated repair.

This observation motivates the transition from repair geometry to
recursive continuation.

Rather than studying isolated repair operators,

we study operators acting repeatedly upon their own outputs.

Recursive continuation is thus the mathematical theory of
self-maintaining persistence.

Throughout this chapter,

\[
\mathcal R
:
\mathcal S
\rightharpoonup
\operatorname{Adm}_{\Phi}
\]

denotes an admissibility-preserving repair operator.

\section{Recursive Continuation Operators}

Repeated application of a repair operator generates a discrete
evolution.

\begin{definition}[Recursive Continuation]

A recursive continuation operator is a mapping

\[
\mathcal T
:
\operatorname{Adm}_{\Phi}
\rightarrow
\operatorname{Adm}_{\Phi}
\]

whose iterates satisfy

\[
s_{n+1}
=
\mathcal T(s_n),
\qquad
n\ge0.
\]

The sequence

\[
\{s_n\}_{n=0}^{\infty}
\]

is called a continuation trajectory.

\end{definition}

Recursive continuation differs fundamentally from ordinary iteration.

The objective is not repeated computation,

but indefinite preservation of admissibility.

\begin{definition}[Recursive Admissibility]

A continuation trajectory

\[
\{s_n\}
\]

is recursively admissible whenever

\[
s_n
\in
\operatorname{Adm}_{\Phi}
\]

for every

\[
n\ge0.
\]

\end{definition}

\section{Invariant Recursive Systems}

Recursive continuation naturally preserves invariant structure.

\begin{definition}[Invariant Recursive Operator]

A recursive continuation operator

\[
\mathcal T
\]

is invariant whenever

\[
\Phi_i(\mathcal T(s))
=
\Phi_i(s)
\]

for every invariant

\[
\Phi_i
\in
\mathbf{\Phi}_{\mathrm{base}}.
\]

\end{definition}

Invariant recursive operators evolve the system while preserving the
structural constraints established in Chapter~2.

\begin{proposition}

If

\[
\mathcal T
\]

is invariant,

then every recursive continuation trajectory preserves the invariant
basis.

\end{proposition}

\begin{proof}

The result follows immediately by induction.

The base case holds by hypothesis.

Assume

\[
\Phi_i(s_n)
=
\Phi_i(s_0).
\]

Then

\[
\Phi_i(s_{n+1})
=
\Phi_i(\mathcal T(s_n))
=
\Phi_i(s_n)
=
\Phi_i(s_0).
\]

Therefore every iterate preserves every invariant.

\end{proof}

\section{Fixed Points}

The simplest recursively persistent systems eventually become
stationary.

\begin{definition}[Persistence Fixed Point]

A configuration

\[
s^*
\]

is a persistence fixed point whenever

\[
\mathcal T(s^*)
=
s^*.
\]

\end{definition}

Fixed points represent completely self-maintaining configurations.

No additional repair is required.

\begin{theorem}[Existence of Fixed Points]

Suppose

\[
(\operatorname{Adm}_{\Phi},d_R)
\]

is complete and

\[
\mathcal T
\]

is a contraction.

Then

\[
\mathcal T
\]

possesses a unique persistence fixed point.

\end{theorem}

\begin{proof}

The result follows immediately from the Banach Fixed Point Theorem.

The completeness of

\[
\operatorname{Adm}_{\Phi}
\]

guarantees convergence of the iterates,

while contraction ensures uniqueness.

\end{proof}

\section{Monotone Recursive Systems}

Many persistence systems are naturally ordered rather than metric.

Examples include logical knowledge bases,

institutional policies,

and recursive admissibility lattices.

These systems are governed by monotone operators.

\begin{definition}[Monotone Continuation Operator]

Let

\[
(L,\le)
\]

be a complete lattice.

A continuation operator

\[
\mathcal T:L\rightarrow L
\]

is monotone whenever

\[
x\le y
\Longrightarrow
\mathcal T(x)
\le
\mathcal T(y).
\]

\end{definition}

\begin{theorem}[Knaster--Tarski Persistence Theorem]

Every monotone continuation operator on a complete lattice possesses a
least fixed point and a greatest fixed point.

\end{theorem}

\begin{proof}

The theorem follows directly from the
Knaster--Tarski Fixed Point Theorem.

The complete lattice guarantees existence of arbitrary meets and joins.

Monotonicity ensures closure under iterative approximation.

Consequently,

least and greatest persistence fixed points exist.

\end{proof}

\section{Recursive Stability}

Not every recursively admissible trajectory is stable.

We therefore distinguish persistence from stability.

\begin{definition}[Recursive Stability]

A persistence fixed point

\[
s^*
\]

is recursively stable if every sufficiently small perturbation

\[
s_0
\]

generates a continuation trajectory satisfying

\[
\lim_{n\rightarrow\infty}
s_n
=
s^*.
\]

\end{definition}

Recursive stability describes asymptotic recovery after finite
perturbations.

\begin{theorem}[Recursive Recovery]

Suppose

\[
\mathcal T
\]

is locally contractive within a repair neighborhood

\[
B_R(s^*,r).
\]

Then every continuation trajectory beginning inside

\[
B_R(s^*,r)
\]

converges to

\[
s^*.
\]

\end{theorem}

\begin{proof}

Since

\[
\mathcal T
\]

acts as a contraction on the repair neighborhood,

Banach's theorem applies locally.

Therefore every iterate converges uniquely toward

\[
s^*.
\]

\end{proof}

\section{Recursive Defect Dynamics}

Reconstruction defect evolves alongside recursive continuation.

\begin{definition}[Recursive Defect Sequence]

Associated with every continuation trajectory is the sequence

\[
d_n
=
\delta_R(s_n).
\]

\end{definition}

The recursive defect sequence measures the changing reconstruction cost
throughout the continuation process.

\begin{proposition}

Suppose

\[
\mathcal T
\]

is defect-reducing.

Then

\[
\{d_n\}
\]

is monotonically non-increasing.

\end{proposition}

\begin{proof}

For every iterate,

\[
\delta_R(\mathcal T(s))
\le
\delta_R(s).
\]

Applying this inequality recursively yields

\[
d_{n+1}
\le
d_n.
\]

Hence the defect sequence decreases monotonically.

\end{proof}

\section{Hierarchies of Recursive Continuation}

The preceding sections considered recursive continuation generated by a
single persistence operator.

Complex systems, however, rarely evolve through a single recursive
mechanism.

Instead, persistence emerges through hierarchies of interacting repair
processes operating across multiple organizational scales.

This motivates the introduction of recursive continuation hierarchies.

\begin{definition}[Recursive Continuation Hierarchy]

A recursive continuation hierarchy is a family of operators

\[
\{
\mathcal T_i
\}_{i\in I}
\]

indexed by an ordered set

\[
I,
\]

together with admissibility-preserving transition maps

\[
\eta_{ij}
:
\mathcal T_i
\Longrightarrow
\mathcal T_j,
\]

whenever

\[
i\le j.
\]

\end{definition}

Each level repairs the persistence defects left unresolved by lower
levels while preserving the invariant structures inherited from higher
levels.

Consequently, persistence becomes distributed rather than localized.

\section{Nested Admissibility}

Hierarchical repair naturally induces nested admissibility manifolds.

\begin{definition}[Nested Admissibility]

A hierarchy of admissibility manifolds is a sequence

\[
\operatorname{Adm}_{\Phi}^{(0)}
\supseteq
\operatorname{Adm}_{\Phi}^{(1)}
\supseteq
\cdots
\supseteq
\operatorname{Adm}_{\Phi}^{(n)}
\]

such that each inclusion is proper and each level imposes additional
invariant constraints.

\end{definition}

Lower levels guarantee coarse persistence.

Higher levels enforce progressively stronger notions of structural
consistency.

\begin{proposition}

Every recursively admissible trajectory on

\[
\operatorname{Adm}_{\Phi}^{(k)}
\]

is recursively admissible on every coarser level

\[
\operatorname{Adm}_{\Phi}^{(j)},
\qquad
j<k.
\]

\end{proposition}

\begin{proof}

Since

\[
\operatorname{Adm}_{\Phi}^{(k)}
\subseteq
\operatorname{Adm}_{\Phi}^{(j)},
\]

membership in the finer admissibility manifold immediately implies
membership in every coarser manifold.

\end{proof}

\section{Coupled Continuation Operators}

Persistence frequently depends upon multiple interacting recursive
systems.

The evolution of one continuation operator may influence another,
producing coupled persistence dynamics.

\begin{definition}[Coupled Continuation System]

A coupled continuation system consists of operators

\[
\mathcal T_1,\ldots,\mathcal T_n
\]

acting on configuration spaces

\[
\mathcal S_1,\ldots,\mathcal S_n
\]

according to

\[
s_i^{(k+1)}
=
\mathcal T_i
\!\left(
s_1^{(k)},
\ldots,
s_n^{(k)}
\right).
\]

\end{definition}

The coupled system evolves on the product configuration manifold

\[
\mathcal S
=
\prod_{i=1}^{n}
\mathcal S_i.
\]

\begin{definition}[Joint Admissibility]

A coupled configuration

\[
(s_1,\ldots,s_n)
\]

is jointly admissible whenever

\[
s_i
\in
\operatorname{Adm}_{\Phi_i}
\]

for every subsystem and all compatibility constraints between
subsystems are simultaneously satisfied.

\end{definition}

Joint admissibility is generally stronger than componentwise
admissibility.

Interactions themselves may introduce persistence constraints.

\begin{theorem}[Existence of Coupled Persistence]

Suppose every continuation operator

\[
\mathcal T_i
\]

is continuous and admissibility preserving.

If the joint admissibility manifold is nonempty and compact, then the
coupled continuation system possesses at least one jointly admissible
continuation trajectory.

\end{theorem}

\begin{proof}

The product operator

\[
\mathcal T
=
(\mathcal T_1,\ldots,\mathcal T_n)
\]

acts continuously on the compact joint admissibility manifold.

Existence of an admissible continuation follows from the continuation
results established in Chapter~3 applied to the product system.

\end{proof}

\section{Multiscale Persistence}

Repair frequently occurs simultaneously across multiple temporal and
spatial scales.

Rapid local corrections coexist with slower global adaptations.

\begin{definition}[Multiscale Continuation]

A multiscale continuation system consists of operators

\[
\{
\mathcal T_{\tau}
\}_{\tau\in\Lambda}
\]

indexed by characteristic scales

\[
\tau.
\]

Each operator governs persistence dynamics at its associated scale.

\end{definition}

The complete persistence process is obtained by composition of the
multiscale continuation operators.

\begin{proposition}

If every scale operator preserves admissibility, then every finite
multiscale composition preserves admissibility.

\end{proposition}

\begin{proof}

Each operator maps admissible configurations into admissible
configurations.

Repeated application therefore preserves admissibility by induction.

\end{proof}

\section{Persistence Attractors}

Recursive continuation naturally generates invariant sets.

\begin{definition}[Persistence Attractor]

A persistence attractor is a closed invariant subset

\[
A
\subseteq
\operatorname{Adm}_{\Phi}
\]

such that every sufficiently nearby recursively admissible trajectory
approaches

\[
A
\]

under repeated continuation.

\end{definition}

Persistence attractors generalize fixed points.

Rather than converging to a single configuration, recursive systems may
stabilize upon invariant manifolds, periodic trajectories, or more
general compact invariant sets.

\begin{theorem}[Persistence of Attractors]

Suppose the continuation operator is continuous and maps a compact
admissible manifold into itself.

Then every persistence attractor is positively invariant.

\end{theorem}

\begin{proof}

Let

\[
A
\]

be a persistence attractor.

Continuity implies

\[
\mathcal T(A)
\]

is compact.

Since every trajectory beginning in

\[
A
\]

remains asymptotically inside

\[
A,
\]

positive invariance follows.

\end{proof}

\section{Recursive Persistence Principle}

The preceding constructions may be summarized by a single structural
principle.

\begin{theorem}[Recursive Persistence Principle]

Persistence is recursively maintained whenever there exists a hierarchy
of admissibility-preserving continuation operators whose joint action
strictly prevents unbounded growth of reconstruction defect while
preserving the invariant basis.

\end{theorem}

\begin{proof}

Continuation preserves admissibility by construction.

Repair bounds reconstruction defect.

Invariant preservation guarantees structural consistency.

Composition of admissibility-preserving operators yields recursively
admissible trajectories.

Therefore the hierarchy maintains persistence indefinitely unless the
admissibility manifold itself becomes empty.

\end{proof}

\section{Transition to Operator Dynamics}

Recursive continuation establishes the existence and stability of
self-maintaining persistence systems.

The remaining question concerns the operators themselves.

Up to this point, persistence operators have been regarded as fixed
objects.

In many systems, however, the operators evolve alongside the
configurations they govern.

Learning systems modify their repair rules.

Adaptive controllers refine their correction strategies.

Scientific theories revise their own inference procedures.

Institutions update their decision mechanisms.

This motivates the transition from recursive continuation to
\emph{operator dynamics}.

Chapter~7 develops the continuous evolution of persistence operators,
introducing operator semigroups, infinitesimal generators, and
continuous flows on the space of admissible transformations.

\chapter{Operator Dynamics}

\epigraph{%
Persistence is not only the evolution of states.
Persistence is the evolution of the rules governing states.
}{}

\section{Dynamic Persistence Operators}

The previous chapter established recursive continuation through
iterated application of fixed persistence operators.

Many natural systems, however, modify their own repair mechanisms over
time.

Learning algorithms update parameters.

Biological organisms adapt regulatory pathways.

Institutions revise procedures.

Scientific theories refine explanatory frameworks.

Persistence therefore requires a mathematical theory in which the
operators themselves evolve.

Throughout this chapter,

\[
\mathfrak P
\]

denotes the space of admissibility-preserving persistence operators.

Each point of

\[
\mathfrak P
\]

represents a complete persistence pipeline rather than an individual
configuration.

Operator dynamics studies trajectories inside

\[
\mathfrak P.
\]

\section{Operator Trajectories}

\begin{definition}[Operator Trajectory]

An operator trajectory is a continuous mapping

\[
\Gamma :
[0,\infty)
\rightarrow
\mathfrak P.
\]

The value

\[
\Gamma(t)
\]

represents the persistence operator active at time

\[
t.
\]

\end{definition}

Unlike continuation trajectories,

whose points are configurations,

operator trajectories consist entirely of persistence operators.

The objects undergoing evolution have therefore shifted one level
higher.

\section{Operator Topology}

To discuss convergence of persistence operators, the operator space
must itself possess a topology.

\begin{definition}[Operator Convergence]

A sequence

\[
\{\mathcal T_n\}
\]

converges strongly to

\[
\mathcal T
\]

if

\[
\lim_{n\to\infty}
\mathcal T_n(s)
=
\mathcal T(s)
\]

for every admissible configuration

\[
s.
\]

\end{definition}

Alternative notions such as uniform convergence or convergence in
operator norm may be employed when additional structure exists.

\section{Operator Semigroups}

Continuous operator evolution naturally satisfies the semigroup
property.

\begin{definition}[Persistence Semigroup]

A family

\[
\{
U_t
\}_{t\ge0}
\]

of persistence operators is called a persistence semigroup whenever

\[
U_0
=
I,
\]

and

\[
U_{t+s}
=
U_t
\circ
U_s
\]

for every

\[
s,t\ge0.
\]

\end{definition}

The semigroup property expresses temporal consistency.

Evolving first for

\[
s
\]

units of time and then for

\[
t
\]

units produces the same result as evolving directly for

\[
s+t.
\]

\begin{theorem}

The persistence semigroups form a category whose morphisms are
semigroup homomorphisms preserving admissibility.

\end{theorem}

\begin{proof}

Identity semigroups exist.

Composition of semigroup homomorphisms preserves both composition and
identity.

Associativity follows immediately from composition of maps.

Hence persistence semigroups form a category.

\end{proof}

\section{Infinitesimal Generators}

Continuous evolution may be described locally through infinitesimal
change.

\begin{definition}[Generator]

The infinitesimal generator of a persistence semigroup is

\[
A
=
\lim_{h\to0}
\frac{U_h-I}{h},
\]

whenever the limit exists.

\end{definition}

The generator determines the instantaneous direction in which
persistence operators evolve.

Knowledge of

\[
A
\]

is sufficient to reconstruct the semigroup under appropriate
analytical hypotheses.

\begin{definition}[Operator Evolution Equation]

The evolution of persistence operators satisfies

\[
\frac{dU_t}{dt}
=
A\,U_t,
\]

subject to

\[
U_0
=
I.
\]

\end{definition}

This equation is the operator analogue of ordinary differential
equations governing state evolution.

\section{Admissibility Preservation Under Flow}

Not every operator flow preserves persistence.

The generator itself must satisfy compatibility conditions with the
admissibility manifold.

\begin{definition}[Admissible Generator]

A generator

\[
A
\]

is admissible whenever every operator

\[
U_t
\]

generated by

\[
A
\]

preserves admissibility.

\end{definition}

\begin{theorem}

Suppose

\[
A
\]

is an admissible generator.

Then

\[
U_t
\]

preserves admissibility for every

\[
t\ge0.
\]

\end{theorem}

\begin{proof}

The proof follows from the definition of admissible generators together
with the semigroup evolution equation.

Every infinitesimal step remains inside the admissible operator space.

Closure under composition implies admissibility for finite evolution
times.

\end{proof}

\section{Operator Stability}

Persistence operators should evolve continuously under small
perturbations.

\begin{definition}[Operator Stability]

An operator trajectory

\[
\Gamma
\]

is stable whenever, for every

\[
\varepsilon>0,
\]

there exists

\[
\delta>0
\]

such that

\[
d_{\mathfrak P}
(
\Gamma(0),
\widetilde{\Gamma}(0)
)
<
\delta
\]

implies

\[
d_{\mathfrak P}
(
\Gamma(t),
\widetilde{\Gamma}(t)
)
<
\varepsilon
\]

for all sufficiently small

\[
t.
\]

\end{definition}

Operator stability guarantees robustness of the persistence mechanism
itself rather than merely the trajectories it produces.

\section{Lyapunov Functionals on Operator Space}

The defect functional may be lifted from configuration space to the
space of persistence operators.

\begin{definition}[Operator Lyapunov Functional]

A functional

\[
\mathcal L
:
\mathfrak P
\rightarrow
\mathbb R_{\ge0}
\]

is an operator Lyapunov functional if

\[
\frac{d}{dt}
\mathcal L(U_t)
\le
0
\]

along every admissible operator flow.

\end{definition}

Operator Lyapunov functionals measure the long-term stability of repair
strategies rather than individual repaired configurations.

They provide sufficient conditions for convergence of adaptive
persistence systems.

\section{Adaptive Persistence}

Operator evolution permits systems to improve their own repair
mechanisms.

\begin{definition}[Adaptive Persistence]

A persistence system is adaptive whenever its operator trajectory

\[
\Gamma(t)
\]

depends upon accumulated continuation history.

\[
\Gamma(t)
=
F\!\left(
\Gamma(0),
H_t
\right),
\]

where

\[
H_t
\]

denotes the persistence history up to time

\[
t.
\]

\end{definition}

Adaptive persistence generalizes recursive continuation by allowing the
repair operator itself to become an evolving state variable.

The resulting systems no longer merely survive perturbation.

They continually refine the mechanisms by which survival is achieved.

\section{Operator Manifolds}

The collection of admissibility-preserving operators possesses more
structure than that of a mere set.

Under appropriate regularity assumptions, it may itself be regarded as
a differentiable manifold.

\begin{definition}[Operator Manifold]

An operator manifold is a smooth manifold

\[
\mathfrak P
\]

whose points are admissibility-preserving persistence operators and
whose tangent spaces describe infinitesimal operator perturbations.

\end{definition}

Local coordinates on

\[
\mathfrak P
\]

represent admissible deformations of persistence pipelines while
maintaining the invariant constraints established in previous
chapters.

Operator evolution therefore becomes a curve on the manifold

\[
\mathfrak P.
\]

\section{Tangent Operators}

Infinitesimal changes in persistence are represented by tangent
vectors.

\begin{definition}[Tangent Operator]

Let

\[
\Gamma(t)
\]

be an operator trajectory.

The tangent operator at time

\[
t
\]

is

\[
\dot{\Gamma}(t)
=
\frac{d}{dt}\Gamma(t).
\]

\end{definition}

The tangent operator measures the instantaneous modification of the
repair mechanism itself.

Whereas ordinary derivatives describe changing configurations,

tangent operators describe changing rules.

\section{Operator Vector Fields}

Just as vector fields generate flows on configuration manifolds,
operator vector fields generate flows on operator manifolds.

\begin{definition}[Operator Vector Field]

An operator vector field is a smooth assignment

\[
X:
\mathfrak P
\rightarrow
T\mathfrak P
\]

mapping each persistence operator to an infinitesimal admissible
deformation.

\end{definition}

The corresponding evolution equation is

\[
\frac{d\Gamma}{dt}
=
X(\Gamma).
\]

Solutions describe continuously adapting persistence mechanisms.

\section{Commuting Operator Flows}

Multiple adaptive processes frequently act simultaneously.

Their interaction depends upon whether their infinitesimal generators
commute.

\begin{definition}[Commuting Flows]

Two operator vector fields

\[
X
\quad\text{and}\quad
Y
\]

commute whenever

\[
[X,Y]
=
0,
\]

where

\[
[\cdot,\cdot]
\]

denotes the Lie bracket.

\end{definition}

Commuting flows may be applied in arbitrary order without changing the
resulting persistence operator.

\begin{proposition}

If

\[
[X,Y]=0,
\]

then

\[
e^{tX}e^{sY}
=
e^{sY}e^{tX}
\]

for all sufficiently small

\[
s,t.
\]

\end{proposition}

\begin{proof}

The result follows from standard Lie theory.

Vanishing Lie bracket implies that the corresponding local flows
commute.

Consequently, their compositions are independent of order.

\end{proof}

\section{Noncommuting Adaptation}

Most persistence mechanisms do not commute.

Learning before repair may produce a different outcome than repairing
before learning.

Likewise, updating an invariant basis prior to modifying observation
may differ fundamentally from performing these operations in reverse
order.

Noncommutativity therefore represents genuine path dependence within
operator space.

\begin{definition}[Operator Curvature]

The operator curvature associated with vector fields

\[
X
\]

and

\[
Y
\]

is measured by their Lie bracket

\[
[X,Y].
\]

Vanishing curvature corresponds to locally order-independent
adaptation.

Nonvanishing curvature measures the failure of adaptive processes to
commute.

\end{definition}

\section{Operator Geodesics}

Not every operator evolution is equally efficient.

Some adaptive trajectories minimize accumulated modification while
preserving admissibility.

\begin{definition}[Operator Geodesic]

An operator trajectory

\[
\Gamma
\]

is an operator geodesic if it is locally minimizing with respect to an
operator metric

\[
g_{\mathfrak P}.
\]

\end{definition}

Operator geodesics represent the most economical continuous
transformation between persistence mechanisms.

They generalize repair geodesics from configuration space to operator
space.

\section{Bifurcations of Persistence}

Adaptive persistence systems may undergo qualitative changes in
behavior.

Such transitions occur when small parameter variations produce
structurally different operator dynamics.

\begin{definition}[Persistence Bifurcation]

A persistence bifurcation occurs whenever a continuous variation of an
operator parameter changes the qualitative structure of admissible
continuation trajectories.

\end{definition}

Typical examples include the creation or destruction of persistence
fixed points, the merging of repair basins, or transitions between
stable and unstable recursive continuation.

\begin{theorem}[Structural Stability]

Suppose an operator trajectory avoids persistence bifurcations on an
interval

\[
[a,b].
\]

Then every sufficiently small perturbation of the initial operator
produces a topologically equivalent persistence evolution on the same
interval.

\end{theorem}

\begin{proof}

In the absence of bifurcations, qualitative properties of the flow vary
continuously under sufficiently small perturbations.

Consequently, the induced persistence dynamics remain topologically
equivalent throughout the interval.

\end{proof}

\section{Persistence Phase Space}

The combined evolution of configurations and persistence operators
defines a higher-dimensional dynamical system.

\begin{definition}[Persistence Phase Space]

The persistence phase space is the product manifold

\[
\mathcal X
=
\mathcal S
\times
\mathfrak P.
\]

A point

\[
(s,\mathcal T)
\]

specifies both the current configuration and the persistence operator
governing its evolution.

\end{definition}

Evolution on

\[
\mathcal X
\]

simultaneously describes changing systems and changing rules.

The dynamics therefore become self-referential without requiring the
operator to act directly upon itself.

\section{The Evolution Principle}

The preceding constructions may be summarized by a single organizing
principle.

\begin{theorem}[Evolution Principle]

Let

\[
(\mathcal S,\mathfrak P)
\]

be a persistence phase space equipped with admissibility-preserving
operator dynamics.

If

\begin{enumerate}

\item every configuration trajectory remains admissible,

\item every operator trajectory preserves admissibility,

\item reconstruction defect remains uniformly bounded, and

\item operator evolution is continuous,

\end{enumerate}

then the coupled evolution defines a well-posed persistence dynamical
system.

\end{theorem}

\begin{proof}

The first hypothesis guarantees admissible continuation on
configuration space.

The second guarantees admissible evolution of the governing operators.

The third prevents divergence of reconstruction cost.

The fourth ensures continuous dependence on initial conditions.

Together these establish existence, consistency, and stability of the
combined evolution on the persistence phase space.

\end{proof}

\section{Transition to Higher Structures}

Operator dynamics describes evolving persistence mechanisms using the
language of differential equations, semigroups, and smooth manifolds.

The next chapter raises the level of abstraction once again.

Rather than studying individual operators or their continuous
evolution, we study relationships between entire persistence systems.

Morphisms become objects of study.

Transformations between morphisms become fundamental.

Local constructions become global through sheaf-theoretic methods.

Universal properties organize admissible extensions through Kan
extensions and higher categorical limits.

This culminates in a higher-dimensional formulation of persistence in
which configurations, operators, operator evolutions, and
inter-system transformations appear within a unified categorical
framework.

\chapter{Higher Structures}

\epigraph{%
Persistence is local.
Understanding persistence is global.
}{}

\section{From Operators to Categories}

The previous chapters developed persistence through increasingly rich
mathematical structures.

Configurations became trajectories.

Trajectories became persistence operators.

Persistence operators evolved through continuous dynamics.

We now take one final abstraction step.

Rather than studying individual persistence operators, we study entire
collections of persistence systems together with the transformations
between them.

The appropriate language for this level of abstraction is higher
category theory.

Throughout this chapter,

\[
\mathbf{Pers}
\]

denotes the category of persistence systems introduced previously.

\section{Persistence 2-Categories}

Morphisms themselves may possess meaningful transformations.

This motivates the introduction of two-dimensional categorical
structure.

\begin{definition}[Persistence 2-Category]

The persistence 2-category

\[
\mathbf{Pers}_2
\]

consists of

\begin{itemize}

\item objects:
persistence systems,

\item 1-morphisms:
admissibility-preserving persistence operators,

\item 2-morphisms:
continuous deformations between persistence operators.

\end{itemize}

\end{definition}

The introduction of 2-morphisms allows operator evolution itself to be
treated as an object of mathematical study.

\section{Vertical and Horizontal Composition}

The two-dimensional structure admits two distinct notions of
composition.

\begin{definition}[Vertical Composition]

Let

\[
\alpha:
F
\Rightarrow
G,
\qquad
\beta:
G
\Rightarrow
H.
\]

Their vertical composition is

\[
\beta
\circ_v
\alpha:
F
\Rightarrow
H.
\]

\end{definition}

Vertical composition combines successive operator deformations.

\begin{definition}[Horizontal Composition]

Suppose

\[
F_1,F_2,G_1,G_2
\]

are persistence operators together with compatible
2-morphisms.

Their horizontal composition is

\[
\beta
\circ_h
\alpha.
\]

Horizontal composition combines independent persistence pipelines.

\end{definition}

The interchange law guarantees compatibility between the two forms of
composition.

\begin{theorem}[Interchange Law]

Let

\[
\alpha_1,\alpha_2,
\beta_1,\beta_2
\]

be composable 2-morphisms.

Then

\[
(\beta_2\circ_v\beta_1)
\circ_h
(\alpha_2\circ_v\alpha_1)
=
(\beta_2\circ_h\alpha_2)
\circ_v
(\beta_1\circ_h\alpha_1).
\]

\end{theorem}

\begin{proof}

The result follows from the axioms defining strict
2-categories.

\end{proof}

\section{Natural Transformations}

Observation itself behaves naturally with respect to persistence
operators.

\begin{definition}[Persistence Natural Transformation]

Let

\[
F,G:
\mathbf{Pers}
\rightarrow
\mathbf{Obs}
\]

be functors.

A persistence natural transformation

\[
\eta:
F
\Rightarrow
G
\]

consists of morphisms

\[
\eta_X:
F(X)
\rightarrow
G(X)
\]

such that every admissibility-preserving operator satisfies the natural
commutative square.

\end{definition}

Naturality expresses compatibility between observation and
continuation.

\section{Limits of Persistence Systems}

Complex persistence architectures are assembled from simpler
components.

Category theory captures this process through limits.

\begin{definition}[Persistence Limit]

A persistence limit is a universal object together with compatible
projection morphisms satisfying the usual universal property.

\end{definition}

Limits describe the integration of multiple persistence systems into a
single coherent architecture.

\begin{definition}[Persistence Colimit]

A persistence colimit is the universal object obtained by gluing
compatible persistence systems together.

\end{definition}

Where limits combine information through consistency,

colimits assemble distributed persistence into larger structures.

\section{Sheaves of Persistence}

Repair and continuation are fundamentally local operations.

Global persistence emerges by consistently combining local data.

\begin{definition}[Persistence Sheaf]

Let

\[
X
\]

be a topological space.

A persistence sheaf assigns to every open set

\[
U
\subseteq
X
\]

a persistence system

\[
\mathcal F(U)
\]

together with restriction morphisms satisfying the standard sheaf
axioms.

\end{definition}

Local repairs therefore become sections of a sheaf.

Global persistence corresponds to compatible gluing of these local
sections.

\begin{theorem}[Gluing Principle]

Suppose

\[
\{U_i\}
\]

forms an open cover of

\[
X.
\]

If persistence sections agree on every overlap

\[
U_i
\cap
U_j,
\]

then there exists a unique global persistence section.

\end{theorem}

\begin{proof}

The result follows directly from the sheaf gluing axiom.

\end{proof}

\section{Fibered Persistence}

Many persistence systems vary continuously over external parameters.

Examples include temporal evolution, spatial organization, and varying
environmental conditions.

\begin{definition}[Persistence Fibration]

A persistence fibration is a projection

\[
\pi:
E
\rightarrow
B
\]

whose fibers

\[
\pi^{-1}(b)
\]

are persistence systems.

\end{definition}

The base space indexes families of persistence structures while the
total space represents their collective organization.

\section{Kan Extensions}

Universal extension problems naturally arise whenever persistence
operators must be extended from local domains to larger admissible
contexts.

\begin{definition}[Left Persistence Kan Extension]

Let

\[
F:
\mathcal C
\rightarrow
\mathbf{Pers}
\]

and

\[
K:
\mathcal C
\rightarrow
\mathcal D.
\]

The left persistence Kan extension

\[
\operatorname{Lan}_K(F)
\]

is the universal extension of

\[
F
\]

along

\[
K.
\]

\end{definition}

Kan extensions provide the canonical method for extending local
persistence constructions while preserving universal properties.

\section{Distributed Persistence Architectures}

The categorical constructions developed above naturally assemble into
distributed persistence networks.

Individual repair systems become objects.

Communication channels become morphisms.

Adaptive modifications become 2-morphisms.

Local repair procedures become sheaf sections.

Global consistency follows from universal gluing.

The resulting architecture is independent of any particular physical
implementation.

It characterizes persistence entirely through structural relations.

\section{The Persistence Correspondence Principle}

The development of this monograph may be summarized by a single
structural correspondence.

\begin{theorem}[Persistence Correspondence Principle]

Every admissible persistence system admits equivalent descriptions at
multiple mathematical levels:

\[
\begin{aligned}
\text{Observation}
&\longleftrightarrow
\text{Projection},
\\
\text{Continuation}
&\longleftrightarrow
\text{Geometry},
\\
\text{Repair}
&\longleftrightarrow
\text{Optimization},
\\
\text{Operators}
&\longleftrightarrow
\text{Algebra},
\\
\text{Recursive Continuation}
&\longleftrightarrow
\text{Fixed-Point Theory},
\\
\text{Operator Dynamics}
&\longleftrightarrow
\text{Semigroup Theory},
\\
\text{Distributed Persistence}
&\longleftrightarrow
\text{Higher Category Theory}.
\end{aligned}
\]

Each description preserves the same admissibility structure while
revealing different mathematical aspects of persistence.

\end{theorem}

\begin{proof}

The preceding chapters establish functorial correspondences between
these mathematical formulations. Observation induces configuration
spaces; admissibility defines geometric constraints; repair yields
variational principles; operator composition forms algebraic
structures; recursive continuation invokes fixed-point theory; operator
evolution generates semigroups; and higher categorical constructions
organize these relationships through universal properties. Each level
extends the previous one without altering the underlying admissibility
relations.

\end{proof}

\chapter{Synthesis}

\epigraph{%
The mathematics is complete.
The interpretation remains open.
}{}

\section{The Architecture of Persistence}

The preceding chapters developed persistence from progressively richer
mathematical viewpoints.

Observation established the passage from high-dimensional field spaces
to reduced configuration spaces.

Invariant theory identified the structural quantities whose
preservation defines admissibility.

Continuation replaced isolated states with admissible trajectories.

Repair introduced a quantitative geometry measuring deviation from
those trajectories.

Operator calculus described admissible transformations.

Recursive continuation explained self-maintaining persistence.

Operator dynamics permitted the repair mechanisms themselves to evolve.

Finally, higher category theory organized entire persistence systems
into a unified mathematical framework.

Although each chapter introduced new mathematical structures, every
construction remained centered upon a single organizing principle:

\[
\text{Persistence}
=
\text{Admissible Continuation}.
\]

Everything else follows as mathematical consequence.

\section{The Layered Structure of the Theory}

The framework developed throughout this monograph may be summarized as
a hierarchy.

\[
\begin{array}{c}
\text{Higher Categories}
\\
\uparrow
\\
\text{Operator Dynamics}
\\
\uparrow
\\
\text{Recursive Continuation}
\\
\uparrow
\\
\text{Operator Calculus}
\\
\uparrow
\\
\text{Repair Geometry}
\\
\uparrow
\\
\text{Continuation}
\\
\uparrow
\\
\text{Invariants}
\\
\uparrow
\\
\text{Observation}
\end{array}
\]

Each level assumes the existence of the preceding structures while
introducing a new mathematical language.

No level replaces another.

Instead, every level refines the preceding description.

\section{The Principle of Structural Sufficiency}

Throughout the development, no appeal has been made to the semantic
meaning of individual configurations.

Only structural relationships have been used.

This motivates the following principle.

\begin{principle}[Structural Sufficiency]

Whenever admissibility, continuation, repair, and recursive
preservation are completely specified, the mathematical theory of
persistence is completely specified.

Additional semantic interpretation may enrich applications but is not
required for the internal mathematical development.

\end{principle}

Structural sufficiency explains why the same formalism applies equally
to physical systems, computational processes, biological regulation,
linguistic evolution, and abstract algebraic constructions.

\section{Universality}

The theory intentionally separates mathematical structure from
application.

A persistence system may describe

\begin{itemize}

\item a differential equation,

\item an algorithm,

\item a proof,

\item a biological organism,

\item a communication protocol,

\item an institutional process,

\item or any system admitting observation, admissibility,
continuation, and repair.

\end{itemize}

Consequently, persistence should be viewed as a mathematical framework
rather than a domain-specific model.

\section{The Persistence Ladder}

The entire theory may be summarized by the following sequence of
questions.

\[
\begin{array}{ll}
\textbf{Observation:}
&
\text{What is distinguishable?}
\\[1ex]
\textbf{Invariants:}
&
\text{What must remain unchanged?}
\\[1ex]
\textbf{Continuation:}
&
\text{Which trajectories are admissible?}
\\[1ex]
\textbf{Repair:}
&
\text{How is admissibility restored?}
\\[1ex]
\textbf{Operators:}
&
\text{How are repairs composed?}
\\[1ex]
\textbf{Recursion:}
&
\text{How is persistence maintained indefinitely?}
\\[1ex]
\textbf{Dynamics:}
&
\text{How do repair mechanisms evolve?}
\\[1ex]
\textbf{Higher Structures:}
&
\text{How do entire persistence systems relate?}
\end{array}
\]

Each question introduces exactly one new mathematical object while
retaining all previous constructions.

\section{The Persistence Theorem}

The principal mathematical result of this monograph may now be stated
informally.

\begin{theorem}[Persistence Theorem]

Suppose

\begin{enumerate}

\item observation induces a well-defined configuration space,

\item an invariant basis determines admissibility,

\item admissible continuations exist,

\item reconstruction defect is measurable,

\item repair operators preserve admissibility,

\item recursive continuation admits bounded defect,

\item operator evolution preserves admissibility,

\item higher categorical compatibility conditions hold.

\end{enumerate}

Then persistence is completely characterized by the admissible
continuation structure generated by these data.

\end{theorem}

\begin{proof}

Each chapter establishes one component of the theorem.

Observation constructs configuration space.

Invariant theory defines admissibility.

Continuation establishes viable trajectories.

Repair geometry minimizes reconstruction defect.

Operator calculus organizes admissible transformations.

Recursive continuation guarantees indefinite preservation.

Operator dynamics permits evolving repair mechanisms.

Higher categorical constructions organize distributed persistence
systems through universal properties.

Together these structures provide a complete mathematical description
of persistence within the axiomatic framework developed throughout the
book.

\end{proof}

\section{Limitations}

The theory deliberately establishes mathematical structure rather than
empirical claims.

Accordingly, it does not prove the existence of persistence mechanisms
within any particular physical, biological, computational, or social
system.

Instead, it provides a language in which such systems may be modeled,
compared, and analyzed whenever the stated axioms are satisfied.

Verification within any specific scientific domain remains an empirical
question beyond the scope of this monograph.

\section{Future Directions}

Several natural extensions remain open.

\begin{enumerate}

\item Stochastic persistence under probabilistic observation.

\item Infinite-dimensional persistence manifolds.

\item Homotopy-theoretic persistence invariants.

\item Derived persistence categories.

\item Persistence in noncommutative geometry.

\item Quantum persistence operators.

\item Persistence over enriched and $\infty$-categories.

\item Computational complexity of optimal repair.

\end{enumerate}

These directions require additional mathematical development beyond the
present axiomatic framework.

\section{Closing Remarks}

The objective of this work has been neither to introduce a new physical
theory nor to replace existing mathematical disciplines.

Instead, the goal has been to identify a common mathematical structure
underlying systems that persist through observation, constraint,
correction, and recursive continuation.

Whether that structure ultimately proves useful in geometry,
computation, biology, physics, linguistics, or other fields will depend
upon future mathematical development and empirical investigation.

The framework presented here is intended as a foundation upon which
such developments may be built.

\bigskip

\begin{center}

\emph{Observe.\\
Distinguish.\\
Constrain.\\
Continue.\\
Repair.\\
Persist.}

\end{center}


\newpage 
\section*{Appendices}

\appendix

\chapter{Metric and Topological Foundations}

This appendix collects the basic topological and metric results used
throughout the monograph. Unless otherwise stated, all spaces are
assumed to be Hausdorff.

%%%%%%%%%%%%%%%%%%%%%%%%%%%%%%%%%%%%%%%%%%%%%%%%%%%%%%%%%%%%%%%

\section{Metric Spaces}

\begin{definition}[Metric Space]

A metric space is an ordered pair

\[
(X,d)
\]

where

\[
d:X\times X\rightarrow\mathbb{R}_{\ge0}
\]

satisfies, for every
\(x,y,z\in X\),

\begin{enumerate}

\item
\[
d(x,y)=0
\iff
x=y,
\]

\item
\[
d(x,y)=d(y,x),
\]

\item
\[
d(x,z)
\le
d(x,y)+d(y,z).
\]

\end{enumerate}

\end{definition}

The function \(d\) is called the metric on \(X\).

%%%%%%%%%%%%%%%%%%%%%%%%%%%%%%%%%%%%%%%%%%%%%%%%%%%%%%%%%%%%%%%

\begin{definition}[Open Ball]

For

\[
x\in X,
\qquad
r>0,
\]

the open ball centered at \(x\) with radius \(r\) is

\[
B_r(x)
=
\{
y\in X:
d(x,y)<r
\}.
\]

\end{definition}

%%%%%%%%%%%%%%%%%%%%%%%%%%%%%%%%%%%%%%%%%%%%%%%%%%%%%%%%%%%%%%%

\begin{definition}[Closed Ball]

The closed ball of radius \(r\) is

\[
\overline B_r(x)
=
\{
y:
d(x,y)\le r
\}.
\]

\end{definition}

%%%%%%%%%%%%%%%%%%%%%%%%%%%%%%%%%%%%%%%%%%%%%%%%%%%%%%%%%%%%%%%

\begin{proposition}

Open balls form a basis for a topology on \(X\).

\end{proposition}

\begin{proof}

If

\[
y
\in
B_r(x)
\cap
B_s(z),
\]

let

\[
\varepsilon
=
\min
\{
r-d(x,y),
s-d(z,y)
\}.
\]

Then

\[
B_\varepsilon(y)
\subseteq
B_r(x)\cap B_s(z).
\]

Hence the basis axioms hold.

\end{proof}

%%%%%%%%%%%%%%%%%%%%%%%%%%%%%%%%%%%%%%%%%%%%%%%%%%%%%%%%%%%%%%%

\section{Convergence}

\begin{definition}

A sequence

\[
(x_n)
\]

converges to

\[
x
\]

if

\[
\forall\varepsilon>0,
\exists N
:
n\ge N
\Longrightarrow
d(x_n,x)<\varepsilon.
\]

We write

\[
x_n\to x.
\]

\end{definition}

%%%%%%%%%%%%%%%%%%%%%%%%%%%%%%%%%%%%%%%%%%%%%%%%%%%%%%%%%%%%%%%

\begin{definition}[Cauchy Sequence]

A sequence

\[
(x_n)
\]

is Cauchy if

\[
\forall\varepsilon>0,
\exists N
:
m,n\ge N
\Longrightarrow
d(x_n,x_m)<\varepsilon.
\]

\end{definition}

%%%%%%%%%%%%%%%%%%%%%%%%%%%%%%%%%%%%%%%%%%%%%%%%%%%%%%%%%%%%%%%

\begin{definition}[Complete Metric Space]

A metric space is complete if every Cauchy sequence converges.

\end{definition}

%%%%%%%%%%%%%%%%%%%%%%%%%%%%%%%%%%%%%%%%%%%%%%%%%%%%%%%%%%%%%%%

\begin{example}

The spaces

\[
\mathbb R^n,
\qquad
\ell^2,
\qquad
L^2(\Omega)
\]

are complete.

The open interval

\[
(0,1)
\]

is not complete under the Euclidean metric.

\end{example}

%%%%%%%%%%%%%%%%%%%%%%%%%%%%%%%%%%%%%%%%%%%%%%%%%%%%%%%%%%%%%%%

\section{Topological Spaces}

\begin{definition}

A topology on a set

\[
X
\]

is a collection

\[
\tau
\]

of subsets of \(X\) satisfying

\begin{enumerate}

\item
\[
\varnothing,
X
\in\tau,
\]

\item
arbitrary unions belong to \(\tau\),

\item
finite intersections belong to \(\tau\).

\end{enumerate}

The pair

\[
(X,\tau)
\]

is called a topological space.

\end{definition}

%%%%%%%%%%%%%%%%%%%%%%%%%%%%%%%%%%%%%%%%%%%%%%%%%%%%%%%%%%%%%%%

\begin{definition}

Members of

\[
\tau
\]

are called open sets.

Their complements are closed.

\end{definition}

%%%%%%%%%%%%%%%%%%%%%%%%%%%%%%%%%%%%%%%%%%%%%%%%%%%%%%%%%%%%%%%

\section{Continuous Maps}

\begin{definition}

A function

\[
f:
X
\rightarrow
Y
\]

between topological spaces is continuous if

\[
f^{-1}(U)
\]

is open for every open set

\[
U
\subseteq
Y.
\]

\end{definition}

%%%%%%%%%%%%%%%%%%%%%%%%%%%%%%%%%%%%%%%%%%%%%%%%%%%%%%%%%%%%%%%

\begin{proposition}

The composition of continuous maps is continuous.

\end{proposition}

\begin{proof}

If

\[
g:Y\to Z
\]

is continuous,

then

\[
(g\circ f)^{-1}(U)
=
f^{-1}(g^{-1}(U)).
\]

Both inverse images are open.

\end{proof}

%%%%%%%%%%%%%%%%%%%%%%%%%%%%%%%%%%%%%%%%%%%%%%%%%%%%%%%%%%%%%%%

\section{Compactness}

\begin{definition}

A topological space is compact if every open cover admits a finite
subcover.

\end{definition}

%%%%%%%%%%%%%%%%%%%%%%%%%%%%%%%%%%%%%%%%%%%%%%%%%%%%%%%%%%%%%%%

\begin{theorem}[Heine--Borel]

A subset of

\[
\mathbb R^n
\]

is compact if and only if it is closed and bounded.

\end{theorem}

%%%%%%%%%%%%%%%%%%%%%%%%%%%%%%%%%%%%%%%%%%%%%%%%%%%%%%%%%%%%%%%

\begin{theorem}

Every continuous image of a compact space is compact.

\end{theorem}

%%%%%%%%%%%%%%%%%%%%%%%%%%%%%%%%%%%%%%%%%%%%%%%%%%%%%%%%%%%%%%%

\begin{corollary}

A continuous real-valued function on a compact space attains both its
maximum and minimum.

\end{corollary}

%%%%%%%%%%%%%%%%%%%%%%%%%%%%%%%%%%%%%%%%%%%%%%%%%%%%%%%%%%%%%%%

\section{Connectedness}

\begin{definition}

A space is connected if it cannot be expressed as the union of two
nonempty disjoint open subsets.

\end{definition}

%%%%%%%%%%%%%%%%%%%%%%%%%%%%%%%%%%%%%%%%%%%%%%%%%%%%%%%%%%%%%%%

\begin{proposition}

Continuous images of connected spaces remain connected.

\end{proposition}

%%%%%%%%%%%%%%%%%%%%%%%%%%%%%%%%%%%%%%%%%%%%%%%%%%%%%%%%%%%%%%%

\section{Hausdorff Spaces}

\begin{definition}

A topological space is Hausdorff if every pair of distinct points
possesses disjoint neighborhoods.

\end{definition}

%%%%%%%%%%%%%%%%%%%%%%%%%%%%%%%%%%%%%%%%%%%%%%%%%%%%%%%%%%%%%%%

\begin{theorem}

Compact subsets of Hausdorff spaces are closed.

\end{theorem}

%%%%%%%%%%%%%%%%%%%%%%%%%%%%%%%%%%%%%%%%%%%%%%%%%%%%%%%%%%%%%%%

\section{Product Topology}

\begin{definition}

Let

\[
(X_i,\tau_i),
\qquad
i\in I,
\]

be topological spaces.

The product topology on

\[
\prod_i X_i
\]

is generated by sets of the form

\[
\prod_i U_i,
\]

where each

\[
U_i
\]

is open and all but finitely many satisfy

\[
U_i=X_i.
\]

\end{definition}

%%%%%%%%%%%%%%%%%%%%%%%%%%%%%%%%%%%%%%%%%%%%%%%%%%%%%%%%%%%%%%%

\section{Quotient Spaces}

\begin{definition}

Let

\[
\sim
\]

be an equivalence relation on

\[
X.
\]

The quotient topology on

\[
X/{\sim}
\]

is the finest topology making the projection

\[
\pi:X\to X/{\sim}
\]

continuous.

\end{definition}

%%%%%%%%%%%%%%%%%%%%%%%%%%%%%%%%%%%%%%%%%%%%%%%%%%%%%%%%%%%%%%%

\section{Topological Manifolds}

\begin{definition}

An \(n\)-dimensional topological manifold is a Hausdorff,
second-countable space locally homeomorphic to

\[
\mathbb R^n.
\]

\end{definition}

%%%%%%%%%%%%%%%%%%%%%%%%%%%%%%%%%%%%%%%%%%%%%%%%%%%%%%%%%%%%%%%

\begin{remark}

Beginning in Chapter~1, configuration spaces are modeled as
topological manifolds or Banach manifolds whenever differentiable
structure is required.

Compactness, continuity, connectedness, and completeness provide the
topological foundation upon which admissibility, reconstruction, and
continuation are later developed.

\end{remark}

\chapter{Analysis and Variational Methods}

This appendix collects the principal analytical results underlying the
variational, optimization, and stability arguments used throughout the
monograph. Unless otherwise stated, all vector spaces are assumed to be
real Banach spaces.

%%%%%%%%%%%%%%%%%%%%%%%%%%%%%%%%%%%%%%%%%%%%%%%%%%%%%%%%%%%%%%%

\section{Lower Semicontinuity}

\begin{definition}[Lower Semicontinuity]

Let

\[
X
\]

be a topological space.

A function

\[
f:X\rightarrow(-\infty,\infty]
\]

is lower semicontinuous if

\[
f(x)
\le
\liminf_{n\rightarrow\infty}
f(x_n)
\]

whenever

\[
x_n\rightarrow x.
\]

\end{definition}

%%%%%%%%%%%%%%%%%%%%%%%%%%%%%%%%%%%%%%%%%%%%%%%%%%%%%%%%%%%%%%%

\begin{proposition}

The sublevel sets

\[
\{x:f(x)\le c\}
\]

of a lower semicontinuous function are closed.

\end{proposition}

\begin{proof}

Suppose

\[
x_n\rightarrow x,
\qquad
f(x_n)\le c.
\]

Then

\[
f(x)
\le
\liminf f(x_n)
\le c.
\]

Hence

\[
x
\]

belongs to the sublevel set.

\end{proof}

%%%%%%%%%%%%%%%%%%%%%%%%%%%%%%%%%%%%%%%%%%%%%%%%%%%%%%%%%%%%%%%

\section{Proper Functionals}

\begin{definition}

A functional

\[
f:X\rightarrow(-\infty,\infty]
\]

is proper if

\begin{enumerate}

\item
\[
f(x)>-\infty
\]

for every

\[
x,
\]

\item

there exists at least one point satisfying

\[
f(x)<\infty.
\]

\end{enumerate}

\end{definition}

%%%%%%%%%%%%%%%%%%%%%%%%%%%%%%%%%%%%%%%%%%%%%%%%%%%%%%%%%%%%%%%

\section{Coercivity}

\begin{definition}

A functional

\[
f:X\rightarrow\mathbb R
\]

is coercive if

\[
\|x\|\rightarrow\infty
\Longrightarrow
f(x)\rightarrow\infty.
\]

\end{definition}

%%%%%%%%%%%%%%%%%%%%%%%%%%%%%%%%%%%%%%%%%%%%%%%%%%%%%%%%%%%%%%%

\begin{proposition}

Every minimizing sequence of a coercive functional is bounded.

\end{proposition}

\begin{proof}

If

\[
\|x_n\|\rightarrow\infty,
\]

then coercivity implies

\[
f(x_n)\rightarrow\infty,
\]

contradicting minimality.

\end{proof}

%%%%%%%%%%%%%%%%%%%%%%%%%%%%%%%%%%%%%%%%%%%%%%%%%%%%%%%%%%%%%%%

\section{Existence of Minimizers}

\begin{theorem}[Direct Method]

Let

\[
X
\]

be reflexive.

Suppose

\[
f:X\rightarrow(-\infty,\infty]
\]

is

\begin{enumerate}

\item proper,

\item coercive,

\item lower semicontinuous.

\end{enumerate}

Then

\[
f
\]

attains its minimum.

\end{theorem}

\begin{proof}

Choose a minimizing sequence.

Coercivity gives boundedness.

Reflexivity yields a weakly convergent subsequence.

Lower semicontinuity preserves the minimum under passage to the limit.

\end{proof}

%%%%%%%%%%%%%%%%%%%%%%%%%%%%%%%%%%%%%%%%%%%%%%%%%%%%%%%%%%%%%%%

\section{Convexity}

\begin{definition}

A subset

\[
C
\]

of a vector space is convex if

\[
tx+(1-t)y
\in C
\]

for every

\[
0\le t\le1.
\]

\end{definition}

%%%%%%%%%%%%%%%%%%%%%%%%%%%%%%%%%%%%%%%%%%%%%%%%%%%%%%%%%%%%%%%

\begin{definition}

A functional

\[
f
\]

is convex if

\[
f(tx+(1-t)y)
\le
tf(x)+(1-t)f(y).
\]

\end{definition}

%%%%%%%%%%%%%%%%%%%%%%%%%%%%%%%%%%%%%%%%%%%%%%%%%%%%%%%%%%%%%%%

\begin{theorem}

Every local minimum of a convex functional is global.

\end{theorem}

%%%%%%%%%%%%%%%%%%%%%%%%%%%%%%%%%%%%%%%%%%%%%%%%%%%%%%%%%%%%%%%

\section{Fréchet Derivatives}

\begin{definition}

Let

\[
X,Y
\]

be Banach spaces.

A map

\[
F:X\rightarrow Y
\]

is Fréchet differentiable at

\[
x
\]

if there exists a bounded linear operator

\[
DF(x)
\]

such that

\[
F(x+h)
=
F(x)
+
DF(x)h
+
o(\|h\|).
\]

\end{definition}

%%%%%%%%%%%%%%%%%%%%%%%%%%%%%%%%%%%%%%%%%%%%%%%%%%%%%%%%%%%%%%%

\begin{proposition}[Chain Rule]

If

\[
F:X\rightarrow Y,
\qquad
G:Y\rightarrow Z
\]

are Fréchet differentiable,

then

\[
D(G\circ F)
=
DG(F(x))
\circ
DF(x).
\]

\end{proposition}

%%%%%%%%%%%%%%%%%%%%%%%%%%%%%%%%%%%%%%%%%%%%%%%%%%%%%%%%%%%%%%%

\section{Gradient Flows}

\begin{definition}

Given a differentiable functional

\[
E:X\rightarrow\mathbb R,
\]

its gradient flow satisfies

\[
\dot x
=
-\nabla E(x).
\]

\end{definition}

%%%%%%%%%%%%%%%%%%%%%%%%%%%%%%%%%%%%%%%%%%%%%%%%%%%%%%%%%%%%%%%

\begin{proposition}[Energy Dissipation]

Along every gradient flow,

\[
\frac{d}{dt}
E(x(t))
=
-
\|\nabla E(x(t))\|^2.
\]

\end{proposition}

\begin{proof}

Apply the chain rule:

\[
\frac{dE}{dt}
=
\langle
\nabla E,
\dot x
\rangle.
\]

Substitute

\[
\dot x
=
-\nabla E.
\]

\end{proof}

%%%%%%%%%%%%%%%%%%%%%%%%%%%%%%%%%%%%%%%%%%%%%%%%%%%%%%%%%%%%%%%

\section{Arzelà--Ascoli}

\begin{definition}

A family of functions is equicontinuous if

for every

\[
\varepsilon>0
\]

there exists

\[
\delta>0
\]

such that

\[
d(x,y)<\delta
\]

implies

\[
d(f(x),f(y))
<
\varepsilon
\]

for every function in the family.

\end{definition}

%%%%%%%%%%%%%%%%%%%%%%%%%%%%%%%%%%%%%%%%%%%%%%%%%%%%%%%%%%%%%%%

\begin{theorem}[Arzelà--Ascoli]

Let

\[
K
\]

be compact.

A uniformly bounded equicontinuous family

\[
\mathcal F
\subseteq
C(K)
\]

has compact closure in the uniform topology.

\end{theorem}

%%%%%%%%%%%%%%%%%%%%%%%%%%%%%%%%%%%%%%%%%%%%%%%%%%%%%%%%%%%%%%%

\section{Grönwall's Inequality}

\begin{theorem}

Suppose

\[
u(t)
\]

satisfies

\[
u(t)
\le
a
+
b
\int_0^t
u(s)\,ds.
\]

Then

\[
u(t)
\le
ae^{bt}.
\]

\end{theorem}

%%%%%%%%%%%%%%%%%%%%%%%%%%%%%%%%%%%%%%%%%%%%%%%%%%%%%%%%%%%%%%%

\section{Barbalat's Lemma}

\begin{theorem}

Suppose

\[
f
\]

is uniformly continuous on

\[
[0,\infty)
\]

and

\[
\int_0^\infty
f(t)\,dt
<
\infty.
\]

Then

\[
f(t)
\rightarrow
0.
\]

\end{theorem}

%%%%%%%%%%%%%%%%%%%%%%%%%%%%%%%%%%%%%%%%%%%%%%%%%%%%%%%%%%%%%%%

\section{Lyapunov Functionals}

\begin{definition}

A Lyapunov functional is a nonnegative function

\[
V:X\rightarrow\mathbb R_{\ge0}
\]

satisfying

\[
\frac{d}{dt}
V(x(t))
\le0
\]

along trajectories.

\end{definition}

%%%%%%%%%%%%%%%%%%%%%%%%%%%%%%%%%%%%%%%%%%%%%%%%%%%%%%%%%%%%%%%

\begin{theorem}

If

\[
V
\]

is strictly decreasing outside an equilibrium,

then every bounded trajectory approaches the invariant set where

\[
\dot V=0.
\]

\end{theorem}

%%%%%%%%%%%%%%%%%%%%%%%%%%%%%%%%%%%%%%%%%%%%%%%%%%%%%%%%%%%%%%%

\begin{remark}

Chapter~2 employs lower semicontinuity, coercivity, Fréchet
differentiability, and gradient-flow dissipation to define repair
functionals. Chapters~3--7 repeatedly invoke Grönwall's inequality,
Arzelà--Ascoli compactness, and Lyapunov methods when proving
continuation, convergence, and stability results.

\end{remark}

\chapter{Fixed Point Theory}

This appendix summarizes the principal fixed-point theorems employed
throughout the monograph. Fixed-point methods provide the analytical
foundation for recursive continuation, repair convergence, and operator
stability.

%%%%%%%%%%%%%%%%%%%%%%%%%%%%%%%%%%%%%%%%%%%%%%%%%%%%%%%%%%%%%%%

\section{Contractions}

\begin{definition}[Contraction]

Let

\[
(X,d)
\]

be a metric space.

A mapping

\[
T:X\rightarrow X
\]

is called a contraction if there exists a constant

\[
0\le k<1
\]

such that

\[
d(Tx,Ty)
\le
k\,d(x,y)
\]

for all

\[
x,y\in X.
\]

\end{definition}

%%%%%%%%%%%%%%%%%%%%%%%%%%%%%%%%%%%%%%%%%%%%%%%%%%%%%%%%%%%%%%%

\begin{lemma}

Every iterate of a contraction satisfies

\[
d(T^n x,T^n y)
\le
k^n d(x,y).
\]

\end{lemma}

\begin{proof}

Apply the contraction inequality inductively.

\end{proof}

%%%%%%%%%%%%%%%%%%%%%%%%%%%%%%%%%%%%%%%%%%%%%%%%%%%%%%%%%%%%%%%

\section{Banach Contraction Principle}

\begin{theorem}[Banach]

Let

\[
(X,d)
\]

be complete.

Every contraction

\[
T:X\rightarrow X
\]

possesses a unique fixed point

\[
x^*
\]

satisfying

\[
T(x^*)=x^*.
\]

Furthermore,

\[
x_{n+1}=Tx_n
\]

converges to

\[
x^*
\]

for every initial point.

\end{theorem}

\begin{proof}

The Picard iteration forms a Cauchy sequence because

\[
d(x_{n+1},x_n)
\le
k^n
d(x_1,x_0).
\]

Completeness yields convergence.

Uniqueness follows immediately from the contraction property.

\end{proof}

%%%%%%%%%%%%%%%%%%%%%%%%%%%%%%%%%%%%%%%%%%%%%%%%%%%%%%%%%%%%%%%

\section{Brouwer Fixed Point Theorem}

\begin{theorem}[Brouwer]

Every continuous mapping

\[
f:B^n\rightarrow B^n
\]

from the closed unit ball into itself possesses at least one fixed
point.

\end{theorem}

\begin{remark}

Brouwer's theorem guarantees existence but not uniqueness.

\end{remark}

%%%%%%%%%%%%%%%%%%%%%%%%%%%%%%%%%%%%%%%%%%%%%%%%%%%%%%%%%%%%%%%

\section{Schauder Fixed Point Theorem}

\begin{definition}

A mapping

\[
T:C\rightarrow C
\]

is compact if it maps bounded subsets into relatively compact subsets.

\end{definition}

%%%%%%%%%%%%%%%%%%%%%%%%%%%%%%%%%%%%%%%%%%%%%%%%%%%%%%%%%%%%%%%

\begin{theorem}[Schauder]

Let

\[
C
\]

be a nonempty compact convex subset of a Banach space.

Every continuous mapping

\[
T:C\rightarrow C
\]

admits a fixed point.

\end{theorem}

%%%%%%%%%%%%%%%%%%%%%%%%%%%%%%%%%%%%%%%%%%%%%%%%%%%%%%%%%%%%%%%

\section{Kakutani Fixed Point Theorem}

\begin{definition}

A set-valued mapping

\[
F:X
\rightrightarrows
X
\]

assigns to each point a subset of

\[
X.
\]

\end{definition}

%%%%%%%%%%%%%%%%%%%%%%%%%%%%%%%%%%%%%%%%%%%%%%%%%%%%%%%%%%%%%%%

\begin{theorem}[Kakutani]

Let

\[
C
\]

be compact and convex.

Suppose

\[
F:C
\rightrightarrows
C
\]

has nonempty compact convex values and closed graph.

Then

\[
F
\]

possesses a fixed point.

\end{theorem}

%%%%%%%%%%%%%%%%%%%%%%%%%%%%%%%%%%%%%%%%%%%%%%%%%%%%%%%%%%%%%%%

\section{Monotone Operators}

\begin{definition}

Let

\[
(P,\le)
\]

be a partially ordered set.

A mapping

\[
T:P\rightarrow P
\]

is monotone if

\[
x\le y
\Longrightarrow
Tx\le Ty.
\]

\end{definition}

%%%%%%%%%%%%%%%%%%%%%%%%%%%%%%%%%%%%%%%%%%%%%%%%%%%%%%%%%%%%%%%

\section{Complete Lattices}

\begin{definition}

A complete lattice is a partially ordered set in which every subset
admits both a supremum and an infimum.

\end{definition}

%%%%%%%%%%%%%%%%%%%%%%%%%%%%%%%%%%%%%%%%%%%%%%%%%%%%%%%%%%%%%%%

\begin{example}

The power set

\[
\mathcal P(X)
\]

ordered by inclusion is a complete lattice.

\end{example}

%%%%%%%%%%%%%%%%%%%%%%%%%%%%%%%%%%%%%%%%%%%%%%%%%%%%%%%%%%%%%%%

\section{Knaster--Tarski Fixed Point Theorem}

\begin{theorem}[Knaster--Tarski]

Let

\[
L
\]

be a complete lattice.

Every monotone operator

\[
T:L\rightarrow L
\]

possesses both a least fixed point and a greatest fixed point.

\end{theorem}

\begin{proof}

The least fixed point is

\[
\inf
\{
x:
Tx\le x
\},
\]

while the greatest fixed point is

\[
\sup
\{
x:
x\le Tx
\}.
\]

Monotonicity guarantees both satisfy

\[
Tx=x.
\]

\end{proof}

%%%%%%%%%%%%%%%%%%%%%%%%%%%%%%%%%%%%%%%%%%%%%%%%%%%%%%%%%%%%%%%

\section{Fixed Point Stability}

\begin{definition}

A fixed point

\[
x^*
\]

is locally stable if every sufficiently small perturbation generates a
trajectory remaining arbitrarily close to

\[
x^*.
\]

\end{definition}

%%%%%%%%%%%%%%%%%%%%%%%%%%%%%%%%%%%%%%%%%%%%%%%%%%%%%%%%%%%%%%%

\begin{proposition}

Every Banach fixed point is locally exponentially stable.

\end{proposition}

\begin{proof}

Since

\[
d(T^n x,x^*)
\le
k^n
d(x,x^*),
\]

the convergence rate is geometric.

\end{proof}

%%%%%%%%%%%%%%%%%%%%%%%%%%%%%%%%%%%%%%%%%%%%%%%%%%%%%%%%%%%%%%%

\section{Iterative Approximation}

\begin{definition}

Given

\[
T:X\rightarrow X,
\]

the Picard iteration is

\[
x_{n+1}=Tx_n.
\]

\end{definition}

%%%%%%%%%%%%%%%%%%%%%%%%%%%%%%%%%%%%%%%%%%%%%%%%%%%%%%%%%%%%%%%

\begin{proposition}

If

\[
T
\]

is a contraction, every Picard sequence converges to the unique fixed
point.

\end{proposition}

%%%%%%%%%%%%%%%%%%%%%%%%%%%%%%%%%%%%%%%%%%%%%%%%%%%%%%%%%%%%%%%

\begin{remark}

Chapter~6 employs the Banach Contraction Principle for recursive
continuation under contraction hypotheses, while the
Knaster--Tarski theorem provides fixed-point existence for monotone
admissibility operators. Schauder and Kakutani serve as more general
existence results when compactness replaces contractivity.

\end{remark}

\chapter{Operator Theory}

This appendix develops the operator-theoretic foundations used
throughout the monograph. Persistence operators are modeled as
endomorphisms on Banach spaces or smooth manifolds, and their
composition, continuity, and evolution are governed by the results
collected here.

%%%%%%%%%%%%%%%%%%%%%%%%%%%%%%%%%%%%%%%%%%%%%%%%%%%%%%%%%%%%%%%

\section{Linear Operators}

\begin{definition}

Let

\[
X,Y
\]

be vector spaces.

A mapping

\[
T:X\rightarrow Y
\]

is linear if

\[
T(ax+by)
=
aT(x)+bT(y)
\]

for every

\[
a,b\in\mathbb R,
\qquad
x,y\in X.
\]

\end{definition}

%%%%%%%%%%%%%%%%%%%%%%%%%%%%%%%%%%%%%%%%%%%%%%%%%%%%%%%%%%%%%%%

\section{Bounded Operators}

\begin{definition}

A linear operator

\[
T:X\rightarrow Y
\]

between normed spaces is bounded if there exists

\[
M>0
\]

such that

\[
\|Tx\|
\le
M\|x\|
\]

for every

\[
x\in X.
\]

\end{definition}

%%%%%%%%%%%%%%%%%%%%%%%%%%%%%%%%%%%%%%%%%%%%%%%%%%%%%%%%%%%%%%%

\begin{theorem}

A linear operator between normed spaces is continuous if and only if it
is bounded.

\end{theorem}

%%%%%%%%%%%%%%%%%%%%%%%%%%%%%%%%%%%%%%%%%%%%%%%%%%%%%%%%%%%%%%%

\section{Operator Norm}

\begin{definition}

The operator norm of

\[
T
\]

is

\[
\|T\|
=
\sup_{\|x\|=1}
\|Tx\|.
\]

\end{definition}

%%%%%%%%%%%%%%%%%%%%%%%%%%%%%%%%%%%%%%%%%%%%%%%%%%%%%%%%%%%%%%%

\begin{proposition}

The operator norm satisfies

\[
\|ST\|
\le
\|S\|
\,
\|T\|
\]

for every composable pair of bounded operators.

\end{proposition}

\begin{proof}

For every unit vector,

\[
\|STx\|
\le
\|S\|
\,
\|Tx\|
\le
\|S\|
\,
\|T\|.
\]

Taking the supremum completes the proof.

\end{proof}

%%%%%%%%%%%%%%%%%%%%%%%%%%%%%%%%%%%%%%%%%%%%%%%%%%%%%%%%%%%%%%%

\section{Endomorphisms}

\begin{definition}

An endomorphism is an operator

\[
T:X\rightarrow X.
\]

\end{definition}

%%%%%%%%%%%%%%%%%%%%%%%%%%%%%%%%%%%%%%%%%%%%%%%%%%%%%%%%%%%%%%%

\begin{definition}

The identity endomorphism is

\[
I(x)=x.
\]

\end{definition}

%%%%%%%%%%%%%%%%%%%%%%%%%%%%%%%%%%%%%%%%%%%%%%%%%%%%%%%%%%%%%%%

\section{Composition}

\begin{proposition}

Composition of endomorphisms is associative:

\[
(RS)T
=
R(ST).
\]

\end{proposition}

%%%%%%%%%%%%%%%%%%%%%%%%%%%%%%%%%%%%%%%%%%%%%%%%%%%%%%%%%%%%%%%

\begin{proposition}

The identity operator satisfies

\[
IT
=
TI
=
T.
\]

\end{proposition}

%%%%%%%%%%%%%%%%%%%%%%%%%%%%%%%%%%%%%%%%%%%%%%%%%%%%%%%%%%%%%%%

\section{Operator Monoids}

\begin{definition}

The collection

\[
\operatorname{End}(X)
\]

of all endomorphisms on

\[
X
\]

forms a monoid under composition.

\end{definition}

%%%%%%%%%%%%%%%%%%%%%%%%%%%%%%%%%%%%%%%%%%%%%%%%%%%%%%%%%%%%%%%

\begin{proof}

Composition is associative and possesses the identity operator.

\end{proof}

%%%%%%%%%%%%%%%%%%%%%%%%%%%%%%%%%%%%%%%%%%%%%%%%%%%%%%%%%%%%%%%

\section{Invertible Operators}

\begin{definition}

An operator

\[
T
\]

is invertible if there exists

\[
T^{-1}
\]

such that

\[
TT^{-1}
=
T^{-1}T
=
I.
\]

\end{definition}

%%%%%%%%%%%%%%%%%%%%%%%%%%%%%%%%%%%%%%%%%%%%%%%%%%%%%%%%%%%%%%%

\begin{proposition}

The invertible bounded operators form a group under composition.

\end{proposition}

%%%%%%%%%%%%%%%%%%%%%%%%%%%%%%%%%%%%%%%%%%%%%%%%%%%%%%%%%%%%%%%

\section{Closed Operators}

\begin{definition}

An operator

\[
T
\]

is closed if its graph

\[
\{
(x,Tx)
\}
\]

is a closed subset of

\[
X\times Y.
\]

\end{definition}

%%%%%%%%%%%%%%%%%%%%%%%%%%%%%%%%%%%%%%%%%%%%%%%%%%%%%%%%%%%%%%%

\section{Strong Operator Convergence}

\begin{definition}

A sequence

\[
T_n
\]

converges strongly to

\[
T
\]

if

\[
T_nx
\rightarrow
Tx
\]

for every

\[
x\in X.
\]

\end{definition}

%%%%%%%%%%%%%%%%%%%%%%%%%%%%%%%%%%%%%%%%%%%%%%%%%%%%%%%%%%%%%%%

\section{Weak Operator Convergence}

\begin{definition}

The sequence

\[
T_n
\]

converges weakly if

\[
\langle
\phi,
T_nx
\rangle
\rightarrow
\langle
\phi,
Tx
\rangle
\]

for every

\[
x\in X,
\qquad
\phi\in X^*.
\]

\end{definition}

%%%%%%%%%%%%%%%%%%%%%%%%%%%%%%%%%%%%%%%%%%%%%%%%%%%%%%%%%%%%%%%

\section{One-Parameter Semigroups}

\begin{definition}

A family

\[
(T_t)_{t\ge0}
\]

of bounded operators is a semigroup if

\[
T_0=I
\]

and

\[
T_{s+t}
=
T_sT_t
\]

for every

\[
s,t\ge0.
\]

\end{definition}

%%%%%%%%%%%%%%%%%%%%%%%%%%%%%%%%%%%%%%%%%%%%%%%%%%%%%%%%%%%%%%%

\begin{definition}

The semigroup is strongly continuous if

\[
\lim_{t\downarrow0}
T_tx
=
x
\]

for every

\[
x\in X.
\]

\end{definition}

%%%%%%%%%%%%%%%%%%%%%%%%%%%%%%%%%%%%%%%%%%%%%%%%%%%%%%%%%%%%%%%

\section{Infinitesimal Generators}

\begin{definition}

The infinitesimal generator of a strongly continuous semigroup is

\[
A
=
\lim_{h\downarrow0}
\frac{T_h-I}{h},
\]

whenever the limit exists.

\end{definition}

%%%%%%%%%%%%%%%%%%%%%%%%%%%%%%%%%%%%%%%%%%%%%%%%%%%%%%%%%%%%%%%

\begin{proposition}

If

\[
A
\]

generates

\[
(T_t),
\]

then

\[
\frac{d}{dt}
T_t
=
AT_t
=
T_tA
\]

on the domain of

\[
A.
\]

\end{proposition}

%%%%%%%%%%%%%%%%%%%%%%%%%%%%%%%%%%%%%%%%%%%%%%%%%%%%%%%%%%%%%%%

\section{Exponentials of Operators}

\begin{definition}

For a bounded operator

\[
A,
\]

the exponential is defined by

\[
e^A
=
\sum_{n=0}^{\infty}
\frac{A^n}{n!}.
\]

\end{definition}

%%%%%%%%%%%%%%%%%%%%%%%%%%%%%%%%%%%%%%%%%%%%%%%%%%%%%%%%%%%%%%%

\begin{proposition}

The operator exponential converges absolutely in operator norm.

\end{proposition}

%%%%%%%%%%%%%%%%%%%%%%%%%%%%%%%%%%%%%%%%%%%%%%%%%%%%%%%%%%%%%%%

\section{Commuting Operators}

\begin{definition}

Two operators commute if

\[
AB
=
BA.
\]

\end{definition}

%%%%%%%%%%%%%%%%%%%%%%%%%%%%%%%%%%%%%%%%%%%%%%%%%%%%%%%%%%%%%%%

\begin{proposition}

If

\[
AB=BA,
\]

then

\[
e^Ae^B
=
e^{A+B}.
\]

\end{proposition}

%%%%%%%%%%%%%%%%%%%%%%%%%%%%%%%%%%%%%%%%%%%%%%%%%%%%%%%%%%%%%%%

\section{Spectral Radius}

\begin{definition}

The spectral radius of

\[
T
\]

is

\[
r(T)
=
\sup
\{
|\lambda|
:
\lambda
\in
\sigma(T)
\},
\]

where

\[
\sigma(T)
\]

denotes the spectrum of

\[
T.
\]

\end{definition}

%%%%%%%%%%%%%%%%%%%%%%%%%%%%%%%%%%%%%%%%%%%%%%%%%%%%%%%%%%%%%%%

\begin{theorem}[Spectral Radius Formula]

For every bounded operator,

\[
r(T)
=
\lim_{n\rightarrow\infty}
\|T^n\|^{1/n}.
\]

\end{theorem}

%%%%%%%%%%%%%%%%%%%%%%%%%%%%%%%%%%%%%%%%%%%%%%%%%%%%%%%%%%%%%%%

\section{Invariant Subspaces}

\begin{definition}

A subspace

\[
M\subseteq X
\]

is invariant under

\[
T
\]

if

\[
T(M)
\subseteq
M.
\]

\end{definition}

%%%%%%%%%%%%%%%%%%%%%%%%%%%%%%%%%%%%%%%%%%%%%%%%%%%%%%%%%%%%%%%

\begin{proposition}

The restriction of

\[
T
\]

to an invariant subspace is itself an endomorphism.

\end{proposition}

%%%%%%%%%%%%%%%%%%%%%%%%%%%%%%%%%%%%%%%%%%%%%%%%%%%%%%%%%%%%%%%

\begin{remark}

Chapters~4 and~7 model persistence operators as endomorphisms acting on
configuration spaces. Primitive operators form operator monoids under
composition, repair flows evolve as strongly continuous semigroups, and
adaptive persistence dynamics are generated by infinitesimal operators.
The results collected here provide the standard operator-theoretic
machinery underlying those constructions.

\end{remark}

\chapter{Rewriting Systems}

This appendix summarizes the theory of abstract rewriting systems used
to establish normalization, confluence, and canonical operator
factorizations.

%%%%%%%%%%%%%%%%%%%%%%%%%%%%%%%%%%%%%%%%%%%%%%%%%%%%%%%%%%%%%%%

\section{Abstract Reduction Systems}

\begin{definition}

An abstract reduction system is a pair

\[
(X,\rightarrow),
\]

where

\[
X
\]

is a set and

\[
\rightarrow
\subseteq
X\times X
\]

is a binary relation called the reduction relation.

\end{definition}

%%%%%%%%%%%%%%%%%%%%%%%%%%%%%%%%%%%%%%%%%%%%%%%%%%%%%%%%%%%%%%%

\begin{definition}

If

\[
x\rightarrow y,
\]

we say that

\[
y
\]

is a one-step reduction of

\[
x.
\]

\end{definition}

%%%%%%%%%%%%%%%%%%%%%%%%%%%%%%%%%%%%%%%%%%%%%%%%%%%%%%%%%%%%%%%

\section{Reduction Sequences}

\begin{definition}

A reduction sequence is a sequence

\[
x_0
\rightarrow
x_1
\rightarrow
x_2
\rightarrow
\cdots
\]

generated by repeated applications of the reduction relation.

\end{definition}

%%%%%%%%%%%%%%%%%%%%%%%%%%%%%%%%%%%%%%%%%%%%%%%%%%%%%%%%%%%%%%%

\begin{definition}

The reflexive-transitive closure of

\[
\rightarrow
\]

is denoted

\[
\rightarrow^{*}.
\]

\end{definition}

%%%%%%%%%%%%%%%%%%%%%%%%%%%%%%%%%%%%%%%%%%%%%%%%%%%%%%%%%%%%%%%

\section{Normal Forms}

\begin{definition}

An element

\[
x
\]

is in normal form if no reduction applies:

\[
\nexists y
\text{ such that }
x\rightarrow y.
\]

\end{definition}

%%%%%%%%%%%%%%%%%%%%%%%%%%%%%%%%%%%%%%%%%%%%%%%%%%%%%%%%%%%%%%%

\begin{definition}

A reduction system is normalizing if every element admits a finite
reduction to a normal form.

\end{definition}

%%%%%%%%%%%%%%%%%%%%%%%%%%%%%%%%%%%%%%%%%%%%%%%%%%%%%%%%%%%%%%%

\section{Termination}

\begin{definition}

A reduction system is terminating if it contains no infinite reduction
sequence.

\end{definition}

%%%%%%%%%%%%%%%%%%%%%%%%%%%%%%%%%%%%%%%%%%%%%%%%%%%%%%%%%%%%%%%

\begin{proposition}

Every terminating reduction system is normalizing.

\end{proposition}

\begin{proof}

Infinite descent is impossible.

Therefore every maximal reduction sequence must terminate.

\end{proof}

%%%%%%%%%%%%%%%%%%%%%%%%%%%%%%%%%%%%%%%%%%%%%%%%%%%%%%%%%%%%%%%

\section{Confluence}

\begin{definition}

A reduction system is confluent if whenever

\[
x
\rightarrow^{*}
y,
\qquad
x
\rightarrow^{*}
z,
\]

there exists

\[
w
\]

such that

\[
y
\rightarrow^{*}
w,
\qquad
z
\rightarrow^{*}
w.
\]

\end{definition}

%%%%%%%%%%%%%%%%%%%%%%%%%%%%%%%%%%%%%%%%%%%%%%%%%%%%%%%%%%%%%%%

\section{Local Confluence}

\begin{definition}

The reduction relation is locally confluent if

\[
x
\rightarrow
y,
\qquad
x
\rightarrow
z
\]

implies the existence of

\[
w
\]

satisfying

\[
y
\rightarrow^{*}
w,
\qquad
z
\rightarrow^{*}
w.
\]

\end{definition}

%%%%%%%%%%%%%%%%%%%%%%%%%%%%%%%%%%%%%%%%%%%%%%%%%%%%%%%%%%%%%%%

\section{Church--Rosser Property}

\begin{definition}

A rewriting system satisfies the Church--Rosser property if

\[
x
\leftrightarrow^{*}
y
\]

implies the existence of

\[
z
\]

such that

\[
x
\rightarrow^{*}
z,
\qquad
y
\rightarrow^{*}
z.
\]

\end{definition}

%%%%%%%%%%%%%%%%%%%%%%%%%%%%%%%%%%%%%%%%%%%%%%%%%%%%%%%%%%%%%%%

\begin{theorem}

Confluence implies the Church--Rosser property.

\end{theorem}

%%%%%%%%%%%%%%%%%%%%%%%%%%%%%%%%%%%%%%%%%%%%%%%%%%%%%%%%%%%%%%%

\section{Newman's Lemma}

\begin{theorem}[Newman]

Every terminating and locally confluent reduction system is confluent.

\end{theorem}

\begin{proof}

Proceed by induction on reduction height.

Termination provides a well-founded ordering.

Local confluence resolves one-step divergences.

Induction extends this to arbitrary reduction sequences.

\end{proof}

%%%%%%%%%%%%%%%%%%%%%%%%%%%%%%%%%%%%%%%%%%%%%%%%%%%%%%%%%%%%%%%

\section{Unique Normal Forms}

\begin{theorem}

If a reduction system is terminating and confluent, then every element
possesses a unique normal form.

\end{theorem}

\begin{proof}

Existence follows from termination.

Uniqueness follows from confluence.

\end{proof}

%%%%%%%%%%%%%%%%%%%%%%%%%%%%%%%%%%%%%%%%%%%%%%%%%%%%%%%%%%%%%%%

\section{Rewrite Equivalence}

\begin{definition}

Two elements

\[
x,
y
\]

are rewrite equivalent if

\[
x
\leftrightarrow^{*}
y.
\]

\end{definition}

%%%%%%%%%%%%%%%%%%%%%%%%%%%%%%%%%%%%%%%%%%%%%%%%%%%%%%%%%%%%%%%

\begin{proposition}

Rewrite equivalence is an equivalence relation.

\end{proposition}

%%%%%%%%%%%%%%%%%%%%%%%%%%%%%%%%%%%%%%%%%%%%%%%%%%%%%%%%%%%%%%%

\section{Critical Pairs}

\begin{definition}

A critical pair consists of two distinct one-step reductions issuing
from a common ancestor.

\end{definition}

%%%%%%%%%%%%%%%%%%%%%%%%%%%%%%%%%%%%%%%%%%%%%%%%%%%%%%%%%%%%%%%

\begin{proposition}

Every unresolved critical pair represents a potential failure of local
confluence.

\end{proposition}

%%%%%%%%%%%%%%%%%%%%%%%%%%%%%%%%%%%%%%%%%%%%%%%%%%%%%%%%%%%%%%%

\section{Canonical Rewrite Systems}

\begin{definition}

A rewrite system is canonical if it is both terminating and confluent.

\end{definition}

%%%%%%%%%%%%%%%%%%%%%%%%%%%%%%%%%%%%%%%%%%%%%%%%%%%%%%%%%%%%%%%

\begin{corollary}

Every canonical rewrite system determines a unique representative of
each rewrite-equivalence class.

\end{corollary}

%%%%%%%%%%%%%%%%%%%%%%%%%%%%%%%%%%%%%%%%%%%%%%%%%%%%%%%%%%%%%%%

\section{Factorization}

\begin{definition}

A factorization of an operator

\[
T
\]

is an expression

\[
T
=
T_n
\circ
\cdots
\circ
T_1
\]

into primitive operators.

\end{definition}

%%%%%%%%%%%%%%%%%%%%%%%%%%%%%%%%%%%%%%%%%%%%%%%%%%%%%%%%%%%%%%%

\begin{proposition}

Suppose primitive operator factorizations form a canonical rewrite
system.

Then every operator possesses a unique normal-form factorization.

\end{proposition}

%%%%%%%%%%%%%%%%%%%%%%%%%%%%%%%%%%%%%%%%%%%%%%%%%%%%%%%%%%%%%%%

\begin{remark}

Chapter~4 uses Newman's Lemma to establish the uniqueness of normal-form
operator decompositions. Primitive persistence operators generate an
abstract rewriting system whose terminating, locally confluent rewrite
rules yield canonical factorizations. This appendix provides the
standard rewriting-theoretic foundation for that construction.

\end{remark}

\chapter{Dynamical Systems}

This appendix summarizes the principal concepts from the theory of
continuous dynamical systems used throughout the monograph.

Unless otherwise stated,

\[
X
\]

denotes a smooth manifold or Banach space.

%%%%%%%%%%%%%%%%%%%%%%%%%%%%%%%%%%%%%%%%%%%%%%%%%%%%%%%%%%%%%%%

\section{Ordinary Differential Equations}

\begin{definition}

An autonomous ordinary differential equation is

\[
\dot{x}
=
f(x),
\]

where

\[
f:X\rightarrow TX
\]

is a vector field.

\end{definition}

%%%%%%%%%%%%%%%%%%%%%%%%%%%%%%%%%%%%%%%%%%%%%%%%%%%%%%%%%%%%%%%

\begin{definition}

A solution is a continuously differentiable curve

\[
x:[0,T)\rightarrow X
\]

satisfying

\[
\dot{x}(t)=f(x(t))
\]

for every

\[
t.
\]

\end{definition}

%%%%%%%%%%%%%%%%%%%%%%%%%%%%%%%%%%%%%%%%%%%%%%%%%%%%%%%%%%%%%%%

\section{Existence and Uniqueness}

\begin{theorem}[Picard--Lindelöf]

Suppose

\[
f
\]

is locally Lipschitz continuous.

Then every initial condition

\[
x(0)=x_0
\]

admits a unique local solution.

\end{theorem}

%%%%%%%%%%%%%%%%%%%%%%%%%%%%%%%%%%%%%%%%%%%%%%%%%%%%%%%%%%%%%%%

\begin{corollary}

If

\[
f
\]

is globally Lipschitz,

then every solution exists for all

\[
t\ge0.
\]

\end{corollary}

%%%%%%%%%%%%%%%%%%%%%%%%%%%%%%%%%%%%%%%%%%%%%%%%%%%%%%%%%%%%%%%

\section{Flows}

\begin{definition}

The flow generated by

\[
f
\]

is the mapping

\[
\varphi_t:X\rightarrow X
\]

defined by

\[
\varphi_t(x_0)=x(t),
\]

where

\[
x(t)
\]

is the solution beginning at

\[
x_0.
\]

\end{definition}

%%%%%%%%%%%%%%%%%%%%%%%%%%%%%%%%%%%%%%%%%%%%%%%%%%%%%%%%%%%%%%%

\begin{proposition}

Flows satisfy

\[
\varphi_0
=
\mathrm{id},
\]

and

\[
\varphi_{t+s}
=
\varphi_t
\circ
\varphi_s.
\]

\end{proposition}

%%%%%%%%%%%%%%%%%%%%%%%%%%%%%%%%%%%%%%%%%%%%%%%%%%%%%%%%%%%%%%%

\section{Invariant Sets}

\begin{definition}

A subset

\[
A\subseteq X
\]

is positively invariant if

\[
\varphi_t(A)
\subseteq
A
\]

for every

\[
t\ge0.
\]

\end{definition}

%%%%%%%%%%%%%%%%%%%%%%%%%%%%%%%%%%%%%%%%%%%%%%%%%%%%%%%%%%%%%%%

\begin{definition}

A subset is invariant if

\[
\varphi_t(A)
=
A
\]

for every

\[
t.
\]

\end{definition}

%%%%%%%%%%%%%%%%%%%%%%%%%%%%%%%%%%%%%%%%%%%%%%%%%%%%%%%%%%%%%%%

\section{Equilibrium Points}

\begin{definition}

A point

\[
x^*
\]

is an equilibrium if

\[
f(x^*)=0.
\]

\end{definition}

%%%%%%%%%%%%%%%%%%%%%%%%%%%%%%%%%%%%%%%%%%%%%%%%%%%%%%%%%%%%%%%

\begin{proposition}

Every equilibrium generates the constant solution

\[
x(t)\equiv x^*.
\]

\end{proposition}

%%%%%%%%%%%%%%%%%%%%%%%%%%%%%%%%%%%%%%%%%%%%%%%%%%%%%%%%%%%%%%%

\section{Linearization}

\begin{definition}

Suppose

\[
f
\]

is differentiable.

The linearization at

\[
x^*
\]

is

\[
A
=
Df(x^*).
\]

\end{definition}

%%%%%%%%%%%%%%%%%%%%%%%%%%%%%%%%%%%%%%%%%%%%%%%%%%%%%%%%%%%%%%%

\begin{theorem}

If every eigenvalue of

\[
A
\]

has negative real part,

then

\[
x^*
\]

is locally asymptotically stable.

\end{theorem}

%%%%%%%%%%%%%%%%%%%%%%%%%%%%%%%%%%%%%%%%%%%%%%%%%%%%%%%%%%%%%%%

\section{Lyapunov Stability}

\begin{definition}

An equilibrium

\[
x^*
\]

is Lyapunov stable if every sufficiently small perturbation remains
arbitrarily close for all future time.

\end{definition}

%%%%%%%%%%%%%%%%%%%%%%%%%%%%%%%%%%%%%%%%%%%%%%%%%%%%%%%%%%%%%%%

\begin{definition}

It is asymptotically stable if it is Lyapunov stable and

\[
x(t)
\rightarrow
x^*.
\]

\end{definition}

%%%%%%%%%%%%%%%%%%%%%%%%%%%%%%%%%%%%%%%%%%%%%%%%%%%%%%%%%%%%%%%

\begin{theorem}[Lyapunov]

Suppose there exists

\[
V:X\rightarrow\mathbb R
\]

such that

\[
V(x)>0,
\qquad
V(x^*)=0,
\]

and

\[
\dot V(x)\le0.
\]

Then

\[
x^*
\]

is Lyapunov stable.

If

\[
\dot V<0
\]

away from

\[
x^*,
\]

then

\[
x^*
\]

is asymptotically stable.

\end{theorem}

%%%%%%%%%%%%%%%%%%%%%%%%%%%%%%%%%%%%%%%%%%%%%%%%%%%%%%%%%%%%%%%

\section{$\omega$-Limit Sets}

\begin{definition}

The

\[
\omega
\]

-limit set of a trajectory is

\[
\omega(x)
=
\left\{
y
:
x(t_n)
\rightarrow
y
\text{ for some }
t_n\rightarrow\infty
\right\}.
\]

\end{definition}

%%%%%%%%%%%%%%%%%%%%%%%%%%%%%%%%%%%%%%%%%%%%%%%%%%%%%%%%%%%%%%%

\begin{proposition}

Every bounded trajectory possesses a nonempty compact

\[
\omega
\]

-limit set.

\end{proposition}

%%%%%%%%%%%%%%%%%%%%%%%%%%%%%%%%%%%%%%%%%%%%%%%%%%%%%%%%%%%%%%%

\section{LaSalle Invariance Principle}

\begin{theorem}[LaSalle]

Let

\[
V
\]

be a Lyapunov function satisfying

\[
\dot V\le0.
\]

Every bounded trajectory approaches the largest invariant subset of

\[
\{
x:
\dot V(x)=0
\}.
\]

\end{theorem}

%%%%%%%%%%%%%%%%%%%%%%%%%%%%%%%%%%%%%%%%%%%%%%%%%%%%%%%%%%%%%%%

\section{Attractors}

\begin{definition}

A compact invariant set

\[
A
\]

is an attractor if there exists an open neighborhood

\[
U
\]

such that

\[
\operatorname{dist}
(\varphi_t(x),A)
\rightarrow0
\]

for every

\[
x\in U.
\]

\end{definition}

%%%%%%%%%%%%%%%%%%%%%%%%%%%%%%%%%%%%%%%%%%%%%%%%%%%%%%%%%%%%%%%

\section{Basins of Attraction}

\begin{definition}

The basin of attraction of

\[
A
\]

is

\[
B(A)
=
\{
x
:
\omega(x)
\subseteq A
\}.
\]

\end{definition}

%%%%%%%%%%%%%%%%%%%%%%%%%%%%%%%%%%%%%%%%%%%%%%%%%%%%%%%%%%%%%%%

\section{Structural Stability}

\begin{definition}

A dynamical system is structurally stable if sufficiently small
perturbations produce topologically equivalent dynamics.

\end{definition}

%%%%%%%%%%%%%%%%%%%%%%%%%%%%%%%%%%%%%%%%%%%%%%%%%%%%%%%%%%%%%%%

\section{Bifurcations}

\begin{definition}

A bifurcation is a qualitative change in the phase portrait resulting
from variation of a parameter.

\end{definition}

%%%%%%%%%%%%%%%%%%%%%%%%%%%%%%%%%%%%%%%%%%%%%%%%%%%%%%%%%%%%%%%

\begin{example}

Common bifurcations include

\begin{itemize}

\item saddle-node,

\item transcritical,

\item pitchfork,

\item Hopf.

\end{itemize}

\end{example}

%%%%%%%%%%%%%%%%%%%%%%%%%%%%%%%%%%%%%%%%%%%%%%%%%%%%%%%%%%%%%%%

\section{Morse Decompositions}

\begin{definition}

A Morse decomposition partitions an invariant set into finitely many
disjoint invariant subsets together with a partial ordering induced by
connecting trajectories.

\end{definition}

%%%%%%%%%%%%%%%%%%%%%%%%%%%%%%%%%%%%%%%%%%%%%%%%%%%%%%%%%%%%%%%

\begin{theorem}

Every compact isolated invariant set admits a Morse decomposition.

\end{theorem}

%%%%%%%%%%%%%%%%%%%%%%%%%%%%%%%%%%%%%%%%%%%%%%%%%%%%%%%%%%%%%%%

\begin{remark}

Chapter~3 relies on flow invariance, existence and uniqueness of
solutions, and LaSalle's invariance principle when studying admissible
continuation. Chapters~6 and~7 employ attractors, $\omega$-limit sets,
and structural stability to analyze recursive continuation and the
evolution of persistence operators under continuous dynamics.

\end{remark}

\chapter{Differential Geometry}

This appendix summarizes the differential geometric concepts employed
throughout the monograph. Unless otherwise stated, all manifolds are
assumed to be smooth (\(C^\infty\)).

%%%%%%%%%%%%%%%%%%%%%%%%%%%%%%%%%%%%%%%%%%%%%%%%%%%%%%%%%%%%%%%

\section{Smooth Manifolds}

\begin{definition}

An \(n\)-dimensional smooth manifold is a Hausdorff,
second-countable topological space

\[
M
\]

equipped with an atlas of coordinate charts

\[
(U_\alpha,\varphi_\alpha)
\]

such that every transition map

\[
\varphi_\beta
\circ
\varphi_\alpha^{-1}
\]

is infinitely differentiable.

\end{definition}

%%%%%%%%%%%%%%%%%%%%%%%%%%%%%%%%%%%%%%%%%%%%%%%%%%%%%%%%%%%%%%%

\begin{example}

Examples include

\[
\mathbb{R}^n,
\qquad
S^n,
\qquad
T^n,
\]

and smooth configuration spaces arising from constrained mechanical
systems.

\end{example}

%%%%%%%%%%%%%%%%%%%%%%%%%%%%%%%%%%%%%%%%%%%%%%%%%%%%%%%%%%%%%%%

\section{Tangent Spaces}

\begin{definition}

For

\[
p\in M,
\]

the tangent space

\[
T_pM
\]

is the vector space consisting of all derivations acting on smooth
functions defined near

\[
p.
\]

\end{definition}

%%%%%%%%%%%%%%%%%%%%%%%%%%%%%%%%%%%%%%%%%%%%%%%%%%%%%%%%%%%%%%%

\begin{definition}

The tangent bundle is

\[
TM
=
\coprod_{p\in M}
T_pM.
\]

\end{definition}

%%%%%%%%%%%%%%%%%%%%%%%%%%%%%%%%%%%%%%%%%%%%%%%%%%%%%%%%%%%%%%%

\section{Cotangent Bundle}

\begin{definition}

The cotangent space

\[
T_p^*M
\]

is the dual vector space

\[
(T_pM)^*.
\]

The cotangent bundle is

\[
T^*M
=
\coprod_{p\in M}
T_p^*M.
\]

\end{definition}

%%%%%%%%%%%%%%%%%%%%%%%%%%%%%%%%%%%%%%%%%%%%%%%%%%%%%%%%%%%%%%%

\section{Vector Fields}

\begin{definition}

A vector field is a smooth section

\[
X:M\rightarrow TM
\]

satisfying

\[
\pi\circ X=\mathrm{id}_M,
\]

where

\[
\pi:TM\rightarrow M
\]

is the natural projection.

\end{definition}

%%%%%%%%%%%%%%%%%%%%%%%%%%%%%%%%%%%%%%%%%%%%%%%%%%%%%%%%%%%%%%%

\begin{definition}

An integral curve of

\[
X
\]

is a curve

\[
\gamma(t)
\]

satisfying

\[
\dot\gamma(t)
=
X(\gamma(t)).
\]

\end{definition}

%%%%%%%%%%%%%%%%%%%%%%%%%%%%%%%%%%%%%%%%%%%%%%%%%%%%%%%%%%%%%%%

\section{Differential Forms}

\begin{definition}

A differential \(k\)-form on

\[
M
\]

is a smooth alternating multilinear mapping

\[
\omega
:
(TM)^k
\rightarrow
\mathbb{R}.
\]

\end{definition}

%%%%%%%%%%%%%%%%%%%%%%%%%%%%%%%%%%%%%%%%%%%%%%%%%%%%%%%%%%%%%%%

\begin{definition}

The exterior derivative

\[
d
\]

maps

\[
k
\]

-forms to

\[
(k+1)
\]

-forms and satisfies

\[
d^2=0.
\]

\end{definition}

%%%%%%%%%%%%%%%%%%%%%%%%%%%%%%%%%%%%%%%%%%%%%%%%%%%%%%%%%%%%%%%

\section{Pullbacks}

\begin{definition}

Let

\[
F:M\rightarrow N
\]

be smooth.

The pullback

\[
F^*
\]

acts on differential forms via

\[
F^*\omega.
\]

\end{definition}

%%%%%%%%%%%%%%%%%%%%%%%%%%%%%%%%%%%%%%%%%%%%%%%%%%%%%%%%%%%%%%%

\begin{proposition}

Pullback commutes with exterior differentiation:

\[
d(F^*\omega)
=
F^*(d\omega).
\]

\end{proposition}

%%%%%%%%%%%%%%%%%%%%%%%%%%%%%%%%%%%%%%%%%%%%%%%%%%%%%%%%%%%%%%%

\section{Lie Derivatives}

\begin{definition}

Let

\[
X
\]

be a vector field.

The Lie derivative of a tensor field

\[
T
\]

along

\[
X
\]

is denoted

\[
\mathcal L_XT.
\]

\end{definition}

%%%%%%%%%%%%%%%%%%%%%%%%%%%%%%%%%%%%%%%%%%%%%%%%%%%%%%%%%%%%%%%

\begin{proposition}[Cartan's Formula]

For every differential form

\[
\omega,
\]

\[
\mathcal L_X\omega
=
d(i_X\omega)
+
i_X(d\omega).
\]

\end{proposition}

%%%%%%%%%%%%%%%%%%%%%%%%%%%%%%%%%%%%%%%%%%%%%%%%%%%%%%%%%%%%%%%

\section{Lie Brackets}

\begin{definition}

Let

\[
X,
Y
\]

be vector fields.

Their Lie bracket is

\[
[X,Y]
=
XY
-
YX.
\]

\end{definition}

%%%%%%%%%%%%%%%%%%%%%%%%%%%%%%%%%%%%%%%%%%%%%%%%%%%%%%%%%%%%%%%

\begin{proposition}

The Lie bracket satisfies

\begin{enumerate}

\item Bilinearity,

\item Antisymmetry,

\item Jacobi Identity.

\end{enumerate}

\end{proposition}

%%%%%%%%%%%%%%%%%%%%%%%%%%%%%%%%%%%%%%%%%%%%%%%%%%%%%%%%%%%%%%%

\section{Riemannian Metrics}

\begin{definition}

A Riemannian metric on

\[
M
\]

is a smoothly varying positive-definite inner product

\[
g_p
:
T_pM
\times
T_pM
\rightarrow
\mathbb R.
\]

\end{definition}

%%%%%%%%%%%%%%%%%%%%%%%%%%%%%%%%%%%%%%%%%%%%%%%%%%%%%%%%%%%%%%%

\begin{definition}

The norm induced by

\[
g
\]

is

\[
\|v\|
=
\sqrt{g(v,v)}.
\]

\end{definition}

%%%%%%%%%%%%%%%%%%%%%%%%%%%%%%%%%%%%%%%%%%%%%%%%%%%%%%%%%%%%%%%

\section{Connections}

\begin{definition}

An affine connection is an operator

\[
\nabla
\]

assigning to vector fields

\[
X,Y
\]

another vector field

\[
\nabla_XY
\]

satisfying linearity and the Leibniz rule.

\end{definition}

%%%%%%%%%%%%%%%%%%%%%%%%%%%%%%%%%%%%%%%%%%%%%%%%%%%%%%%%%%%%%%%

\section{Geodesics}

\begin{definition}

A geodesic is a curve

\[
\gamma(t)
\]

satisfying

\[
\nabla_{\dot\gamma}
\dot\gamma
=
0.
\]

\end{definition}

%%%%%%%%%%%%%%%%%%%%%%%%%%%%%%%%%%%%%%%%%%%%%%%%%%%%%%%%%%%%%%%

\begin{proposition}

Every point of a smooth manifold possesses a neighborhood in which
geodesics exist uniquely for prescribed initial position and velocity.

\end{proposition}

%%%%%%%%%%%%%%%%%%%%%%%%%%%%%%%%%%%%%%%%%%%%%%%%%%%%%%%%%%%%%%%

\section{Exponential Map}

\begin{definition}

The exponential map at

\[
p
\]

is

\[
\exp_p
:
T_pM
\rightarrow
M,
\]

defined by

\[
\exp_p(v)
=
\gamma_v(1),
\]

where

\[
\gamma_v
\]

is the geodesic satisfying

\[
\gamma_v(0)=p,
\qquad
\dot\gamma_v(0)=v.
\]

\end{definition}

%%%%%%%%%%%%%%%%%%%%%%%%%%%%%%%%%%%%%%%%%%%%%%%%%%%%%%%%%%%%%%%

\section{Curvature}

\begin{definition}

The Riemann curvature tensor is

\[
R(X,Y)Z
=
\nabla_X\nabla_YZ
-
\nabla_Y\nabla_XZ
-
\nabla_{[X,Y]}Z.
\]

\end{definition}

%%%%%%%%%%%%%%%%%%%%%%%%%%%%%%%%%%%%%%%%%%%%%%%%%%%%%%%%%%%%%%%

\section{Embedded Submanifolds}

\begin{definition}

A subset

\[
N\subseteq M
\]

is an embedded submanifold if every point possesses local coordinates
in which

\[
N
\]

is described by the vanishing of a collection of coordinate functions.

\end{definition}

%%%%%%%%%%%%%%%%%%%%%%%%%%%%%%%%%%%%%%%%%%%%%%%%%%%%%%%%%%%%%%%

\begin{theorem}[Regular Level Set]

Let

\[
F:M\rightarrow\mathbb R^k
\]

be smooth.

If

\[
c
\]

is a regular value of

\[
F,
\]

then

\[
F^{-1}(c)
\]

is an embedded submanifold of codimension

\[
k.
\]

\end{theorem}

%%%%%%%%%%%%%%%%%%%%%%%%%%%%%%%%%%%%%%%%%%%%%%%%%%%%%%%%%%%%%%%

\begin{remark}

Chapters~3 and~5 model admissibility manifolds as embedded submanifolds
or Banach manifolds. Tangent spaces, vector fields, Lie derivatives,
connections, and geodesics provide the geometric language for
continuation, repair trajectories, and operator dynamics developed in
the main text.

\end{remark}

\chapter{Category Theory}

This appendix summarizes the categorical framework underlying the
operator calculus developed throughout the monograph. Categories provide
an abstract language for compositional systems, persistence operators,
and recursive constructions.

%%%%%%%%%%%%%%%%%%%%%%%%%%%%%%%%%%%%%%%%%%%%%%%%%%%%%%%%%%%%%%%

\section{Categories}

\begin{definition}

A category

\[
\mathcal C
\]

consists of

\begin{enumerate}

\item
a class of objects,

\item
for every pair of objects

\[
A,B,
\]

a set

\[
\operatorname{Hom}_{\mathcal C}(A,B)
\]

of morphisms,

\item
an associative composition law,

\item
identity morphisms

\[
\mathrm{id}_A
\]

for every object.

\end{enumerate}

\end{definition}

%%%%%%%%%%%%%%%%%%%%%%%%%%%%%%%%%%%%%%%%%%%%%%%%%%%%%%%%%%%%%%%

\begin{proposition}

For every morphism

\[
f:A\rightarrow B,
\]

\[
f\circ\mathrm{id}_A
=
f
=
\mathrm{id}_B\circ f.
\]

\end{proposition}

%%%%%%%%%%%%%%%%%%%%%%%%%%%%%%%%%%%%%%%%%%%%%%%%%%%%%%%%%%%%%%%

\section{Examples}

\begin{example}

Typical categories include

\[
\mathbf{Set},
\qquad
\mathbf{Top},
\qquad
\mathbf{Grp},
\qquad
\mathbf{Vect},
\qquad
\mathbf{Ban},
\]

and

\[
\mathbf{Man}.
\]

\end{example}

%%%%%%%%%%%%%%%%%%%%%%%%%%%%%%%%%%%%%%%%%%%%%%%%%%%%%%%%%%%%%%%

\section{Isomorphisms}

\begin{definition}

A morphism

\[
f:A\rightarrow B
\]

is an isomorphism if there exists

\[
g:B\rightarrow A
\]

such that

\[
gf=\mathrm{id}_A,
\qquad
fg=\mathrm{id}_B.
\]

\end{definition}

%%%%%%%%%%%%%%%%%%%%%%%%%%%%%%%%%%%%%%%%%%%%%%%%%%%%%%%%%%%%%%%

\section{Monomorphisms}

\begin{definition}

A morphism

\[
m
\]

is monic if

\[
mf=mg
\Longrightarrow
f=g.
\]

\end{definition}

%%%%%%%%%%%%%%%%%%%%%%%%%%%%%%%%%%%%%%%%%%%%%%%%%%%%%%%%%%%%%%%

\section{Epimorphisms}

\begin{definition}

A morphism

\[
e
\]

is epic if

\[
fe=ge
\Longrightarrow
f=g.
\]

\end{definition}

%%%%%%%%%%%%%%%%%%%%%%%%%%%%%%%%%%%%%%%%%%%%%%%%%%%%%%%%%%%%%%%

\section{Functors}

\begin{definition}

A functor

\[
F:\mathcal C\rightarrow\mathcal D
\]

assigns

\begin{enumerate}

\item
objects to objects,

\item
morphisms to morphisms,

\end{enumerate}

such that

\[
F(gf)
=
F(g)F(f),
\]

and

\[
F(\mathrm{id})
=
\mathrm{id}.
\]

\end{definition}

%%%%%%%%%%%%%%%%%%%%%%%%%%%%%%%%%%%%%%%%%%%%%%%%%%%%%%%%%%%%%%%

\begin{example}

The tangent bundle construction

\[
T:\mathbf{Man}
\rightarrow
\mathbf{Man}
\]

is a functor.

\end{example}

%%%%%%%%%%%%%%%%%%%%%%%%%%%%%%%%%%%%%%%%%%%%%%%%%%%%%%%%%%%%%%%

\section{Natural Transformations}

\begin{definition}

Let

\[
F,G:
\mathcal C
\rightarrow
\mathcal D
\]

be functors.

A natural transformation

\[
\eta:F\Rightarrow G
\]

assigns morphisms

\[
\eta_A:
F(A)
\rightarrow
G(A)
\]

such that

\[
G(f)\eta_A
=
\eta_BF(f)
\]

for every morphism

\[
f:A\rightarrow B.
\]

\end{definition}

%%%%%%%%%%%%%%%%%%%%%%%%%%%%%%%%%%%%%%%%%%%%%%%%%%%%%%%%%%%%%%%

\section{Commutative Diagrams}

\begin{definition}

A diagram is commutative if every directed path having identical
endpoints determines the same composite morphism.

\end{definition}

%%%%%%%%%%%%%%%%%%%%%%%%%%%%%%%%%%%%%%%%%%%%%%%%%%%%%%%%%%%%%%%

\section{Products}

\begin{definition}

A product of

\[
A
\]

and

\[
B
\]

consists of an object

\[
A\times B
\]

together with projections satisfying the universal property.

\end{definition}

%%%%%%%%%%%%%%%%%%%%%%%%%%%%%%%%%%%%%%%%%%%%%%%%%%%%%%%%%%%%%%%

\section{Coproducts}

\begin{definition}

A coproduct

\[
A\amalg B
\]

is defined dually by injections satisfying the corresponding universal
property.

\end{definition}

%%%%%%%%%%%%%%%%%%%%%%%%%%%%%%%%%%%%%%%%%%%%%%%%%%%%%%%%%%%%%%%

\section{Equalizers}

\begin{definition}

Given

\[
f,g:A\rightarrow B,
\]

the equalizer is the universal object through which every morphism

satisfying

\[
fh=gh
\]

factors uniquely.

\end{definition}

%%%%%%%%%%%%%%%%%%%%%%%%%%%%%%%%%%%%%%%%%%%%%%%%%%%%%%%%%%%%%%%

\section{Pullbacks}

\begin{definition}

A pullback is the universal limit of two morphisms having a common
codomain.

\end{definition}

%%%%%%%%%%%%%%%%%%%%%%%%%%%%%%%%%%%%%%%%%%%%%%%%%%%%%%%%%%%%%%%

\section{Pushouts}

\begin{definition}

A pushout is the dual construction of a pullback.

\end{definition}

%%%%%%%%%%%%%%%%%%%%%%%%%%%%%%%%%%%%%%%%%%%%%%%%%%%%%%%%%%%%%%%

\section{Limits}

\begin{definition}

A limit of a diagram is a universal cone over that diagram.

\end{definition}

%%%%%%%%%%%%%%%%%%%%%%%%%%%%%%%%%%%%%%%%%%%%%%%%%%%%%%%%%%%%%%%

\section{Colimits}

\begin{definition}

A colimit is a universal cocone.

\end{definition}

%%%%%%%%%%%%%%%%%%%%%%%%%%%%%%%%%%%%%%%%%%%%%%%%%%%%%%%%%%%%%%%

\section{Adjoint Functors}

\begin{definition}

Functors

\[
F
\dashv
G
\]

are adjoint if there exists a natural bijection

\[
\operatorname{Hom}(F(A),B)
\cong
\operatorname{Hom}(A,G(B)).
\]

\end{definition}

%%%%%%%%%%%%%%%%%%%%%%%%%%%%%%%%%%%%%%%%%%%%%%%%%%%%%%%%%%%%%%%

\section{Monads}

\begin{definition}

A monad consists of

\[
(T,\eta,\mu)
\]

where

\[
T
\]

is an endofunctor,

\[
\eta
\]

is the unit,

and

\[
\mu
\]

is the multiplication satisfying the associativity and unit axioms.

\end{definition}

%%%%%%%%%%%%%%%%%%%%%%%%%%%%%%%%%%%%%%%%%%%%%%%%%%%%%%%%%%%%%%%

\section{Yoneda Lemma}

\begin{theorem}[Yoneda]

For every locally small category

\[
\mathcal C,
\]

natural transformations

\[
\operatorname{Hom}(A,-)
\Rightarrow
F
\]

are in natural bijection with

\[
F(A).
\]

\end{theorem}

%%%%%%%%%%%%%%%%%%%%%%%%%%%%%%%%%%%%%%%%%%%%%%%%%%%%%%%%%%%%%%%

\begin{remark}

The Yoneda Lemma characterizes an object entirely by its relationships
to every other object in the category.

\end{remark}

%%%%%%%%%%%%%%%%%%%%%%%%%%%%%%%%%%%%%%%%%%%%%%%%%%%%%%%%%%%%%%%

\section{Equivalences of Categories}

\begin{definition}

Categories

\[
\mathcal C
\]

and

\[
\mathcal D
\]

are equivalent if there exist functors

\[
F
\]

and

\[
G
\]

whose composites are naturally isomorphic to the identity functors.

\end{definition}

%%%%%%%%%%%%%%%%%%%%%%%%%%%%%%%%%%%%%%%%%%%%%%%%%%%%%%%%%%%%%%%

\section{Categorical Semantics}

\begin{proposition}

Associativity of morphism composition implies that composite operator
pipelines are independent of parenthesization.

\end{proposition}

%%%%%%%%%%%%%%%%%%%%%%%%%%%%%%%%%%%%%%%%%%%%%%%%%%%%%%%%%%%%%%%

\begin{remark}

The operator calculus developed in Chapters~4--8 naturally forms a
category whose objects are admissible configuration spaces and whose
morphisms are persistence operators. Functors describe translations
between representations, natural transformations encode coherent changes
of interpretation, and universal constructions organize recursive repair
architectures. Appendix~J extends these ideas to Kan extensions and
higher categorical structures.

\end{remark}

\chapter{Sheaf Theory}

This appendix summarizes the basic notions of sheaf theory used to
formalize local-to-global reconstruction, consistency, and information
integration throughout the monograph.

%%%%%%%%%%%%%%%%%%%%%%%%%%%%%%%%%%%%%%%%%%%%%%%%%%%%%%%%%%%%%%%

\section{Presheaves}

\begin{definition}

Let

\[
X
\]

be a topological space.

A presheaf

\[
\mathcal F
\]

on

\[
X
\]

assigns

\begin{enumerate}

\item
to every open set

\[
U\subseteq X,
\]

an object

\[
\mathcal F(U),
\]

called the sections over

\[
U,
\]

\item
to every inclusion

\[
V\subseteq U,
\]

a restriction morphism

\[
\rho^U_V:
\mathcal F(U)
\rightarrow
\mathcal F(V),
\]

such that

\[
\rho^U_U=\mathrm{id},
\]

and

\[
\rho^W_V
=
\rho^U_V
\circ
\rho^W_U.
\]

\end{enumerate}

\end{definition}

%%%%%%%%%%%%%%%%%%%%%%%%%%%%%%%%%%%%%%%%%%%%%%%%%%%%%%%%%%%%%%%

\section{Sheaves}

\begin{definition}

A presheaf is a sheaf if it satisfies

\begin{enumerate}

\item
(Locality)

If two sections agree on every member of an open cover, they are equal.

\item
(Gluing)

Compatible local sections admit a unique global section.

\end{enumerate}

\end{definition}

%%%%%%%%%%%%%%%%%%%%%%%%%%%%%%%%%%%%%%%%%%%%%%%%%%%%%%%%%%%%%%%

\section{Open Covers}

\begin{definition}

An open cover of

\[
U
\]

is a collection

\[
\{U_i\}_{i\in I}
\]

satisfying

\[
U
=
\bigcup_{i\in I}
U_i.
\]

\end{definition}

%%%%%%%%%%%%%%%%%%%%%%%%%%%%%%%%%%%%%%%%%%%%%%%%%%%%%%%%%%%%%%%

\section{Compatible Families}

\begin{definition}

A family

\[
\{s_i\}
\]

of sections over an open cover is compatible if

\[
s_i|_{U_i\cap U_j}
=
s_j|_{U_i\cap U_j}
\]

for every pair

\[
i,j.
\]

\end{definition}

%%%%%%%%%%%%%%%%%%%%%%%%%%%%%%%%%%%%%%%%%%%%%%%%%%%%%%%%%%%%%%%

\section{Global Sections}

\begin{definition}

A global section is an element of

\[
\mathcal F(X).
\]

\end{definition}

%%%%%%%%%%%%%%%%%%%%%%%%%%%%%%%%%%%%%%%%%%%%%%%%%%%%%%%%%%%%%%%

\section{Stalks}

\begin{definition}

The stalk of

\[
\mathcal F
\]

at

\[
x\in X
\]

is the direct limit

\[
\mathcal F_x
=
\varinjlim_{x\in U}
\mathcal F(U).
\]

\end{definition}

%%%%%%%%%%%%%%%%%%%%%%%%%%%%%%%%%%%%%%%%%%%%%%%%%%%%%%%%%%%%%%%

\begin{remark}

The stalk represents all local information arbitrarily close to

\[
x.
\]

\end{remark}

%%%%%%%%%%%%%%%%%%%%%%%%%%%%%%%%%%%%%%%%%%%%%%%%%%%%%%%%%%%%%%%

\section{Morphisms of Sheaves}

\begin{definition}

A morphism

\[
\varphi:
\mathcal F
\rightarrow
\mathcal G
\]

consists of compatible morphisms

\[
\varphi_U:
\mathcal F(U)
\rightarrow
\mathcal G(U)
\]

commuting with every restriction map.

\end{definition}

%%%%%%%%%%%%%%%%%%%%%%%%%%%%%%%%%%%%%%%%%%%%%%%%%%%%%%%%%%%%%%%

\section{Constant Sheaves}

\begin{definition}

Given a set

\[
A,
\]

the constant sheaf

\[
\underline{A}
\]

assigns locally constant functions into

\[
A.
\]

\end{definition}

%%%%%%%%%%%%%%%%%%%%%%%%%%%%%%%%%%%%%%%%%%%%%%%%%%%%%%%%%%%%%%%

\section{Sheafification}

\begin{theorem}

Every presheaf admits a sheafification.

\end{theorem}

%%%%%%%%%%%%%%%%%%%%%%%%%%%%%%%%%%%%%%%%%%%%%%%%%%%%%%%%%%%%%%%

\begin{remark}

Sheafification is left adjoint to the forgetful functor from sheaves to
presheaves.

\end{remark}

%%%%%%%%%%%%%%%%%%%%%%%%%%%%%%%%%%%%%%%%%%%%%%%%%%%%%%%%%%%%%%%

\section{Exact Sequences}

\begin{definition}

A sequence

\[
A
\rightarrow
B
\rightarrow
C
\]

is exact if

\[
\operatorname{im}
=
\ker
\]

at every intermediate object.

\end{definition}

%%%%%%%%%%%%%%%%%%%%%%%%%%%%%%%%%%%%%%%%%%%%%%%%%%%%%%%%%%%%%%%

\section{Sheaf Cohomology}

\begin{definition}

The cohomology groups

\[
H^k(X,\mathcal F)
\]

measure the obstruction to extending compatible local sections into
global sections.

\end{definition}

%%%%%%%%%%%%%%%%%%%%%%%%%%%%%%%%%%%%%%%%%%%%%%%%%%%%%%%%%%%%%%%

\begin{remark}

When

\[
H^1
=
0,
\]

local compatibility is sufficient for global reconstruction.

\end{remark}

%%%%%%%%%%%%%%%%%%%%%%%%%%%%%%%%%%%%%%%%%%%%%%%%%%%%%%%%%%%%%%%

\section{Soft and Fine Sheaves}

\begin{definition}

A sheaf is soft if sections over closed subsets extend to global
sections.

\end{definition}

%%%%%%%%%%%%%%%%%%%%%%%%%%%%%%%%%%%%%%%%%%%%%%%%%%%%%%%%%%%%%%%

\begin{definition}

A sheaf is fine if it admits partitions of unity subordinate to every
locally finite open cover.

\end{definition}

%%%%%%%%%%%%%%%%%%%%%%%%%%%%%%%%%%%%%%%%%%%%%%%%%%%%%%%%%%%%%%%

\begin{theorem}

Every fine sheaf is acyclic.

\end{theorem}

%%%%%%%%%%%%%%%%%%%%%%%%%%%%%%%%%%%%%%%%%%%%%%%%%%%%%%%%%%%%%%%

\section{Local Systems}

\begin{definition}

A local system is a sheaf that is locally constant.

\end{definition}

%%%%%%%%%%%%%%%%%%%%%%%%%%%%%%%%%%%%%%%%%%%%%%%%%%%%%%%%%%%%%%%

\section{Constructible Sheaves}

\begin{definition}

A constructible sheaf is locally constant on each member of a finite
stratification of the underlying space.

\end{definition}

%%%%%%%%%%%%%%%%%%%%%%%%%%%%%%%%%%%%%%%%%%%%%%%%%%%%%%%%%%%%%%%

\section{Čech Cohomology}

\begin{definition}

Given an open cover

\[
\mathcal U,
\]

the Čech complex is formed from compatible sections on finite
intersections of members of

\[
\mathcal U.
\]

Its cohomology approximates the sheaf cohomology of

\[
\mathcal F.
\]

\end{definition}

%%%%%%%%%%%%%%%%%%%%%%%%%%%%%%%%%%%%%%%%%%%%%%%%%%%%%%%%%%%%%%%

\section{Local-to-Global Principle}

\begin{theorem}

Suppose

\[
\mathcal F
\]

is a sheaf on

\[
X.
\]

If

\begin{enumerate}

\item
local sections exist,

\item
they agree on all overlaps,

\end{enumerate}

then there exists a unique global section restricting to each local
section.

\end{theorem}

%%%%%%%%%%%%%%%%%%%%%%%%%%%%%%%%%%%%%%%%%%%%%%%%%%%%%%%%%%%%%%%

\begin{proof}

Apply the gluing axiom of sheaves.

\end{proof}

%%%%%%%%%%%%%%%%%%%%%%%%%%%%%%%%%%%%%%%%%%%%%%%%%%%%%%%%%%%%%%%

\section{Functoriality}

\begin{proposition}

Every continuous map

\[
f:X\rightarrow Y
\]

induces inverse-image and direct-image functors

\[
f^{-1}
\qquad\text{and}\qquad
f_*,
\]

between categories of sheaves.

\end{proposition}

%%%%%%%%%%%%%%%%%%%%%%%%%%%%%%%%%%%%%%%%%%%%%%%%%%%%%%%%%%%%%%%

\begin{remark}

Throughout the main text, local observations, repair constraints, and
partial reconstructions are modeled as sections over open subsets of a
configuration space. Consistency on overlapping domains corresponds to
compatibility of sections, while successful global persistence is
expressed as the existence of a unique glued section. Sheaf cohomology
provides a natural language for describing obstructions to global
reconstruction, and Chapter~8 builds on these ideas using categorical
extensions and universal constructions.

\end{remark}

\chapter{Kan Extensions and Higher Categories}

This appendix summarizes the universal constructions underlying the
higher-level categorical framework employed in the final chapters of the
monograph. Particular emphasis is placed on Kan extensions, adjunctions,
and higher morphisms.

%%%%%%%%%%%%%%%%%%%%%%%%%%%%%%%%%%%%%%%%%%%%%%%%%%%%%%%%%%%%%%%

\section{Universal Properties}

\begin{definition}

A universal property specifies an object uniquely (up to unique
isomorphism) by the existence of a unique factorization through every
competing object.

\end{definition}

%%%%%%%%%%%%%%%%%%%%%%%%%%%%%%%%%%%%%%%%%%%%%%%%%%%%%%%%%%%%%%%

\begin{proposition}

Objects characterized by the same universal property are uniquely
isomorphic.

\end{proposition}

%%%%%%%%%%%%%%%%%%%%%%%%%%%%%%%%%%%%%%%%%%%%%%%%%%%%%%%%%%%%%%%

\section{Comma Categories}

\begin{definition}

Given functors

\[
F:\mathcal A\rightarrow\mathcal C,
\qquad
G:\mathcal B\rightarrow\mathcal C,
\]

the comma category

\[
(F\downarrow G)
\]

has objects

\[
(A,B,\alpha),
\]

where

\[
\alpha:
F(A)\rightarrow G(B).
\]

\end{definition}

%%%%%%%%%%%%%%%%%%%%%%%%%%%%%%%%%%%%%%%%%%%%%%%%%%%%%%%%%%%%%%%

\section{Kan Extensions}

\begin{definition}

Let

\[
F:\mathcal C\rightarrow\mathcal D,
\qquad
K:\mathcal C\rightarrow\mathcal E.
\]

A left Kan extension of

\[
F
\]

along

\[
K
\]

is a functor

\[
\operatorname{Lan}_K F
:
\mathcal E
\rightarrow
\mathcal D
\]

satisfying the corresponding universal property.

\end{definition}

%%%%%%%%%%%%%%%%%%%%%%%%%%%%%%%%%%%%%%%%%%%%%%%%%%%%%%%%%%%%%%%

\begin{definition}

The right Kan extension

\[
\operatorname{Ran}_K F
\]

is defined dually.

\end{definition}

%%%%%%%%%%%%%%%%%%%%%%%%%%%%%%%%%%%%%%%%%%%%%%%%%%%%%%%%%%%%%%%

\begin{remark}

Left Kan extensions preserve colimit-like constructions, whereas right
Kan extensions preserve limit-like constructions.

\end{remark}

%%%%%%%%%%%%%%%%%%%%%%%%%%%%%%%%%%%%%%%%%%%%%%%%%%%%%%%%%%%%%%%

\section{Pointwise Kan Extensions}

\begin{definition}

A Kan extension is pointwise if it may be computed using limits or
colimits over comma categories.

\end{definition}

%%%%%%%%%%%%%%%%%%%%%%%%%%%%%%%%%%%%%%%%%%%%%%%%%%%%%%%%%%%%%%%

\begin{theorem}

Whenever the required limits or colimits exist, pointwise Kan
extensions exist.

\end{theorem}

%%%%%%%%%%%%%%%%%%%%%%%%%%%%%%%%%%%%%%%%%%%%%%%%%%%%%%%%%%%%%%%

\section{Density}

\begin{definition}

A functor

\[
K
\]

is dense if every object of the codomain is obtained canonically as a
colimit of objects in the image of

\[
K.
\]

\end{definition}

%%%%%%%%%%%%%%%%%%%%%%%%%%%%%%%%%%%%%%%%%%%%%%%%%%%%%%%%%%%%%%%

\section{Ends}

\begin{definition}

Given

\[
F:
\mathcal C^{op}
\times
\mathcal C
\rightarrow
\mathcal D,
\]

an end is the universal object

\[
\int_C
F(C,C).
\]

\end{definition}

%%%%%%%%%%%%%%%%%%%%%%%%%%%%%%%%%%%%%%%%%%%%%%%%%%%%%%%%%%%%%%%

\section{Coends}

\begin{definition}

The dual construction

\[
\int^{C}
F(C,C)
\]

is called the coend.

\end{definition}

%%%%%%%%%%%%%%%%%%%%%%%%%%%%%%%%%%%%%%%%%%%%%%%%%%%%%%%%%%%%%%%

\section{Enriched Categories}

\begin{definition}

A category enriched over a monoidal category

\[
\mathcal V
\]

replaces hom-sets by objects of

\[
\mathcal V.
\]

\end{definition}

%%%%%%%%%%%%%%%%%%%%%%%%%%%%%%%%%%%%%%%%%%%%%%%%%%%%%%%%%%%%%%%

\begin{example}

Metric spaces may be viewed as categories enriched over

\[
([0,\infty],\ge,+).
\]

\end{example}

%%%%%%%%%%%%%%%%%%%%%%%%%%%%%%%%%%%%%%%%%%%%%%%%%%%%%%%%%%%%%%%

\section{2-Categories}

\begin{definition}

A 2-category consists of

\begin{enumerate}

\item objects,

\item 1-morphisms,

\item 2-morphisms between 1-morphisms.

\end{enumerate}

Both horizontal and vertical composition satisfy associativity and unit
laws.

\end{definition}

%%%%%%%%%%%%%%%%%%%%%%%%%%%%%%%%%%%%%%%%%%%%%%%%%%%%%%%%%%%%%%%

\section{Bicategories}

\begin{definition}

A bicategory weakens the associativity and identity axioms so that they
hold only up to coherent natural isomorphism.

\end{definition}

%%%%%%%%%%%%%%%%%%%%%%%%%%%%%%%%%%%%%%%%%%%%%%%%%%%%%%%%%%%%%%%

\section{Natural Isomorphisms}

\begin{definition}

A natural transformation is a natural isomorphism if every component is
an isomorphism.

\end{definition}

%%%%%%%%%%%%%%%%%%%%%%%%%%%%%%%%%%%%%%%%%%%%%%%%%%%%%%%%%%%%%%%

\section{Limits in 2-Categories}

\begin{definition}

A bilimit is the higher-categorical analogue of an ordinary categorical
limit, satisfying its universal property up to coherent isomorphism.

\end{definition}

%%%%%%%%%%%%%%%%%%%%%%%%%%%%%%%%%%%%%%%%%%%%%%%%%%%%%%%%%%%%%%%

\section{Higher Adjunctions}

\begin{definition}

Adjunctions in a 2-category are determined by units and counits together
with coherence conditions expressed by higher morphisms.

\end{definition}

%%%%%%%%%%%%%%%%%%%%%%%%%%%%%%%%%%%%%%%%%%%%%%%%%%%%%%%%%%%%%%%

\section{Coherence}

\begin{theorem}[Mac Lane Coherence]

Every diagram generated by associators and unit constraints in a
monoidal category commutes.

\end{theorem}

%%%%%%%%%%%%%%%%%%%%%%%%%%%%%%%%%%%%%%%%%%%%%%%%%%%%%%%%%%%%%%%

\section{Universal Computation}

\begin{proposition}

Many recursive constructions are naturally expressed as universal
extensions rather than explicit algorithms.

\end{proposition}

%%%%%%%%%%%%%%%%%%%%%%%%%%%%%%%%%%%%%%%%%%%%%%%%%%%%%%%%%%%%%%%

\section{Recursive Diagrams}

\begin{definition}

A recursive diagram is a directed diagram whose objects are generated by
iterated applications of one or more endofunctors.

\end{definition}

%%%%%%%%%%%%%%%%%%%%%%%%%%%%%%%%%%%%%%%%%%%%%%%%%%%%%%%%%%%%%%%

\begin{proposition}

If the corresponding colimit exists, the recursive diagram admits a
canonical universal completion.

\end{proposition}

%%%%%%%%%%%%%%%%%%%%%%%%%%%%%%%%%%%%%%%%%%%%%%%%%%%%%%%%%%%%%%%

\section{Functorial Continuation}

\begin{definition}

A continuation process is functorial if every admissible morphism
commutes with continuation under composition.

\end{definition}

%%%%%%%%%%%%%%%%%%%%%%%%%%%%%%%%%%%%%%%%%%%%%%%%%%%%%%%%%%%%%%%

\begin{proposition}

Functorial continuation preserves identities and composite morphisms.

\end{proposition}

%%%%%%%%%%%%%%%%%%%%%%%%%%%%%%%%%%%%%%%%%%%%%%%%%%%%%%%%%%%%%%%

\section{Categorical Fixed Points}

\begin{definition}

An algebra

\[
(A,\alpha)
\]

for an endofunctor

\[
F
\]

is initial if every other

\[
F
\]

-algebra admits a unique morphism from

\[
(A,\alpha).
\]

\end{definition}

%%%%%%%%%%%%%%%%%%%%%%%%%%%%%%%%%%%%%%%%%%%%%%%%%%%%%%%%%%%%%%%

\begin{remark}

Initial algebras provide categorical models of recursively generated
objects.

\end{remark}

%%%%%%%%%%%%%%%%%%%%%%%%%%%%%%%%%%%%%%%%%%%%%%%%%%%%%%%%%%%%%%%

\section{Summary}

\begin{remark}

The final chapters interpret persistence operators as morphisms within
higher categorical structures. Kan extensions formalize optimal
extensions of partial constructions, enriched categories encode
additional geometric structure, and higher morphisms capture coherent
transformations between operator systems. Universal properties ensure
that recursive constructions are characterized intrinsically rather than
by arbitrary implementation details.

\end{remark}

\chapter{Optimization and Projection Theory}

This appendix collects the principal optimization results underlying the
repair, admissibility, and continuation functionals introduced
throughout the monograph.

Unless otherwise stated,

\[
X
\]

denotes a real Hilbert space.

%%%%%%%%%%%%%%%%%%%%%%%%%%%%%%%%%%%%%%%%%%%%%%%%%%%%%%%%%%%%%%%

\section{Optimization Problems}

\begin{definition}

An optimization problem consists of finding

\[
x^*
\in X
\]

such that

\[
f(x^*)
=
\inf_{x\in X}
f(x).
\]

\end{definition}

%%%%%%%%%%%%%%%%%%%%%%%%%%%%%%%%%%%%%%%%%%%%%%%%%%%%%%%%%%%%%%%

\section{Feasible Sets}

\begin{definition}

A feasible set is a subset

\[
C
\subseteq
X
\]

containing all admissible solutions.

\end{definition}

%%%%%%%%%%%%%%%%%%%%%%%%%%%%%%%%%%%%%%%%%%%%%%%%%%%%%%%%%%%%%%%

\section{Convex Optimization}

\begin{definition}

A convex optimization problem consists of minimizing a convex functional
over a convex feasible set.

\end{definition}

%%%%%%%%%%%%%%%%%%%%%%%%%%%%%%%%%%%%%%%%%%%%%%%%%%%%%%%%%%%%%%%

\begin{theorem}

Every local minimizer of a convex optimization problem is globally
optimal.

\end{theorem}

%%%%%%%%%%%%%%%%%%%%%%%%%%%%%%%%%%%%%%%%%%%%%%%%%%%%%%%%%%%%%%%

\section{Projection onto Convex Sets}

\begin{definition}

Let

\[
C
\subseteq
X
\]

be closed and convex.

The metric projection

\[
P_C:X\rightarrow C
\]

is defined by

\[
P_C(x)
=
\arg\min_{y\in C}
\|x-y\|.
\]

\end{definition}

%%%%%%%%%%%%%%%%%%%%%%%%%%%%%%%%%%%%%%%%%%%%%%%%%%%%%%%%%%%%%%%

\begin{theorem}[Projection Theorem]

Every point

\[
x\in X
\]

admits a unique metric projection onto every nonempty closed convex set.

\end{theorem}

%%%%%%%%%%%%%%%%%%%%%%%%%%%%%%%%%%%%%%%%%%%%%%%%%%%%%%%%%%%%%%%

\begin{proposition}

The projection operator satisfies

\[
\|P_Cx-P_Cy\|
\le
\|x-y\|.
\]

\end{proposition}

%%%%%%%%%%%%%%%%%%%%%%%%%%%%%%%%%%%%%%%%%%%%%%%%%%%%%%%%%%%%%%%

\section{Orthogonality}

\begin{theorem}

The projection

\[
p=P_C(x)
\]

is characterized by

\[
\langle
x-p,
y-p
\rangle
\le
0
\]

for every

\[
y\in C.
\]

\end{theorem}

%%%%%%%%%%%%%%%%%%%%%%%%%%%%%%%%%%%%%%%%%%%%%%%%%%%%%%%%%%%%%%%

\section{Variational Inequalities}

\begin{definition}

Given

\[
F:C\rightarrow X,
\]

the variational inequality problem is to find

\[
x^*
\in C
\]

such that

\[
\langle
F(x^*),
y-x^*
\rangle
\ge
0
\]

for every

\[
y\in C.
\]

\end{definition}

%%%%%%%%%%%%%%%%%%%%%%%%%%%%%%%%%%%%%%%%%%%%%%%%%%%%%%%%%%%%%%%

\section{Euler--Lagrange Equations}

\begin{theorem}

Let

\[
J[\gamma]
=
\int_a^b
L(\gamma,\dot\gamma)\,dt.
\]

If

\[
\gamma
\]

minimizes

\[
J,
\]

then

\[
\frac{d}{dt}
\frac{\partial L}{\partial\dot\gamma}
-
\frac{\partial L}{\partial\gamma}
=
0.
\]

\end{theorem}

%%%%%%%%%%%%%%%%%%%%%%%%%%%%%%%%%%%%%%%%%%%%%%%%%%%%%%%%%%%%%%%

\section{Gradient Descent}

\begin{definition}

The gradient iteration is

\[
x_{n+1}
=
x_n
-
\alpha_n
\nabla f(x_n).
\]

\end{definition}

%%%%%%%%%%%%%%%%%%%%%%%%%%%%%%%%%%%%%%%%%%%%%%%%%%%%%%%%%%%%%%%

\begin{proposition}

If

\[
\nabla f
\]

is Lipschitz continuous and

\[
0<\alpha<\frac{2}{L},
\]

then

\[
f(x_{n+1})
\le
f(x_n).
\]

\end{proposition}

%%%%%%%%%%%%%%%%%%%%%%%%%%%%%%%%%%%%%%%%%%%%%%%%%%%%%%%%%%%%%%%

\section{Projected Gradient Methods}

\begin{definition}

Projected gradient descent is

\[
x_{n+1}
=
P_C
\left(
x_n
-
\alpha
\nabla f(x_n)
\right).
\]

\end{definition}

%%%%%%%%%%%%%%%%%%%%%%%%%%%%%%%%%%%%%%%%%%%%%%%%%%%%%%%%%%%%%%%

\section{Lagrange Multipliers}

\begin{theorem}

Suppose

\[
g(x)=0
\]

defines a regular constraint.

Then every constrained extremum satisfies

\[
\nabla f
=
\lambda
\nabla g.
\]

\end{theorem}

%%%%%%%%%%%%%%%%%%%%%%%%%%%%%%%%%%%%%%%%%%%%%%%%%%%%%%%%%%%%%%%

\section{Karush--Kuhn--Tucker Conditions}

\begin{theorem}

Suppose suitable regularity assumptions hold.

Then every constrained optimum satisfies the KKT conditions consisting
of

\begin{enumerate}

\item primal feasibility,

\item dual feasibility,

\item complementary slackness,

\item stationarity.

\end{enumerate}

\end{theorem}

%%%%%%%%%%%%%%%%%%%%%%%%%%%%%%%%%%%%%%%%%%%%%%%%%%%%%%%%%%%%%%%

\section{Duality}

\begin{definition}

The Lagrangian is

\[
L(x,\lambda)
=
f(x)
+
\lambda g(x).
\]

The dual problem maximizes the corresponding dual functional.

\end{definition}

%%%%%%%%%%%%%%%%%%%%%%%%%%%%%%%%%%%%%%%%%%%%%%%%%%%%%%%%%%%%%%%

\section{Proximal Operators}

\begin{definition}

The proximal operator associated with

\[
f
\]

is

\[
\operatorname{prox}_f(x)
=
\arg\min_y
\left(
f(y)
+
\frac12
\|x-y\|^2
\right).
\]

\end{definition}

%%%%%%%%%%%%%%%%%%%%%%%%%%%%%%%%%%%%%%%%%%%%%%%%%%%%%%%%%%%%%%%

\begin{remark}

Metric projection is the proximal operator of the indicator functional
of a closed convex set.

\end{remark}

%%%%%%%%%%%%%%%%%%%%%%%%%%%%%%%%%%%%%%%%%%%%%%%%%%%%%%%%%%%%%%%

\section{Alternating Projections}

\begin{definition}

Given closed convex sets

\[
A,
B,
\]

the alternating projection iteration is

\[
x_{n+1}
=
P_BP_Ax_n.
\]

\end{definition}

%%%%%%%%%%%%%%%%%%%%%%%%%%%%%%%%%%%%%%%%%%%%%%%%%%%%%%%%%%%%%%%

\begin{theorem}

If

\[
A\cap B\neq\varnothing,
\]

alternating projections converge under standard convexity hypotheses.

\end{theorem}

%%%%%%%%%%%%%%%%%%%%%%%%%%%%%%%%%%%%%%%%%%%%%%%%%%%%%%%%%%%%%%%

\section{Optimization on Manifolds}

\begin{definition}

Optimization on a smooth manifold replaces Euclidean gradients with
Riemannian gradients and Euclidean straight lines with geodesics.

\end{definition}

%%%%%%%%%%%%%%%%%%%%%%%%%%%%%%%%%%%%%%%%%%%%%%%%%%%%%%%%%%%%%%%

\section{Projection Operators and Repair}

\begin{proposition}

Suppose

\[
C
\]

is the admissible configuration manifold.

Then the repair operator

\[
R
=
P_C
\]

is nonexpansive and minimizes reconstruction distance among admissible
configurations.

\end{proposition}

%%%%%%%%%%%%%%%%%%%%%%%%%%%%%%%%%%%%%%%%%%%%%%%%%%%%%%%%%%%%%%%

\section{Variational Continuation}

\begin{proposition}

Suppose

\[
\Gamma
\]

is the collection of admissible continuation paths.

Then

\[
\gamma^*
=
\arg\min_{\gamma\in\Gamma}
J[\gamma]
\]

determines an optimal continuation trajectory whenever the hypotheses of
the direct method of the calculus of variations are satisfied.

\end{proposition}

%%%%%%%%%%%%%%%%%%%%%%%%%%%%%%%%%%%%%%%%%%%%%%%%%%%%%%%%%%%%%%%

\section{Summary}

\begin{remark}

This appendix supplies the optimization machinery supporting Chapters
2--5. Projection theorems justify repair operators, variational
principles characterize continuation paths, and proximal methods unify
repair, admissibility, and optimization within a common mathematical
framework. Together with the preceding appendices, these results provide
the analytical foundation for the persistence calculus developed
throughout the monograph.

\end{remark}

\chapter{Mathematical Dependency Graphs and Reference Tables}

This appendix summarizes the logical structure of the monograph. Rather
than introducing new mathematics, it records the dependency relations
among definitions, propositions, theorems, and appendices.

%%%%%%%%%%%%%%%%%%%%%%%%%%%%%%%%%%%%%%%%%%%%%%%%%%%%%%%%%%%%%%%

\section{Dependency Convention}

\begin{definition}

Let

\[
A
\]

and

\[
B
\]

denote mathematical statements.

We write

\[
A\Longrightarrow B
\]

whenever the proof of

\[
B
\]

depends directly upon

\[
A.
\]

\end{definition}

%%%%%%%%%%%%%%%%%%%%%%%%%%%%%%%%%%%%%%%%%%%%%%%%%%%%%%%%%%%%%%%

\section{Foundational Hierarchy}

\[
\begin{array}{c}
\text{Topology}
\\
\downarrow
\\
\text{Analysis}
\\
\downarrow
\\
\text{Optimization}
\\
\downarrow
\\
\text{Operator Theory}
\\
\downarrow
\\
\text{Dynamics}
\\
\downarrow
\\
\text{Geometry}
\\
\downarrow
\\
\text{Category Theory}
\\
\downarrow
\\
\text{Higher Categories}
\end{array}
\]

%%%%%%%%%%%%%%%%%%%%%%%%%%%%%%%%%%%%%%%%%%%%%%%%%%%%%%%%%%%%%%%

\section{Chapter Dependency Graph}

\[
\begin{array}{c}
\text{Chapter 1}
\\
\downarrow
\\
\text{Chapter 2}
\\
\downarrow
\\
\text{Chapter 3}
\\
\downarrow
\\
\text{Chapter 4}
\\
\downarrow
\\
\text{Chapter 5}
\\
\downarrow
\\
\text{Chapter 6}
\\
\downarrow
\\
\text{Chapter 7}
\\
\downarrow
\\
\text{Chapter 8}
\end{array}
\]

%%%%%%%%%%%%%%%%%%%%%%%%%%%%%%%%%%%%%%%%%%%%%%%%%%%%%%%%%%%%%%%

\section{Appendix Dependencies}

\[
\begin{array}{lll}
A &\rightarrow& B
\\
B &\rightarrow& C
\\
B &\rightarrow& K
\\
C &\rightarrow& D
\\
D &\rightarrow& E
\\
B &\rightarrow& F
\\
A &\rightarrow& G
\\
H &\rightarrow& I
\\
H &\rightarrow& J
\\
I &\rightarrow& J
\end{array}
\]

%%%%%%%%%%%%%%%%%%%%%%%%%%%%%%%%%%%%%%%%%%%%%%%%%%%%%%%%%%%%%%%

\section{Definition Dependencies}

\begin{center}

\begin{tabular}{lll}

Definition &
Depends Upon &
Used In
\\
\hline

Observation Operator &
Topology &
Ch.~2--8
\\

Defect Functional &
Analysis &
Ch.~2--7
\\

Repair Flow &
Dynamics &
Ch.~2--7
\\

Continuation Operator &
Operator Theory &
Ch.~3--8
\\

Persistence Operator &
Operator Theory &
Ch.~4--8
\\

Recursive Continuation &
Fixed Points &
Ch.~6--8
\\

Operator Category &
Category Theory &
Ch.~8

\end{tabular}

\end{center}

%%%%%%%%%%%%%%%%%%%%%%%%%%%%%%%%%%%%%%%%%%%%%%%%%%%%%%%%%%%%%%%

\section{Principal Theorem Dependencies}

\begin{center}

\begin{tabular}{lll}

Theorem &
Depends Upon &
Appendices
\\
\hline

Repair Existence &
Lower Semicontinuity &
A,B,K
\\

Repair Stability &
Lyapunov &
B,F
\\

Continuation &
Nagumo &
F,G
\\

Persistence Factorization &
Newman's Lemma &
E
\\

Recursive Continuation &
Banach Fixed Point &
C,D
\\

Operator Evolution &
Semigroup Theory &
D,F
\\

Functorial Persistence &
Category Theory &
H
\\

Universal Extension &
Kan Extension &
J

\end{tabular}

\end{center}

%%%%%%%%%%%%%%%%%%%%%%%%%%%%%%%%%%%%%%%%%%%%%%%%%%%%%%%%%%%%%%%

\section{Analytical Prerequisites}

\begin{center}

\begin{tabular}{ll}

Result &
Primary Use
\\
\hline

Lower Semicontinuity &
Existence
\\

Coercivity &
Compactness
\\

Fréchet Derivative &
Repair Flows
\\

Grönwall &
Stability Bounds
\\

Barbalat &
Asymptotic Convergence
\\

Arzelà--Ascoli &
Compactness
\\

Projection Theorem &
Repair Operators
\\

Euler--Lagrange &
Continuation Paths

\end{tabular}

\end{center}

%%%%%%%%%%%%%%%%%%%%%%%%%%%%%%%%%%%%%%%%%%%%%%%%%%%%%%%%%%%%%%%

\section{Geometric Prerequisites}

\begin{center}

\begin{tabular}{ll}

Concept &
Application
\\
\hline

Tangent Bundle &
Repair Velocity
\\

Geodesics &
Continuation
\\

Connections &
Operator Evolution
\\

Curvature &
Configuration Geometry
\\

Submanifolds &
Admissibility

\end{tabular}

\end{center}

%%%%%%%%%%%%%%%%%%%%%%%%%%%%%%%%%%%%%%%%%%%%%%%%%%%%%%%%%%%%%%%

\section{Categorical Prerequisites}

\begin{center}

\begin{tabular}{ll}

Concept &
Application
\\
\hline

Functor &
Representation Change
\\

Natural Transformation &
Operator Translation
\\

Adjunction &
Duality
\\

Sheaf &
Local Reconstruction
\\

Kan Extension &
Universal Continuation
\\

Initial Algebra &
Recursive Generation

\end{tabular}

\end{center}

%%%%%%%%%%%%%%%%%%%%%%%%%%%%%%%%%%%%%%%%%%%%%%%%%%%%%%%%%%%%%%%

\section{Notation by Chapter}

\begin{center}

\begin{tabular}{ll}

Chapter &
Principal Symbols
\\
\hline

1 &
$\mathcal O,\mathcal S$
\\

2 &
$\Delta_R,E,R$
\\

3 &
$\Gamma,\Phi,\omega$
\\

4 &
$T,\circ,I$
\\

5 &
$\mathcal M,g,\nabla$
\\

6 &
$T^n,x^*$
\\

7 &
$S_t,A,e^{tA}$
\\

8 &
$\mathcal C,F,\eta,\operatorname{Lan}$

\end{tabular}

\end{center}

%%%%%%%%%%%%%%%%%%%%%%%%%%%%%%%%%%%%%%%%%%%%%%%%%%%%%%%%%%%%%%%

\section{Logical Flow}

The principal logical development of the monograph may be summarized as

\[
\begin{aligned}
\text{Observation}
&\Longrightarrow
\text{Configuration}
\\[2mm]
&\Longrightarrow
\text{Defect}
\\[2mm]
&\Longrightarrow
\text{Repair}
\\[2mm]
&\Longrightarrow
\text{Continuation}
\\[2mm]
&\Longrightarrow
\text{Persistence}
\\[2mm]
&\Longrightarrow
\text{Recursion}
\\[2mm]
&\Longrightarrow
\text{Categories}
\\[2mm]
&\Longrightarrow
\text{Universal Extensions.}
\end{aligned}
\]

%%%%%%%%%%%%%%%%%%%%%%%%%%%%%%%%%%%%%%%%%%%%%%%%%%%%%%%%%%%%%%%

\section{Dependency Principle}

\begin{proposition}

Every theorem appearing in the main text depends only upon definitions,
lemmas, propositions, or appendices that precede it in the dependency
graph.

\end{proposition}

%%%%%%%%%%%%%%%%%%%%%%%%%%%%%%%%%%%%%%%%%%%%%%%%%%%%%%%%%%%%%%%

\begin{remark}

The dependency graph is intentionally acyclic. Consequently, the logical
development of the monograph is well-founded: every proof ultimately
reduces to previously established results from topology, analysis,
geometry, operator theory, or category theory. This organization allows
the text to function both as a sequential exposition and as a reference
work in which the provenance of every major theorem can be traced
systematically.

\end{remark}

\chapter{Notation and Symbol Index}

This appendix collects the principal notation employed throughout the
monograph. Unless explicitly redefined, every symbol retains the meaning
assigned here.

%%%%%%%%%%%%%%%%%%%%%%%%%%%%%%%%%%%%%%%%%%%%%%%%%%%%%%%%%%%%%%%

\section{General Sets}

\begin{center}

\begin{tabular}{ll}

Symbol & Meaning
\\
\hline

$\mathbb N$ &
Natural numbers
\\

$\mathbb Z$ &
Integers
\\

$\mathbb Q$ &
Rational numbers
\\

$\mathbb R$ &
Real numbers
\\

$\mathbb C$ &
Complex numbers
\\

$\varnothing$ &
Empty set
\\

$\mathcal P(X)$ &
Power set
\\

$|X|$ &
Cardinality

\end{tabular}

\end{center}

%%%%%%%%%%%%%%%%%%%%%%%%%%%%%%%%%%%%%%%%%%%%%%%%%%%%%%%%%%%%%%%

\section{Topology}

\begin{center}

\begin{tabular}{ll}

Symbol & Meaning
\\
\hline

$X$ &
Topological space
\\

$\tau$ &
Topology
\\

$\overline A$ &
Closure
\\

$\operatorname{int}(A)$ &
Interior
\\

$\partial A$ &
Boundary
\\

$B(x,r)$ &
Open ball
\\

$d(x,y)$ &
Metric

\end{tabular}

\end{center}

%%%%%%%%%%%%%%%%%%%%%%%%%%%%%%%%%%%%%%%%%%%%%%%%%%%%%%%%%%%%%%%

\section{Analysis}

\begin{center}

\begin{tabular}{ll}

Symbol & Meaning
\\
\hline

$\|\cdot\|$ &
Norm
\\

$\langle\cdot,\cdot\rangle$ &
Inner product
\\

$X^*$ &
Dual space
\\

$Df$ &
Fréchet derivative
\\

$\nabla f$ &
Gradient
\\

$\operatorname{argmin}$ &
Minimizing argument
\\

$\inf,\sup$ &
Infimum, supremum

\end{tabular}

\end{center}

%%%%%%%%%%%%%%%%%%%%%%%%%%%%%%%%%%%%%%%%%%%%%%%%%%%%%%%%%%%%%%%

\section{Geometry}

\begin{center}

\begin{tabular}{ll}

Symbol & Meaning
\\
\hline

$M$ &
Smooth manifold
\\

$TM$ &
Tangent bundle
\\

$T^*M$ &
Cotangent bundle
\\

$T_pM$ &
Tangent space
\\

$g$ &
Riemannian metric
\\

$\nabla$ &
Affine connection
\\

$\exp_p$ &
Exponential map
\\

$R$ &
Curvature tensor

\end{tabular}

\end{center}

%%%%%%%%%%%%%%%%%%%%%%%%%%%%%%%%%%%%%%%%%%%%%%%%%%%%%%%%%%%%%%%

\section{Dynamical Systems}

\begin{center}

\begin{tabular}{ll}

Symbol & Meaning
\\
\hline

$\varphi_t$ &
Flow
\\

$\omega(x)$ &
$\omega$-limit set
\\

$A$ &
Attractor or operator
\\

$V$ &
Lyapunov function
\\

$f$ &
Vector field
\\

$x^*$ &
Equilibrium point

\end{tabular}

\end{center}

%%%%%%%%%%%%%%%%%%%%%%%%%%%%%%%%%%%%%%%%%%%%%%%%%%%%%%%%%%%%%%%

\section{Operator Theory}

\begin{center}

\begin{tabular}{ll}

Symbol & Meaning
\\
\hline

$T$ &
Operator
\\

$I$ &
Identity operator
\\

$T^{-1}$ &
Inverse operator
\\

$T_t$ &
Operator semigroup
\\

$e^{tA}$ &
Operator exponential
\\

$\sigma(T)$ &
Spectrum
\\

$r(T)$ &
Spectral radius
\\

$\operatorname{End}(X)$ &
Endomorphism monoid

\end{tabular}

\end{center}

%%%%%%%%%%%%%%%%%%%%%%%%%%%%%%%%%%%%%%%%%%%%%%%%%%%%%%%%%%%%%%%

\section{Optimization}

\begin{center}

\begin{tabular}{ll}

Symbol & Meaning
\\
\hline

$J$ &
Objective functional
\\

$P_C$ &
Metric projection
\\

$\operatorname{prox}$ &
Proximal operator
\\

$\Gamma$ &
Admissible path family
\\

$\gamma$ &
Continuation path
\\

$\lambda$ &
Lagrange multiplier

\end{tabular}

\end{center}

%%%%%%%%%%%%%%%%%%%%%%%%%%%%%%%%%%%%%%%%%%%%%%%%%%%%%%%%%%%%%%%

\section{Category Theory}

\begin{center}

\begin{tabular}{ll}

Symbol & Meaning
\\
\hline

$\mathcal C$ &
Category
\\

$\operatorname{Hom}$ &
Hom-set
\\

$F,G$ &
Functors
\\

$\eta$ &
Natural transformation
\\

$\operatorname{Lan}$ &
Left Kan extension
\\

$\operatorname{Ran}$ &
Right Kan extension
\\

$\circ$ &
Composition

\end{tabular}

\end{center}

%%%%%%%%%%%%%%%%%%%%%%%%%%%%%%%%%%%%%%%%%%%%%%%%%%%%%%%%%%%%%%%

\section{Sheaf Theory}

\begin{center}

\begin{tabular}{ll}

Symbol & Meaning
\\
\hline

$\mathcal F$ &
Sheaf
\\

$\rho^U_V$ &
Restriction morphism
\\

$\mathcal F(U)$ &
Sections
\\

$\mathcal F_x$ &
Stalk
\\

$H^k(X,\mathcal F)$ &
Sheaf cohomology

\end{tabular}

\end{center}

%%%%%%%%%%%%%%%%%%%%%%%%%%%%%%%%%%%%%%%%%%%%%%%%%%%%%%%%%%%%%%%

\section{Persistence Framework}

\begin{center}

\begin{tabular}{ll}

Symbol & Meaning
\\
\hline

$\mathcal O_{\mathrm{obs}}$ &
Observation operator
\\

$\mathcal S$ &
Configuration space
\\

$\Delta_R$ &
Repair defect functional
\\

$\mathcal E$ &
Evaluation functional
\\

$R$ &
Repair operator
\\

$\Phi_t$ &
Continuation flow
\\

$\mathcal P$ &
Persistence operator
\\

$\mathcal A$ &
Admissibility manifold
\\

$\mathcal M$ &
Configuration manifold
\\

$\delta_R$ &
Local repair defect

\end{tabular}

\end{center}

%%%%%%%%%%%%%%%%%%%%%%%%%%%%%%%%%%%%%%%%%%%%%%%%%%%%%%%%%%%%%%%

\section{Logical Symbols}

\begin{center}

\begin{tabular}{ll}

Symbol & Meaning
\\
\hline

$\forall$ &
For all
\\

$\exists$ &
There exists
\\

$\exists !$ &
There exists exactly one
\\

$\Rightarrow$ &
Implies
\\

$\Leftrightarrow$ &
If and only if
\\

$\subseteq$ &
Subset
\\

$\subset$ &
Proper subset
\\

$\cup$ &
Union
\\

$\cap$ &
Intersection
\\

$\circ$ &
Composition
\\

$\longrightarrow$ &
Map
\\

$\mapsto$ &
Element mapping

\end{tabular}

\end{center}

%%%%%%%%%%%%%%%%%%%%%%%%%%%%%%%%%%%%%%%%%%%%%%%%%%%%%%%%%%%%%%%

\section{Abbreviations}

\begin{center}

\begin{tabular}{ll}

Abbreviation & Meaning
\\
\hline

ODE &
Ordinary Differential Equation
\\

PDE &
Partial Differential Equation
\\

KKT &
Karush--Kuhn--Tucker
\\

ARS &
Abstract Reduction System
\\

l.s.c. &
Lower semicontinuous
\\

RHS &
Right-hand side
\\

LHS &
Left-hand side

\end{tabular}

\end{center}

%%%%%%%%%%%%%%%%%%%%%%%%%%%%%%%%%%%%%%%%%%%%%%%%%%%%%%%%%%%%%%%

\section{Indexing Conventions}

\begin{itemize}

\item Roman letters denote points, vectors, operators, and functions.

\item Greek letters denote paths, parameters, transformations, and
multipliers.

\item Script letters denote spaces, categories, sheaves, and structured
collections.

\item Blackboard bold symbols denote standard numerical systems.

\item Boldface is reserved for matrices and finite-dimensional vectors,
where required.

\end{itemize}

%%%%%%%%%%%%%%%%%%%%%%%%%%%%%%%%%%%%%%%%%%%%%%%%%%%%%%%%%%%%%%%

\section{Reading Guide}

Readers interested primarily in applications may proceed directly through
the main chapters, consulting this appendix whenever notation is
unclear. Readers interested in the mathematical foundations are
encouraged to use Appendix~L together with this appendix to trace each
symbol to its originating definition and each theorem to its supporting
results.

%%%%%%%%%%%%%%%%%%%%%%%%%%%%%%%%%%%%%%%%%%%%%%%%%%%%%%%%%%%%%%%

\section{Closing Remark}

The notation has been chosen to satisfy three design principles:

\begin{enumerate}

\item Symbols remain consistent throughout the monograph.

\item Related mathematical structures employ related notation whenever
possible.

\item Every nonstandard symbol is explicitly introduced before use and
listed in this appendix.

\end{enumerate}

Consequently, the notation forms a coherent mathematical vocabulary
supporting the development of persistence, repair, continuation, and
operator calculus from elementary topological concepts through higher
categorical constructions.

\newpage

\begin{thebibliography}{99}

\newcommand{\bibheading}[1]{%
  \vspace{1.2em}%
  \item[]\hspace*{-\leftmargin}\textbf{\large #1}%
  \vspace{0.4em}%
}

%%%%%%%%%%%%%%%%%%%%%%%%%%%%%%%%%%%%%%%%%%%%%%%%%%%%
\bibheading{Foundations of Mathematics}
%%%%%%%%%%%%%%%%%%%%%%%%%%%%%%%%%%%%%%%%%%%%%%%%%%%%

\bibitem{bourbaki}
Nicolas Bourbaki.
\textit{Elements of Mathematics}.
Springer.

\bibitem{halmos}
Paul R. Halmos.
\textit{Naive Set Theory}.
Springer, 1974.

\bibitem{lang}
Serge Lang.
\textit{Basic Mathematics}.
Springer.

\bibitem{munkres}
James R. Munkres.
\textit{Topology}.
2nd ed.
Prentice Hall, 2000.

%%%%%%%%%%%%%%%%%%%%%%%%%%%%%%%%%%%%%%%%%%%%%%%%%%%%
\bibheading{Metric Spaces and Analysis}
%%%%%%%%%%%%%%%%%%%%%%%%%%%%%%%%%%%%%%%%%%%%%%%%%%%%

\bibitem{rudin}
Walter Rudin.
\textit{Principles of Mathematical Analysis}.
3rd ed.
McGraw--Hill, 1976.

\bibitem{royden}
Halsey Royden and Patrick Fitzpatrick.
\textit{Real Analysis}.
Pearson.

\bibitem{conway}
John B. Conway.
\textit{A Course in Functional Analysis}.
Springer.

\bibitem{brezis}
Haim Brézis.
\textit{Functional Analysis, Sobolev Spaces and Partial Differential Equations}.
Springer.

%%%%%%%%%%%%%%%%%%%%%%%%%%%%%%%%%%%%%%%%%%%%%%%%%%%%
\bibheading{Convex Analysis and Optimization}
%%%%%%%%%%%%%%%%%%%%%%%%%%%%%%%%%%%%%%%%%%%%%%%%%%%%

\bibitem{rockafellar}
R. Tyrrell Rockafellar.
\textit{Convex Analysis}.
Princeton University Press.

\bibitem{boyd}
Stephen Boyd and Lieven Vandenberghe.
\textit{Convex Optimization}.
Cambridge University Press.

\bibitem{bertsekas}
Dimitri P. Bertsekas.
\textit{Nonlinear Programming}.
Athena Scientific.

%%%%%%%%%%%%%%%%%%%%%%%%%%%%%%%%%%%%%%%%%%%%%%%%%%%%
\bibheading{Differential Geometry}
%%%%%%%%%%%%%%%%%%%%%%%%%%%%%%%%%%%%%%%%%%%%%%%%%%%%

\bibitem{lee}
John M. Lee.
\textit{Introduction to Smooth Manifolds}.
Springer.

\bibitem{docarmo}
Manfredo P. do Carmo.
\textit{Riemannian Geometry}.
Birkhäuser.

\bibitem{spivak}
Michael Spivak.
\textit{A Comprehensive Introduction to Differential Geometry}.
Publish or Perish.

%%%%%%%%%%%%%%%%%%%%%%%%%%%%%%%%%%%%%%%%%%%%%%%%%%%%
\bibheading{Differential Equations and Dynamical Systems}
%%%%%%%%%%%%%%%%%%%%%%%%%%%%%%%%%%%%%%%%%%%%%%%%%%%%

\bibitem{hirsch}
Morris W. Hirsch, Stephen Smale, and Robert L. Devaney.
\textit{Differential Equations, Dynamical Systems, and an Introduction to Chaos}.
Academic Press.

\bibitem{hale}
Jack K. Hale.
\textit{Ordinary Differential Equations}.
Dover.

\bibitem{perko}
Lawrence Perko.
\textit{Differential Equations and Dynamical Systems}.
Springer.

%%%%%%%%%%%%%%%%%%%%%%%%%%%%%%%%%%%%%%%%%%%%%%%%%%%%
\bibheading{Viability Theory}
%%%%%%%%%%%%%%%%%%%%%%%%%%%%%%%%%%%%%%%%%%%%%%%%%%%%

\bibitem{nagumo}
Mitio Nagumo.
``Über die Lage der Integralkurven gewöhnlicher Differentialgleichungen.''
\textit{Proceedings of the Physico-Mathematical Society of Japan},
1942.

\bibitem{aubin}
Jean-Pierre Aubin.
\textit{Viability Theory}.
Birkhäuser.

%%%%%%%%%%%%%%%%%%%%%%%%%%%%%%%%%%%%%%%%%%%%%%%%%%%%
\bibheading{Calculus of Variations}
%%%%%%%%%%%%%%%%%%%%%%%%%%%%%%%%%%%%%%%%%%%%%%%%%%%%

\bibitem{giaquinta}
Mariano Giaquinta and Stefan Hildebrandt.
\textit{Calculus of Variations}.
Springer.

\bibitem{dacorogna}
Bernard Dacorogna.
\textit{Direct Methods in the Calculus of Variations}.
Springer.

%%%%%%%%%%%%%%%%%%%%%%%%%%%%%%%%%%%%%%%%%%%%%%%%%%%%
\bibheading{Operator Theory}
%%%%%%%%%%%%%%%%%%%%%%%%%%%%%%%%%%%%%%%%%%%%%%%%%%%%

\bibitem{reed}
Michael Reed and Barry Simon.
\textit{Methods of Modern Mathematical Physics}.
Academic Press.

\bibitem{yosida}
Kôsaku Yosida.
\textit{Functional Analysis}.
Springer.

\bibitem{engel}
Klaus-Jochen Engel and Rainer Nagel.
\textit{One-Parameter Semigroups for Linear Evolution Equations}.
Springer.

%%%%%%%%%%%%%%%%%%%%%%%%%%%%%%%%%%%%%%%%%%%%%%%%%%%%
\bibheading{Fixed Point Theory}
%%%%%%%%%%%%%%%%%%%%%%%%%%%%%%%%%%%%%%%%%%%%%%%%%%%%

\bibitem{banach}
Stefan Banach.
``Sur les opérations dans les ensembles abstraits.''
\textit{Fundamenta Mathematicae},
1922.

\bibitem{tarski}
Alfred Tarski.
``A Lattice-Theoretical Fixed Point Theorem.''
\textit{Pacific Journal of Mathematics},
5 (1955), 285--309.

%%%%%%%%%%%%%%%%%%%%%%%%%%%%%%%%%%%%%%%%%%%%%%%%%%%%
\bibheading{Term Rewriting}
%%%%%%%%%%%%%%%%%%%%%%%%%%%%%%%%%%%%%%%%%%%%%%%%%%%%

\bibitem{baader}
Franz Baader and Tobias Nipkow.
\textit{Term Rewriting and All That}.
Cambridge University Press.

\bibitem{newman}
M. H. A. Newman.
``On Theories with a Combinatorial Definition of Equivalence.''
\textit{Annals of Mathematics},
43 (1942), 223--243.

%%%%%%%%%%%%%%%%%%%%%%%%%%%%%%%%%%%%%%%%%%%%%%%%%%%%
\bibheading{Category Theory}
%%%%%%%%%%%%%%%%%%%%%%%%%%%%%%%%%%%%%%%%%%%%%%%%%%%%

\bibitem{maclane}
Saunders Mac Lane.
\textit{Categories for the Working Mathematician}.
Springer.

\bibitem{riehl}
Emily Riehl.
\textit{Category Theory in Context}.
Dover.

%%%%%%%%%%%%%%%%%%%%%%%%%%%%%%%%%%%%%%%%%%%%%%%%%%%%
\bibheading{Sheaf Theory}
%%%%%%%%%%%%%%%%%%%%%%%%%%%%%%%%%%%%%%%%%%%%%%%%%%%%

\bibitem{bredon}
Glen E. Bredon.
\textit{Sheaf Theory}.
Springer.

\bibitem{maclane-moerdijk}
Saunders Mac Lane and Ieke Moerdijk.
\textit{Sheaves in Geometry and Logic}.
Springer.

%%%%%%%%%%%%%%%%%%%%%%%%%%%%%%%%%%%%%%%%%%%%%%%%%%%%
\bibheading{Information Theory}
%%%%%%%%%%%%%%%%%%%%%%%%%%%%%%%%%%%%%%%%%%%%%%%%%%%%

\bibitem{shannon}
Claude E. Shannon.
``A Mathematical Theory of Communication.''
\textit{Bell System Technical Journal},
27 (1948), 379--423, 623--656.

\bibitem{cover}
Thomas M. Cover and Joy A. Thomas.
\textit{Elements of Information Theory}.
Wiley.

%%%%%%%%%%%%%%%%%%%%%%%%%%%%%%%%%%%%%%%%%%%%%%%%%%%%
\bibheading{Graph Theory and Discrete Geometry}
%%%%%%%%%%%%%%%%%%%%%%%%%%%%%%%%%%%%%%%%%%%%%%%%%%%%

\bibitem{ollivier}
Yann Ollivier.
``Ricci Curvature of Markov Chains on Metric Spaces.''
\textit{Journal of Functional Analysis},
256 (2009), 810--864.

\bibitem{forman}
Robin Forman.
``Bochner's Method for Cell Complexes and Combinatorial Ricci Curvature.''
\textit{Discrete and Computational Geometry},
29 (2003), 323--374.

%%%%%%%%%%%%%%%%%%%%%%%%%%%%%%%%%%%%%%%%%%%%%%%%%%%%
\bibheading{Optimization and Machine Learning}
%%%%%%%%%%%%%%%%%%%%%%%%%%%%%%%%%%%%%%%%%%%%%%%%%%%%

\bibitem{polyak}
Boris T. Polyak.
\textit{Introduction to Optimization}.
Optimization Software.

\bibitem{lojasiewicz}
Stanisław Łojasiewicz.
``A Topological Property of Real Analytic Subsets.''
In \textit{Coll. du CNRS, Les Équations aux Dérivées Partielles},
1963.

%%%%%%%%%%%%%%%%%%%%%%%%%%%%%%%%%%%%%%%%%%%%%%%%%%%%
\bibheading{Classical References}
%%%%%%%%%%%%%%%%%%%%%%%%%%%%%%%%%%%%%%%%%%%%%%%%%%%%

\bibitem{hilbert}
David Hilbert.
\textit{Foundations of Geometry}.
Open Court.

\bibitem{pontryagin}
L. S. Pontryagin, V. G. Boltyanskii,
R. V. Gamkrelidze, and E. F. Mishchenko.
\textit{The Mathematical Theory of Optimal Processes}.
Interscience.

\bibitem{arnold}
Vladimir I. Arnold.
\textit{Mathematical Methods of Classical Mechanics}.
Springer.

\end{thebibliography}

\end{document}